\documentclass[11pt,a4paper]{article}
\usepackage[top=2.5cm,margin=3cm,bottom=2.5cm]{geometry}

\usepackage[italian]{babel}

\usepackage[utf8]{inputenc}
\usepackage[T1]{fontenc}
\usepackage{amsfonts,amsmath,amsthm,lmodern}

\usepackage{accents}
\usepackage{bbm}
\usepackage{cancel}
\usepackage{graphicx}
\usepackage{braket}
\usepackage{blindtext}
\usepackage{mathtools}
\usepackage{titlesec}
\usepackage{tikz}
\usepackage{caption}

\usetikzlibrary{patterns}
\usetikzlibrary{arrows.meta}

\DeclareMathOperator{\Tr}{Tr}

\graphicspath{ {./images/} }

\title{\textbf{Appunti di Meccanica Quantistica}}
\author{}
\date{}

\begin{document}
\maketitle

\tableofcontents
\pagebreak

\section*{Riferimenti bibliografici}
\addcontentsline{toc}{section}{Riferimenti bibliografici}
\begin{enumerate}
\item J.J. Sakurai, \textit{Meccanica Quantistica Moderna}, Zanichelli, 1996.
\item F. Schwabl, \textit{Quantum Mechanics}, Springer, 2007.
\end{enumerate}
\pagebreak

\section{Primo modulo}
\subsection{Esperimento mentale}
\begin{figure}
\centering
\begin{tikzpicture}
  \node (A) at (5,0) {};
  \node (B) at (5,-2) {\(B_1\)};
  \node (C) at (5,-2.1) {};
  \node (D) at (5,-4.2) {\(B_2\)};
  \node (E) at (5,-4.3) {};
  \node (F) at (5,-6.3) {parete impermeabile};
  \node (Circle) at (-.3,-3.1) {\(A\)};

  \node (G) at (10,0) {};
  \node (H) at (10, -6.3) {schermo};

  \node (I) at (.1,-3.0) {};
  \node (J) at (2.5,-2.2) {granelli di sabbia};
  \node (K) at (.15,-3.1) {};
  \node (L) at (2.0,-3) {};
  \node (M) at (.1,-3.2) {};
  \node (N) at (1.5,-3.8) {};

  \draw[thick] (A) -- (B);
  \draw[thick] (C) -- (D);
  \draw[thick] (E) -- (F);
  \draw[thick] (G) -- (H);

  \draw (0,-3.1) circle (1.5pt);
  \draw (Circle);

  \draw (I) -- (J);
  \draw (K) -- (L);
  \draw (M) -- (N);
\end{tikzpicture}
\caption{Una sorgente (\(A\)) emette granelli di sabbia verso una parete permeabile con due fenditure (\(B_1\) e \(B_2\)). Lo schermo registra per ogni punto quanti granelli sono arrivati.}
\label{doppia_fenditura}
\end{figure}
Lo schema dell'apparato sperimentale \`e riportato in figura \ref{doppia_fenditura}. Una fenditura aperta, piccata sulla fenditura (figura \ref{granelli_una_fenditura}). Due fenditure aperte, ho la somma (figura \ref{granelli_due_fenditure}). Ho un'equazione del moto deterministica (anche se non la so risolvere). Per\`o ho una probabilit\`a.
\[
p(A \rightarrow C_k) = p(A \rightarrow B_1) \; p(B_1 \rightarrow C_k) + p(A \rightarrow B_2) \; p(B_2 \rightarrow C_k)
\]
perch\'e la particella passa o in \(B_1\) o in \(B_2\). Quindi
\begin{align*}
N(A \rightarrow C_k) &= N_{\text{TOT}} \sum_{i = 1}^2 p(A \rightarrow B_i) \; p(B_i \rightarrow C_k) =  \\
&= \sum_{i = 1}^2 N(A \rightarrow B_i) \; p(B_i \rightarrow C_k) = \\
&= \sum_{i = 1}^2 N(B_i \rightarrow C_k)
\end{align*}
\begin{figure}
\centering
\begin{tikzpicture}
  \draw[smooth, thick] plot file {./grafici/granelli_una_fend_1.txt};
  \draw[smooth, thick] plot file {./grafici/granelli_una_fend_2.txt};

  \node (A) at (-.5,-3) {};
  \node (B) at (-.5,3) {};

  \node (C) at (-3.5, 3) {};
  \node (D) at (-3.5, 1.05) {};

  \node (E) at (-3.5, .95) {};
  \node (F) at (-3.5, -.95) {};

  \node (G) at (-3.5, -1.05) {};
  \node (H) at (-3.5, -3) {};

  \draw (A) -- (B);
  \draw (C) -- (D);
  \draw (E) -- (F);
  \draw (G) -- (H);
\end{tikzpicture}
\caption{le due distribuzioni nel caso di granelli di sabbia in cui si lasci aperta una sola fenditura per volta. Sono piccate sulle fenditure.}
\label{granelli_una_fenditura}
\end{figure}
\begin{figure}
\centering
\begin{tikzpicture}
  \draw[smooth, thick] plot file {./grafici/granelli_due_fend.txt};

  \node (A) at (-.5,-3) {};
  \node (B) at (-.5,3) {};

  \node (C) at (-3.5, 3) {};
  \node (D) at (-3.5, 1.05) {};

  \node (E) at (-3.5, .95) {};
  \node (F) at (-3.5, -.95) {};

  \node (G) at (-3.5, -1) {};
  \node (H) at (-3.5, -3) {};

  \node (I) at (-.3, -1.2) {};
  \node (J) at (-.7, -1.2) {};

  \node (K) at (-.3, -1.3) {};
  \node (L) at (-.7, -1.3) {};

  \node (M) at (-.1, -1.3) {\(C_k\)};

  \draw(A) -- (B);
  \draw(C) -- (D);
  \draw(E) -- (F);
  \draw(G) -- (H);
  \draw(I) -- (J);
  \draw(K) -- (L);
\end{tikzpicture}
\caption{distribuzione di granelli di sabbia ottenuta con entrambe le fenditure aperte. \(p(A \to C_k)\) rappresenta la probabilit\`a che un granello di sabbia arrivi sullo schermo in posizione \(C_k\).}
\label{granelli_due_fenditure}
\end{figure}

\bigskip
\noindent
Quando la sorgente emette fullereni, le cose cambiano, bench\'e siano oggetti piuttosto grandi (diametro 1 nm). L'esempio \`e stato fatto anche con molecole pi\`u grandi (\( \text{C}_{60} \text{F}_{48} \) e TPP). L'esperimento \`e stato fatto da M.~Arndt e A.~Zeilinger (figura \ref{arndt_zeilinger}).

\bigskip
\noindent
\textbf{N.B.:} spara molecole singole.
\begin{figure}
\centering
\begin{tikzpicture}
  \draw (-5, 2.7) circle (1.5pt);

  \draw[pattern=north west lines, pattern color=black, draw=black] (0,0) rectangle ++(0.1,2.0);
  \draw[pattern=north west lines, pattern color=black, draw=black] (0,2.2) rectangle ++(0.1,1.1);
  \draw[pattern=north west lines, pattern color=black, draw=black] (0,3.5) rectangle ++(0.1,2.0);

  \draw [line width=0.25mm] (3, 5.5) -- (3, 0) node [right] {fascio laser};
  \draw [->] (3.5, 1.0) -- (3.5, 2.5);

  \node (A) at (0, 5.7) {200 nm};
  \node (B) at (1, 2.7) {100 nm};
  \node (C) at (1, 2.1) {52 nm};
  \node (D) at (0, -.3) {Si N};
  \node (E) at (4.2, 1.7) {8 \(\mu\)m};
  \node (F) at (3.7, 0.7) {2 \(\mu\)m};
  \node (G) at (-5.0, 3.0) {C\(_{60}\)};
  \node (H) at (-5.0, 2.3) {\(v = 220\) km/h};
\end{tikzpicture}
\caption{Schema dell'esperimento di Arndt e Zeilinger.}
\label{arndt_zeilinger}
\end{figure}

\bigskip
\noindent
L'istogramma ha un andamento come in figura \ref{arndt_zeilinger_distro}.
\begin{figure}
\centering
\begin{tikzpicture}
  \draw (0, 0) circle (1.5pt);

  \node (A) at (3.5, 3) {};
  \node (B) at (3.5, 1.05) {};

  \node (C) at (3.5, .95) {};
  \node (D) at (3.5, -.95) {};

  \node (E) at (3.5, -1) {};
  \node (F) at (3.5, -3) {};

  \draw (A) -- (B);
  \draw (C) -- (D);
  \draw (E) -- (F);

  \begin{scope}[shift={(4,0)}]
    \draw[smooth, thick] plot file {./grafici/doppia_fenditura_1.txt};
    \draw[smooth, dashed] plot file {./grafici/doppia_fenditura_2.txt};
  \end{scope}
\end{tikzpicture}
\caption{distribuzione nell'esperimento di Arndt e Zeilinger.}
\label{arndt_zeilinger_distro}
\end{figure}
\[
p(A \rightarrow C) \sim \left| \sum_i R(A \rightarrow B_i) \; R(B_i \rightarrow C_k)\right|^2 
\]
\[
R = \cos(\omega t + \varphi)
\]
figura di interferenza data, per\`o, da una sorgente di onde per il principio di Fresnel. \`E come se sommassi due onde
\[
R_1 = A \cos \left( \omega t + \varphi_1 \right) \qquad R_2 = A \cos \left( \omega t + \varphi_2 \right)
\]
con uguali ampiezza e frequenza e diverse fasi, dove so che
\[
\varphi = 2 \pi \frac{d}{\lambda} - \left[ \frac{2 \pi d}{\lambda} \right] \qquad \varphi_1 - \varphi_2 = 2 \pi \frac{\Delta d}{\lambda}
\]
\begin{align*}
R &= R_1 + R_2 = \\
&= A \; \mathrm{Re} \left( e^{i (\omega t + \varphi_1)} + e^{i (\omega t + \varphi_2)} \right) = \\
&= A \; \mathrm{Re} \left( e^{i \omega t}\left( e^{i \varphi_1} + e^{i \varphi_2} \right) \right) = \\
&= A \; \mathrm{Re} \left( e^{i \frac{\varphi_1 + \varphi_2}{2} + i \omega t} \left( e^{i \frac{\varphi_1 - \varphi_2}{2}} + e^{i \frac{\varphi_2 - \varphi_1}{2}} \right) \right) = \\
&= A \; \mathrm{Re} \left( e^{i \frac{\varphi_1 + \varphi_2}{2} + i  \omega t} \left( e^{i \frac{\varphi_1 - \varphi_2}{2}} + e^{-i \frac{\varphi_1 - \varphi_2}{2}} \right) \right) = \\
&= 2 A \; \cos \frac{\varphi_1 - \varphi_2}{2} \; \cos \left( \omega t + \frac{\varphi_1 + \varphi_2}{2}  \right)
\end{align*}
e la nuova ampiezza \`e
\[
2 A \; \cos \frac{\varphi_1 - \varphi_2}{2} = \sqrt{\left| R_1 + R_2 \right|^2}
\]
Il fatto \`e che la molecola dovrebbe passare o in una o nell'altra delle fenditure, quindi dovrebbe essere come i granelli di sabbia. Invece ho interferenza, quindi l'onda \`e data dal fatto che sto misurando. Non c'\`e un'onda fisica perch\'e sto lavorando con molecole e perch\'e la figura di interferenza \`e tanto pi\`u visibile quanto pi\`u si riescono ad isolare singole molecole con una certa velocit\`a.

\bigskip
\noindent
Se per\`o inserisco un rilevatore che dica in quale fenditura passa la singola molecola, sparisce l'interferenza (figura \ref{rilevatore}).
\begin{figure}
\centering
\begin{tikzpicture}
  \draw[smooth, thick] plot file {./grafici/granelli_due_fend.txt};

  \node (C) at (-1.0, 3) {};
  \node (D) at (-1.0, 1.05) {};

  \node (E) at (-1.0, .95) {};
  \node (F) at (-1.0, -.95) {};

  \node (G) at (-1.0, -1) {};
  \node (H) at (-1.0, -3) {};
  \node (I) at (-.4,-1.8) {\(R\)};

  \draw[thick] (C) -- (D);
  \draw[thick] (E) -- (F);
  \draw[thick] (G) -- (H);

  \draw[draw=black] (-.8,-2.0) rectangle ++(.2,.5);
  \draw (-4.0, 0) circle (1.5pt);
\end{tikzpicture}
\caption{Se viene utilizzato un rilevatore \(R\) per individuare da quale fenditura passa la molecola, si ritrova la distribuzione che ci si aspetterebbe nel caso classico.}
\label{rilevatore}
\end{figure}
Allora l'esperimento non viola la causalit\`a e il principio di localit\`a. Secondo l'interpretazione di Copenhagen:
\begin{quote}
Esiste un mondo classico ed esiste un mondo quantistico, ma l'oggetto su cui faccio la misura fa parte del mondo quantistico e il rilevatore \`e un oggetto classico. Le leggi che descrivono gli oggetti quindi sono ondulatorie, ma per misurare una grandezza quantistica uso un oggetto classico che fa precipitare il sistema dallo stato quantistico a quello classico.
\end{quote}
\subsection{Sistema QUBIT}
Sistema che pu\`o esistere in due stati \(\rightarrow\) notazione di Dirac. \(\ket{A}\) ket \(\rightarrow\) stati del sistema \(\sim\) vettore.
\begin{figure}
\centering
\begin{tikzpicture}
  \node (A) at (0, 3) {};
  \node (B) at (0, 1.05) {};

  \node (C) at (0, .95) {};
  \node (D) at (0, -.95) {};

  \node (E) at (0, -1) {};
  \node (F) at (0, -3) {};

  \node (G) at (2, 3) {};
  \node (H) at (2, -3) {};

  \node (I) at (-2, 0) {\( \ket{\psi} \)};
  \node (J) at (.7, .95) {\( \ket{+} \)};
  \node (K) at (.7,-1) {\( \ket{-} \)};

  \draw[thick] (A) -- (B);
  \draw[thick] (C) -- (D);
  \draw[thick] (E) -- (F);
  \draw[thick] (G) -- (H);
\end{tikzpicture}
\caption{Nell'esperimento della doppia fenditura, lo stato \( \ket{\psi} \) \`e la combinazione degli stati \( \ket{+} \) e \( \ket{-} \).}
\label{stati_qubit}
\end{figure}
In questo caso indico gli stati con \(\ket{+}\) e \(\ket{-}\) \(\Rightarrow \ket{\pm}\) (figura \ref{stati_qubit}).

\bigskip
\noindent
Se non misuro da quale fenditura passa la particella allora ho sovrapposizione, cio\`e
\[
\ket{\psi} = a_+ \ket{+} + a_- \ket{-} \qquad \text{con } a_+, a_- \in \mathbb{C}, \mathrm{ampiezze}
\]
e \(\ket{\psi}\) \`e ancora uno stato del sistema. La probabilit\`a allora \`e definita come
\[
P_\pm = |a_\pm|^2
\]
quando sto misurando.

\bigskip
\noindent
Questa probabilit\`a \`e tale che \(P_+ + P_- = 1\), condizione di normalizzazione su \(\ket{\psi}\). A priori, \(a_\pm\) non sono determinabili.

\bigskip
\noindent
Una peculiarit\`a \`e quella di poter applicare il principio  di sovrapposizione. Quindi (senza fare misure) gli stati si sovrappongono linearmente (figura \ref{sovrapposizione_stati}).
\begin{figure}
\centering
\begin{tikzpicture}
  \node (A) at (0, 3) {};
  \node (B) at (0, 1.05) {};

  \node (C) at (0, .95) {};
  \node (D) at (0, -.95) {};

  \node (E) at (0, -1) {};
  \node (F) at (0, -3) {};

  \node (G) at (2, 3) {};
  \node (H) at (2, 1.05) {};

  \node (I) at (2, .95) {};
  \node (J) at (2, -.95) {};

  \node (K) at (2, -1) {};
  \node (L) at (2, -3) {};

  \node (M) at (-2, 0) {\( \ket{\psi} \)};
  \node (N) at (.7, .95) {\( \ket{+} \)};
  \node (O) at (.7,-1) {\( \ket{-} \)};
  \node (P) at (2.7, .95) {\( \ket{\psi_1} \)};
  \node (Q) at (2.7,-1) {\( \ket{\psi_2} \)};

  \draw[thick] (A) -- (B);
  \draw[thick] (C) -- (D);
  \draw[thick] (E) -- (F);
  \draw[thick] (G) -- (H);
  \draw[thick] (I) -- (J);
  \draw[thick] (K) -- (L);
\end{tikzpicture}
\caption{Esempio di sovrapposizione lineare degli stati.}
\label{sovrapposizione_stati}
\end{figure}
\[
\ket{\psi_1} = a_+^{(1)} \ket{+} + a_-^{(1)} \ket{-}
\]
\[
\ket{\psi_2} = a_+^{(2)} \ket{+} + a_-^{(2)} \ket{-}
\]
Sovrapposizione
\[
\ket{\psi} = c_1 \ket{\psi_1} + c_2 \ket{\psi_2}
\]
Quando misuro, sapendo che
\[
c_1 = c_2 = \frac{1}{\sqrt{2}}
\]
per semplicit\`a, e 
\[
\ket{\psi} = \frac{1}{\sqrt{2}} \left( \left( a_+^{(1)} + a_+^{(2)} \right) \ket{+} + \left( a_-^{(1)} + a_-^{(2)} \right) \ket{-} \right)
\]
ricavo le probabilit\`a
\begin{align*}
P_\pm &= \left| a_\pm^{(1)} + a_\pm^{(2)} \right|^2 = \\
&= \left| a_\pm^{(1)} \right|^2 + \left| a_\pm^{(2)} \right|^2 + \left( a_+^{(1)}\right)^* a_\pm^{(2)} + \left( a_\pm^{(2)} \right)^* a_\pm^{(1)} \neq \\
&\neq P_\pm^{(1)} + P_\pm^{(2)}
\end{align*}
Allora suppongo che ci sia un prodotto scalare: bra-ket.
\[
\ket{\psi}
\]
\[
\ket{\varphi} \rightarrow \mathrm{bra } \bra{\varphi}
\]
\[
\text{prodotto scalare} \braket{\varphi|\psi}
\]
e lo definisco con le propriet\`a
\begin{enumerate}
\item \(\braket{\varphi|\psi} = (\braket{\psi|\varphi})^*\)
\item \(\braket{\varphi|\varphi} \geq 0\)
\item se \(\ket{\psi} = a \ket{1} + b \ket{2}\) allora
\[
\braket{\varphi|\psi} = a \braket{\varphi|1} + b \braket{\varphi|2}
\]
\end{enumerate}
Cos\`{i} \(\ket{\psi} = a \ket{1} + b \ket{2}\) d\`a
\begin{align*}
\braket{\psi|w} &\stackrel{(1)}{=} \left( \braket{w|\psi} \right)^* \\
&\stackrel{(3)}{=} \left( a \braket{w|1} + b \braket{w|2} \right)^* \\
&\stackrel{(2)}{=} a^* \braket{1|w} + b^* \braket{2|w} \\
&\stackrel{(3)}{=} \left( a^* \bra{1} + b^* \bra{2} \right) \ket{w}
\end{align*}
cio\`e
\[
\bra{\psi} = a^* \bra{1} + b^* \bra{2}
\]
Stati esaustivi ed esclusivi \`e dato da
\[
\braket{+|+} = \braket{-|-}  = 1 \quad \text{esaustivo}
\]
\[
\braket{+|-} = \braket{-|+} = 0 \quad \text{esclusivo}
\]
Quindi
\[
P_\pm = \left| \braket{\pm|\psi} \right|^2 = \left| \Braket{\pm|\left( a_+ \ket{+} + a_- \ket{-} \right)} \right|^2 = |a_\pm|^2
\]
Ad esempio,
\[
\ket{\psi_1} = \frac{1}{\sqrt{2}} (\ket{+} + \ket{-})
\]
\[
\ket{\psi_2} = \frac{1}{\sqrt{2}} (\ket{+} - \ket{-})
\]
\[
P_\pm = \frac{1}{2}
\]
tuttavia non so se \(\psi_1\) e \(\psi_2\) sono uguali. Infatti,
\[
\braket{\psi_1|\psi_2} = \frac{1}{2} (\bra{+} + \bra{-})(\ket{+} - \ket{-}) = \frac{1}{2} (1 + 0 + 0 - 1) = 0
\]
cio\`e sono ortogonali, quindi sono diversi pur avendo la stessa probabilit\`a. Inoltre,
\[
\braket{\psi_1|\psi_1} = \braket{\psi_2|\psi_2} = 1
\]
Quindi se dallo stato \(\ket{\psi_1}\) misuro \(\ket{+}\) e da qui misuro la probabilit\`a associata allo stato \(\ket{\psi_2}\), classicamente otterrei 0, mentre qui ho la rigenerazione di \(\ket{\psi_2}\).
\[
\left| \braket{\psi_2|+} \right|^2 = \frac{1}{2} \left| \left( \bra{+} + \bra{-} \right) \ket{+} \right|^2 = \frac{1}{2}, \quad \text{cio\`e il 50\%}.
\]
La misura quindi modifica lo stato. Allo stesso modo, se non eseguo nessuna misura, la sovrapposizione \`e data da
\begin{align*}
\ket{\varphi} &= \frac{1}{\sqrt{2}} \left( \frac{1}{\sqrt{2}} \left( \ket{+} + \ket{-} \right) + \frac{1}{\sqrt{2}} \left( \ket{+} - \ket{-} \right) \right) = \\
&= \frac{1}{2} \left( \ket{+} + \ket{+} + \ket{-} - \ket{-} \right) = \\
&= \ket{+}
\end{align*}
cio\`e \`e sparito \(\ket{-}\) (cfr. le onde in controfase).

\bigskip
\noindent
Potrei avere
\[
\ket{\psi_A} = \frac{1}{\sqrt{2}} \left( \ket{+} + i \ket{-} \right), \qquad \ket{\psi_B} = \frac{1}{\sqrt{2}} \left( \ket{+} - i \ket{-} \right)
\]
e allora
\[
\braket{\psi_A|\psi_B} = \frac{1}{2} \left( \bra{+} - i \bra{-} \right) \left( \ket{+} - i \ket{-} \right) = \frac{1}{2} (1 - 1) = 0
\]
\[
\left| \braket{\psi_A|+} \right|^2 = \left| \frac{1}{\sqrt{2}} \left( \bra{+} - i \bra{-} \right) \ket{+} \right|^2 = \frac{1}{2}
\]
\[
\ket{\varphi} = \frac{1}{\sqrt{2}} \left( \ket{\psi_A} + \ket{\psi_B} \right) = \frac{1}{2} \left( \ket{+} + i \ket{-} + \ket{+} - i \ket{-} \right) = \ket{+}
\]
proprio come prima.

\bigskip
\noindent
Quindi se ho due stati
\[
\ket{\psi} = a_+ \ket{+} + a_- \ket{-}
\]
\[
\ket{\varphi} = b_+ \ket{+} + b_- \ket{-}
\]
allora
\[
\braket{\varphi|\psi} = \left( b_+^* \bra{+} + b_-^* \bra{-} \right) \left( a_+ \ket{+} + a_- \ket{-} \right) = a_+ b_+^* + a_- b_-^*
\]
e questo \`e leggibile in termini di sovrapposizione
\[
\ket{w} = N \left( c \ket{\varphi} + d \ket{\psi} \right)
\]
Allora \(\left| \braket{\varphi|\psi} \right|^2\) \`e una probabilit\`a che ci sia la sovrapposizione degli stati.

\bigskip
\noindent
Un altro esempio \`e
\[
\ket{\psi} = \frac{1}{\sqrt{2}} \left( \ket{+} + \ket{-} \right)
\]
\[
\ket{\varphi} = \frac{1}{\sqrt{2}} \left( \ket{+} + e^{i \vartheta} \ket{-} \right)
\]
\[
\left| \braket{\varphi|\psi} \right|^2 = \frac{1}{4} \left| 1 + e^{-i \vartheta} \right|^2
\]
Se faccio la misura
\[
\left| \braket{\varphi|+} \right|^2 = \frac{1}{2}
\]
\`e sparita la fase, quindi \`e cambiato lo stato.
\subsection{Nuova scrittura per gli stati}
\begin{align*}
\ket{\psi} &= \braket{+|\psi} \ket{+} + \braket{-|\psi} \ket{-} = \ket{+} \braket{+|\psi} \ket{-} \braket{-|\psi} = \\
&= \underbrace{\left( \ket{+} \bra{+} + \ket{-} \bra{-} \right)}_{\mathbbm{1}} \ket{\psi}
\end{align*}
In generale, gli stati \(\ket{e_i}\) tali che \(\braket{e_i|e_j} = \delta_{ij}\) danno
\[
\mathbbm{1} = \sum_i \ket{e_i}\bra{e_i} \qquad \text{relazione di completezza}
\]
Ad esempio,
\[
\braket{\varphi|\psi} = \sum_i \braket{\varphi|e_i} \braket{e_i|\psi}
\]
e allora
\[
\ket{\psi} = \sum_i \ket{e_i} \braket{e_i|\psi} = \sum_i \ket{e_i} c_i^{(\psi)}
\]
\[
\ket{\varphi} = \sum_i \ket{e_i} \braket{e_i|\varphi} = \sum_i \ket{e_i} c_i^{(\varphi)}
\]
\[
\braket{\varphi|\psi} = \sum_i \left( c_i^{(\varphi)} \right)^* c_i^{(\psi)}
\]
e quindi
\[
\braket{\psi|\psi} = \sum_i \braket{\psi|e_i}\braket{e_i|\psi} = \sum_i \left| c_i^{(\psi)} \right|^2
\]
\subsection{Operatori}

\`E un \(A \ket{\psi} = \ket{\psi'}\). Mi interesso agli operatori lineari. Quindi sono applicazioni, definite dallo spazio degli stati in s\'e.

Ho la decomposizione nella base \(\ket{e_i}\).
\[
A\ket{\psi} = \sum_i A \ket{e_i} c_i^{(\psi)}, \qquad \text{dove } A \ket{e_i} = \ket{\psi_i'}
\]
Ma anche
\[
\ket{\psi_i'} = \sum_j \ket{e_j}c_{ji}, \qquad \text{con } c_{ji} = \braket{e_j|\psi_i'} = \braket{e_j|A|e_i}
\]
Allora per lo stato \(\ket{\psi}\)
\[
A\ket{\psi} = \sum_{ij} \ket{e_j} c_{ij} c_i^{(\psi)}
\]
Nello stesso modo
\[
A\ket{\psi} = \sum_i A \ket{e_i} \braket{e_i|\psi} = \sum_{ij} \ket{e_j} \braket{e_j|A|e_i} \braket{e_i|\psi}
\]
\[
\braket{e_j|A|e_i} \text{elemento } ji \text{ della matrice}
\]
Quindi
\[
\ket{\psi} = \begin{pmatrix}
\braket{e_1|\psi} \\
\braket{e_2|\psi} \\
\vdots
\end{pmatrix}
\]
\[
A = \left( \braket{e_i|A|e_j} \right)_{i,j}
\]

L'operatore aggiunto \`e \(A^\dag\) e agisce su
\[
\bra{\psi}A = \bra{\psi'}
\]
Inoltre
\[
\braket{\varphi|A|\psi} = \braket{\varphi|\psi'}
\]
\[
\braket{\psi|A^\dag|\varphi} = \braket{\psi'|\varphi} = \left( \braket{\varphi|\psi'} \right)^*
\]
\[
A_{ij}^\dag = \braket{e_i|A^\dag|e_j} = \left( A_{ji} \right)^*
\]
\subsection{Esprimere le osservabili come operatori}
Il sistema si presenta come combinazione di stati \(\ket{e_i}\) e quando misuro ottengo un \(\lambda_i\) reale associato a \(\ket{e_i}\). Chiamo \(O\) questa osservabile. Allora
\begin{align*}
\text{valore medio} \quad \braket{O}_\psi &= \sum_i \lambda_i P_i = \sum_i \lambda_i \left| \braket{e_i|\psi} \right|^2 \\
&\stackrel{\text{def}}{=} \braket{\psi|O|\psi}
\end{align*}
\[
\text{operatore} \quad O = \sum_i \lambda_i \ket{e_i} \bra{e_i}
\]
\[
O \ket{\psi} = \sum_i \lambda_i \ket{e_i} \braket{e_i|\psi}
\]
\begin{align*}
\braket{O}_\psi &= \sum_{ij} \braket{\psi|e_i} \braket{e_i|O|e_j} \braket{e_j|\psi} = \\
&= \sum_{ijm} \braket{e_i|\psi} \underbrace{\braket{e_i|\lambda_m|e_m}}_{\lambda_m \delta_{im}} \underbrace{\braket{e_m|e_j}}_{\delta_{mj}} \braket{e_j|\psi} = \\
&= \sum_i \lambda_i \braket{\psi|e_i} \braket{e_i|\psi} = \\
&= \sum_i \lambda_i \left| \braket{\psi|e_i} \right|^2
\end{align*}
\[
\braket{O}_\psi = \sum_i \lambda_i \left| \braket{\psi|e_i} \right|^2 \qquad \braket{O}_\psi = \braket{\psi|O|\psi}
\]
Inoltre, la matrice \(O\) è diagonale
\[
O = \begin{pmatrix}
\lambda_1 & 0 \\
0 & \ddots
\end{pmatrix}
\]
In particolare
\[
\lambda_i \text{ è un autovalore reale}
\]
\[
\ket{e_i} \text{ è un autostato}
\]
Quindi ogni osservabile è un operatore tale che ammette autovalori reali e una BON di autostati e viceversa.

\subsubsection{Esempio: operatore come matrice}
\[
O = \ket{\psi_1}\bra{\psi_2} + \ket{\psi_2}\bra{\psi_1}
\]
\[
\ket{\psi_1} = \frac{1}{\sqrt{2}} \left( \ket{+} + i \ket{-} \right) \qquad \ket{\psi_2} = \frac{1}{\sqrt{2}} \left( \ket{+} - i \ket{-} \right)
\]
Poiché \(\braket{\pm|\pm} = 1\) e \(\braket{\pm|\mp} = 0\), ho la BON e
\begin{align*}
O_{11} &= \braket{+|\frac{1}{\sqrt{2}} \left( \ket{+} + i \ket{-} \right) \frac{1}{\sqrt{2}} \left( \bra{+} + i \bra{-} \right)|+} + \\
&+ \braket{+|\frac{1}{\sqrt{2}} \left( \ket{+} - i \ket{-} \right) \frac{1}{\sqrt{2}} \left( \bra{+} - i \bra{-} \right)|+} = \\
&= \frac{1}{2} + \frac{1}{2} = 1
\end{align*}
\[
O_{12} = -\frac{1}{2} + \frac{1}{2} = 0 \qquad O_{21} = \frac{1}{2} - \frac{1}{2} = 0 \qquad O_{22} = -1
\]
\[
O = \begin{pmatrix}
1 & 0 \\
0 & -1
\end{pmatrix}
\]

\subsubsection{Esempio: proiezione di uno stato}
\[
\ket{\psi} = \frac{1}{2} \ket{1} + \frac{1}{\sqrt{2}} \ket{2} + \frac{1}{\sqrt{2}} \ket{3}
\]
Normalizzo
\[
P_i = \frac{|a_i|^2}{\sum_{k=1}^3 |a_k|^2} = \begin{cases}
1/5 \\
2/5 \\
2/5
\end{cases}
\]
e ho le singole probabilità. La probabilità che non si trovi in \(\ket{2}\) è \(1 - P_2 = 3/5\).

\bigskip
\noindent
Potrei definire un'osservabile
\[
O = \lambda \ket{1}\bra{1} + \lambda \ket{3}\bra{3} + \lambda' \ket{2}\bra{2}
\]
e se misuro che il sistema si trova in \(\ket{2}\) so già, altrimenti se non si trova in \(\ket{2}\) è nello stato
\[
\ket{\psi'} = \frac{1}{2} \ket{1} + \frac{1}{\sqrt{2}} \ket{3}
\]
cioè proietto \(\ket{\psi}\) attraverso un operatore \(P\) tale che \(P^2 = P\) attraverso
\[
P_{\ket{S}} = \ket{S}\bra{S}, \qquad \text{non } S \rightarrow P = \mathbbm{1} - \ket{S}\bra{S}
\]

\bigskip
\noindent
Se si trova in \(\ket{2}\) e voglio calcolare la probabiltià con
\[
\ket{\varphi} = \frac{1}{\sqrt{2}} \left( \ket{2} - \ket{3} \right)
\]
\[
\left| \braket{2|\varphi} \right|^2 = \frac{1}{2} = P
\]
Inoltre,
\[
\left| \braket{\varphi|\psi} \right|^2 = 0
\]
cioè \(\ket{\varphi} \perp \ket{\psi}\) ed è l'informazione sul sistema \(\rightarrow\) il proiettore lascia invariato lo \(\ket{\varphi}\), infatti
\[
P = \mathbbm{1} - \ket{\varphi}\bra{\varphi}, \qquad P \ket{\psi} = \ket{\psi}
\]
e quindi la probabilità che il sistema sia in \(\ket{2}\).

\bigskip
\noindent
Quindi il proiettore
\[
P_\varphi = \ket{\varphi}\bra{\varphi}
\]
% \subsection{Buone osservabili}

Un operatore \`e autoaggiunto (hermitiano) quando
\[
A^\dag = A
\]
Le osservabili
\[
O = \sum_i \lambda_i \ket{e_i}\bra{e_i} \qquad \text{con } \braket{e_i|e_j} = \delta_{ij}
\]
sono tutti e soli gli operatori autoaggiunti. Infatti,
\[
O^\dag = \left( \sum_i \lambda_i \ket{e_i} \bra{e_i} \right)^\dag = \sum_i \lambda_i^* \ket{e_i} \bra{e_i} = O
\]
e quindi una l'ho dimostrata.

Inoltre, se cambio autostati (nuova BON)
\[
O = \sum_i \mu_i \ket{e_i'} \bra{e_i'} \qquad \text{con } \braket{e_i'|e_j'} = \delta_{ij}
\]
allora
\[
\braket{e_i'|\left( \sum_m \mu_m \ket{e_m'} \bra{e_m'} \right)|e_j'} = \begin{pmatrix}
\mu_1 & 0 \\
0 & \ddots
\end{pmatrix}
\]
Se vale
\[
O \ket{e_i} = \lambda_i \ket{e_i} \quad \Rightarrow \quad \braket{e_i|O|e_i} = \lambda_i \braket{e_i|e_i}\quad \Rightarrow
\]
\[
\quad \Rightarrow \qquad \lambda_i = \frac{\braket{e_i|O|e_i}}{\braket{e_i|e_i}}
\]
\[
\lambda_i^* = \frac{\braket{e_i|O^\dag|e_i}}{\braket{e_i|e_i}} = \frac{\braket{e_i|O|e_i}}{\braket{e_i|e_i}} = \lambda_i \quad \Rightarrow \quad \lambda_i \in \mathbb{R}
\]
Se invece vale
\[
\braket{e_j|O|e_i} = \lambda_i \braket{e_j|e_i}
\]
\[
\braket{e_i|O^\dag|e_j} = \lambda_i \braket{e_j|e_i}
\]
\[
\braket{e_i|O|e_j} = \lambda_j \braket{e_i|e_j} = \lambda_i \braket{e_i|e_j}
\]
\[
\left( \lambda_i - \lambda_j \right) \braket{e_i|e_j} = 0
\]
Allora se i $\lambda_i$ sono tutti diversi, $\braket{e_i|e_j} = 0$; se $\lambda_i = \lambda_j$, ho $\braket{e_i|e_j} \neq 0$. Cio\`e ho ortigonalit\`a, poi normalizzo per avere $= 1$.

Nel caso degenere $\ket{e_i^{(1)}}, \ket{e_i^{(2)}}, \dots$ associati a $\lambda_i$, la chiamo $\ket{e_i^{(a)}}$ e la proiezione
\[
\braket{e_j|e_i^{(a)}} = 0
\]
Poich\'e invece
\[
O \ket{e_i^{(a)}} = \lambda_i \ket{e_i^{(a)}} \quad \Rightarrow \quad O \left( \sum_a c_a \ket{e_i^{(a)}} \right) = \sum_a c_a O \ket{e_i^{(a)}} \quad \text{etc.}
\]

Riepilogo i due risultati ottenuti.
\begin{enumerate}
\item $O$ \`e un'osservabile se e solo se \`e un operatore autoaggiunto.
\item Se un operatore \`e autoaggiunto allora
\begin{enumerate}
\item ha autovalori reali,
\item ammette BON.
\end{enumerate}
\end{enumerate}
% \subsection{Costruire una BON}
\[
\ket{\psi_1}, \ket{\psi_2}, \dots, \ket{\psi_n}
\]
\[
\ket{e_1} = \ket{\psi_1}, \quad \ket{e_j} = \ket{\psi_j} - \sum_{i = 1}^{j - 1} \frac{\braket{e_i|\psi_j}}{\braket{e_i|e_i}} \ket{e_i}
\]
% \subsection{Postulati della Meccanica Quantistica}
\begin{enumerate}
\item Uno stato fisico associato ad un sistema \`e rappresentato da un vettore nello spazio di Hilbert, per cui vale il principio di sovrapposizione.
\item I possibili risultati associati ad una misura eseguita su un sistema sono rappresentati da stati esclusivi ed esaustivi.
\item La probabilit\`a della misura di uno stato del sistema \`e data dal modulo quadrato dell'ampiezza di probabilit\`a associata a quello stato.
\end{enumerate}
% \subsection{Cambiamenti di base ed operatori unitari}
\[
\ket{\psi} = \sum_i \ket{e_i} \braket{e_i|\psi} = \sum_i \ket{e_i'}\braket{e_i'|\psi}
\]
\[
\braket{e_i'|\psi} = \sum_j \braket{e_i'|e_j} \braket{e_j|\psi}
\]
e chiamo
\[
U_{ij} = \braket{e_i'|e_j}
\]
e \(U\) \`e la matrice di cambiamento di base, con
\[
U = \sum_m \ket{e_m} \bra{e_m'}
\]

Inoltre, se \(U \ket{\psi} = \ket{\psi'}\), vale che
\begin{enumerate}
\item \(\braket{e_i|\psi'} = \sum_m \braket{e_i|e_m'} \braket{e_m'|\psi} = \braket{e_i'|\psi}\)
\item \(U_{ij} = \sum_m \braket{e_i|e_m} \braket{e_m'|e_j} = \sum_m \braket{e_i'|e_m'} \braket{e_m'|e_j'} = \braket{e_i'|U|e_j'} = \braket{e_i'|e_j}\)
\end{enumerate}

\(U\) \`e unitario se
\[
U^\dag = U^{-1}, \qquad U^\dag U = \mathbbm{1}
\]
e scrivo
\[
U = \sum_m \ket{e_m} \bra{e_m'}
\]
\[
U^\dag = \sum_m \ket{e_m'} \bra{e_m}
\]
\[
U U^\dag = \sum_{n,m} \ket{e_n} \braket{e_n'|e_m'} \bra{e_m} = \sum_m \ket{e_m} \bra{e_m} = \mathbbm{1} = U^\dag U
\]
Gli operatori unitari sono matrici di cambio base.

Prendo due stati \(\ket{\psi}\) e \(\ket{\varphi}\), allora per \(U\) unitario
\[
\ket{\psi'} = U \ket{\psi} \qquad \ket{\varphi'} = U \ket{\varphi}
\]
\[
\braket{\psi'|\varphi'} = \braket{\psi|UU^\dag|\varphi} = \braket{\psi|\varphi}
\]
cio\`e i cambiamenti di base attraverso operatori unitari mantiene invariato il prodotto scalare (si conserva la probabilit\`a). Ho a che fare con la reversibilit\`a: in particolare, l'operatore di misura non \`e reversibile.
\[
\ket{\psi} \xrightarrow{\text{evoluzione}} U \ket{\psi} \qquad U \text{ operatore unitario}
\]
\[
\ket{\psi} \xrightarrow{\text{misura}} \ket{e_k} = N \, P_k \ket{\psi} \qquad P_k \text{ operatore di proiezione}
\]
Infatti gli operatori di proiezione non sono unitari perch\'e se esistesse l'inverso \(P^{-1}\) avrei \(P^2 = P\), cio\`e \(P = \mathbbm{1}\).
% \subsection{Trasformazioni unitarie sugli operatori}

Ho un operatore qualunque \(A_{ij} = \braket{e_i|A|e_j}\) e voglio trovare \(A_{ij}' = \braket{e_i'|A|e_j'}\). Prendo
\[
\bra{e_i} U = \bra{e_i} \sum_k \ket{e_k} \bra{e_k'} = \bra{e_i'}
\]
\[
U^\dag \ket{e_j} = \left( \sum_k \ket{e_k} \bra{e_k'} \right)^\dag \ket{e_j} = \ket{e_j'}
\]
quindi
\[
A' = U A U^\dag
\]

Così definisco gli operatori unitari equivalenti. Siano \(A\) e \(B\) operatori tali che \(B = U A U^\dag\) con \(U\) unitario. Allora per gli autovettori
\begin{align*}
A \ket{e_k} = \mu_k \ket{e_k} \quad &\Leftrightarrow \quad U A \ket{e_k} = \mu_k U \ket{e_k} \\
&\Leftrightarrow \quad U A U^\dag \left[ U \ket{e_k} \right] = \mu_k U \ket{e_k} \\
&\Leftrightarrow \quad B \ket{e_k'} = \mu_k \ket{e_k'} \\
&\Leftrightarrow \quad \ket{e_k'} = U \ket{e_k}
\end{align*}
% \subsection{Operatori compatibili e commutatore}
\(A_1\) e \(A_2\) hermitiani sono
\begin{enumerate}
\item compatibili se ammettono una base comune di autovettori;
\item incompatibili altrimenti.
\end{enumerate}
Teorema sul commutatore: \(A\) e \(B\) hermitiani sono compatibili se e solo se \(BA = AB\).
\[
\left[ A, B \right] = AB - BA \qquad \text{commutatore}
\]
\textbf{Dimostrazione.} Caso non degenere (\(\lambda_i \neq \lambda_j\)).

\bigskip
\noindent
\(\Leftarrow\) |

\bigskip
\noindent
\(A \ket{e_k} = \lambda_k \ket{e_k}\) sono autovettori e autovalori e
\begin{align*}
0 &= \Braket{e_k|\left[ A, B \right]|e_k} = \braket{e_k|AB|e_k} - \braket{e_k|BA|e_k} = \\
&= \lambda_k \left( \braket{e_k|B|e_k} - \braket{e_k|B|e_k} \right) = 0
\end{align*}
\[
0 = \braket{e_i|\left[ A, B \right]|e_k} = \braket{e_i|B|e_k} \left( \lambda_i - \lambda_k \right) \quad \implies \quad  \braket{e_i|B|e_k} = 0
\]
cioè \(B\) è diagonale con autovettore \(\ket{e_k}\), infatti
\[
B\ket{e_k} = \left( \sum_j \lambda_j \ket{e_j} \bra{e_j} \right) \ket{e_k} = \lambda_k \ket{e_k}
\]
\(\implies\) |
\[
\left( AB - BA \right) \ket{e_j} = \left( \lambda_i \mu_j - \lambda_i \mu_j \right) \ket{e_j} = 0
\]
\textbf{Dimostrazione.} Caso degenere

\bigskip
\noindent
\(\Leftarrow\) |
\[
A \ket{e_i} = \lambda_i \ket{e_i} \quad \text{con } \ket{e_i^a} \text{ tale che} 
\]
\[
A \ket{e_i^a} = \lambda_i \ket{e_i^a}
\]
\[
O = \braket{e_i^a|[A,B]|e_i^b} = (\lambda_i - \lambda_i) \braket{e_i^a|B|e_i^a} = 0
\]
Fabbrico un nuovo
\[
\ket{e_i'^c} = \sum_{a \text{ degeneri}} \ket{e_i^a}\braket{e_i^a|e_i'^c}
\]
e allora
\[
A \ket{e_i'^c} = \lambda_i \ket{e_i'^c}
\]
\[
B \ket{e_i'^c} = B \sum_{a \text{ degeneri}} \ket{e_i^a}\braket{e_i^a|e_i'^c} = \mu_i^c \ket{e_i'^c}
\]
\[
\braket{e_i'^d|B|e_i'^c} = \mu_i^c \braket{e_i'^d|e_i'^c}
\]
\[
B \ket{e_i'^c} = \mu_i^c \ket{e_i'^c} \qquad \square
\]

\subsubsection{Esempio: operatori unitari e generatori}
\(U = e^{iH}\) operatore unitario.
\[
U = \sum_{k = 0}^\infty \frac{1}{k!} \left( iH \right)^k \qquad \implies \qquad U^\dag = \sum_{k = 0}^\infty \frac{1}{k!} \left[ \left( iH \right)^k \right]^\dag
\]
Sapendo che \((z A)^\dag = z^* A^\dag\) e che
\[
AB \ket{\psi} = \left( B \ket{\psi} \right)^\dag A^\dag = \bra{\psi} B^\dag A^\dag
\]
risulta
\[
U^\dag = \sum_{k = 0}^\infty \frac{1}{k!} \left( -iH^\dag \right)^k = e^{-iH^\dag}
\]
Allora
\[
U \, U^\dag = \mathbbm{1}.
\]
% \subsection{Insieme completo di operatori}

Un insieme completo di operatori \`e un gruppo di operatori osservabili (cio\`e operatori hermitiani) tali che:
\begin{enumerate}
\item gli operatori commutano tra loro (\( [A, B] = 0 \) per ogni coppia di operatori \( A \) e \( B \) nell'insieme);
\item il loro insieme di autovalori (spettro combinato) determina completamente lo stato del sistema, fino a una fase globale.
\end{enumerate}
Questi insiemi sono anche chiamati insiemi completi di osservabili compatibili.
% \subsection{Deviazione standard dell'operatore}
\`E l'incertezza dell'operatore su \(\ket{\psi}\)
\begin{align*}
\Delta^2 A_\psi & \equiv \braket{\psi|A^2|\psi} - \left( \braket{\psi|A|\psi} \right)^2 = \\
& = \braket{\psi|\left( A - \braket{\psi|A|\psi} \mathbbm{1} \right)^2|\psi}
\end{align*}
\[
\Delta^2 A_\psi = 0 \quad \Leftrightarrow \quad A \ket{\psi} = \lambda \ket{\psi}
\]
\(\Leftarrow\) | 
\[
A^2 \ket{\psi} = \lambda^2 \ket{\psi}
\]
\[
\Delta^2 A = \braket{\psi|A^2|\psi} - \left( \braket{\psi|A|\psi} \right)^2 = \lambda^2 \braket{\psi|\psi} - \lambda^2 \braket{\psi|\psi} = 0
\]
\(\implies\) |
\begin{align*}
0 & = \braket{e_i|A^2|e_i} - \left( \braket{e_i|A|e_i} \right)^2 = \\
& = \sum_j \braket{e_i|A|e_j}\braket{e_j|A|e_i} - \left( \underbrace{\braket{e_i|A|e_i}}_{\text{matrice diagonale}} \right)^2 = \\
& = \sum_j \left| \braket{e_i|A|e_j} \right|^2 - \left| \braket{e_i|A|e_i} \right|^2 = \\
& = \sum_{j \ne i} \left| \braket{e_i|A|e_j} \right|^2
\end{align*}
Somma di numeri positivi, quindi \(\braket{e_i|A|e_j} = 0\). Quindi solo \(\braket{e_i|A|e_i}\) può essere diverso da 0, cioè \(\ket{e_i}\) è autostato. \(\square\)
% \subsection{Principio di indeterminazione}

Prendo due operatori $A$ e $B$ applicati allo stato $\ket{\psi}$. Allora
$$
\Delta^2 A_\psi \; \Delta^2 B_\psi \geq \frac{1}{4} \left| \braket{\psi|[A, B]|\psi} \right|^2
$$
1) Disuguaglianza di Schwarz
$$
\left| \braket{\psi|AB|\psi} \right|^2 \leq \braket{\psi|A^2|\psi} \braket{\psi|B^2|\psi}
$$
Dimostrazione: definisco $A\ket{\psi} = \ket{\alpha}$ e $B\ket{\psi} = \ket{\beta}$ tali che esiste $\ket{\delta}$ tale che
$$
\ket{\alpha} = z\ket{\beta} + \ket{\delta} \qquad \text{e} \qquad \braket{\beta|\delta} = 0
$$
Quindi
$$
\braket{\beta|\alpha} = z \braket{\beta|\beta} \quad \Rightarrow \quad z = \frac{\braket{\beta|\alpha}}{\braket{\beta|\beta}}
$$
Calcolo
\begin{align*}
\braket{\alpha|\alpha} & = |z|^2 \braket{\beta|\beta} + \braket{\delta|\delta} = \\
& = \frac{|\braket{\beta|\alpha}|^2}{(\braket{\beta|\beta})^2} \braket{\beta|\beta}  + \braket{\delta|\delta} = \\
& = \frac{|\braket{\beta|\alpha}|^2}{\braket{\beta|\beta}}  + \braket{\delta|\delta}
\end{align*}
Quindi
$$
\braket{\alpha|\alpha} \braket{\beta|\beta} = |\braket{\beta|\alpha}|^2 + \braket{\delta|\delta} \braket{\beta|\beta} \geq |\braket{\beta|\alpha}|^2 \qquad \square
$$
2)
$$
|\braket{\psi|AB|\psi}|^2 \geq \frac{1}{4} |\braket{\psi|[A,B]|\psi}|^2
$$
Dimostrazione:
$$
AB = \frac{1}{2} (AB - BA) + \frac{1}{2} (AB + BA) = \frac{1}{2} [A,B] + \frac{1}{2} \{A,B\}
$$
$$
z = \braket{\psi|AB|\psi} = \braket{\psi|\frac{1}{2}[A,B]|\psi} + \braket{\psi|\frac{1}{2}\{A,B\}|\psi} = x + iy
$$
$$
x = \braket{\psi|\frac{1}{2}\{A,B\}|\psi} \qquad \Rightarrow \qquad x^* = \dots
$$
$$
y = \braket{\psi|\frac{1}{2}[A,B]|\psi} \qquad \Rightarrow \qquad (iy)^* = \dots
$$
Quindi
$$
|z|^2 = x^2 + y^2 \geq y^2 = \frac{1}{4} |\braket{\psi|[A,B]|\psi}|^2 \qquad \square
$$

Le disuguaglianze 1) e 2) danno
$$
\begin{cases}
\braket{\psi|A^2|\psi} \braket{\psi|B^2|\psi} \geq \frac{1}{4} |\braket{\psi|[A,B]|\psi}|^2 \\
A = C - \braket{C} \mathbbm{1} \\
B = D - \braket{D} \mathbbm{1}
\end{cases} \qquad \Rightarrow
$$
$$
\Rightarrow \qquad
\begin{cases}
\braket{\psi|A^2|\psi} \braket{\psi|B^2|\psi} = \Delta^2 C_\psi \Delta^2 D_\psi \\
[A,B] = [C - \braket{C} \mathbbm{1}, D - \braket{D} \mathbbm{1}] = [C, D]
\end{cases} \qquad \Rightarrow
$$
$$
\Rightarrow \qquad
\Delta^2 C_\psi \; \Delta^2 D_\psi \geq \frac{1}{4} |\braket{\psi|[C,D]|\psi}|^2 \qquad \square
$$
% \subsection{Matrice densit\`a}

Sistema a due livelli
\[
\ket{\psi} = a_+ \ket{+} + a_- \ket{-}
\]
e costruisco
\[
\ket{\psi'} = e^{i \theta} \ket{\psi}
\]
Allora le probabilit\`a (fase globale)
\[
P_i = |\braket{e_i|\psi}|^2 = |\braket{e_i|\psi'}|^2
\]
fanno perdere informazione  sulla fase di $\ket{\psi'}$ rispetto a  $\ket{\psi}$ e anche per un generico operatore $A$
\[
\braket{\psi|A|\psi} = \braket{\psi'|A|\psi'}
\]

Ho una classe di equivalenze (\textit{raggio}) fra quelli che differiscono per la fase globale $\rightarrow$ sono lo stesso stato.

Parto dalla considerazione che se non posso ricostruire il sistema con le sue condizioni iniziali,  non ho informazioni sul sistema. Devo poter riprodurre pi\`u volte le codizioni iniziali.

Suppongo di avere un'osservabile e voglio capire quanta informazione sullo stato posso ricostruire partendo dalle caratteristiche dell'osservabile.

Da $A$ ottengo
\[
\begin{cases}
\braket{A} = \lambda_1 P_1 + \lambda_2 P_2 \\
P_1 + P_2 = 1 \\
P_1 = |a_1|^2 \\
P_2 = |a_2|^2
\end{cases} \qquad \Rightarrow \qquad P_1, P_2 \text{ deterministiche}
\]
Scrivo
\[
A = \lambda_1 \ket{1}\bra{1} + \lambda_2 \ket{2}\bra{2}
\]
\[
A' = \lambda_1' \ket{1}\bra{1} + \lambda_2' \ket{2}\bra{2}
\]
\[
\begin{cases}
\braket{A'} = \lambda_1' P_1 + \lambda_2' P_2 \\
P_1 + P_2 = 1 \\
\qquad \vdots
\end{cases} \qquad \Rightarrow \qquad P_1, P_2 \text{ deterministiche}
\]
Quindi due variabili compatibili danno la stessa informazione. Dovr\`o prenderne due incompatibili.

Introduco l'operatore densit\`a
\[
\rho_\psi = P_\psi = \ket{\psi}\bra{\psi} \qquad \text{stati puri}
\]
Allora
\[
\braket{A} = \Tr(A\rho)
\]
e ricordo che
\begin{enumerate}
\item \`e invariante per cambio di base,
\item $\Tr(A\rho) = \Tr(\rho A)$,
\item $\Tr(A) = \sum_i \mu_i$ autovalori.
\end{enumerate}
Dimostrazione che $\braket{A} = \Tr(A\rho)$
\[
\begin{cases}
\Tr(A\rho) = \sum_i \braket{e_i|A\rho|e_i} \\
A \ket{e_i} = \lambda_i \ket{e_i}
\end{cases}
\]
\[
\Tr(A\rho) = \sum_i \braket{e_i|A|\psi} \braket{\psi|e_i} = \sum_i \lambda_i \braket{e_i|\psi} \braket{\psi|e_i} = \sum_i \lambda_i |a_i|^2 = \braket{A}
\]

In particolare, l'operatore densit\`a \`e utile quando di un sistema conosco solo le probabilit\`a associate agli stati.
\[
p_i \longrightarrow \ket{\psi_i}
\]
\[
\braket{\psi_i|\psi_j} = 0
\]
\[
\braket{\psi_i|\psi_i} = 1
\]

Allora $\rho$ \`e definita da
\[
\rho = \sum_i p_i \ket{\psi_i} \bra{\psi_i} \qquad \text{stati misti}
\]
quindi
\[
\braket{\psi_i|\rho|\psi_i} = p_i \qquad \braket{\psi_i|\rho|\psi_j} = 0
\]
inoltre, poich\'e $\braket{A} = \Tr(A\rho)$, vale
\begin{align*}
\Tr(A\rho) &= \sum_i \braket{\psi_i|A\rho|\psi_i} = \\
&= \sum_i \braket{\psi_i|A \left( \sum_j p_j \ket{\psi_j} \bra{\psi_j} \right)|\psi_i} = \\
&= \sum_i \braket{\psi_i|A p_i|\psi_i} = \\
&= \sum_i p_i \braket{\psi_i|A|\psi_i} \qquad \Rightarrow
\end{align*}
\[
\Rightarrow \qquad \braket{A} = \sum_i p_i \braket{\psi_i|A|\psi_i}
\]
In generale, l'operatore densit\`a associato allo stato puro  \`e diverso da quello associato allo stato misto: nel primo sopravvivono i termini di interferenza, nel secondo no.

\subsubsection{Esempio: matrice densit\`a e stati}
\[
\ket{\psi_1} = \frac{1}{\sqrt{2}} (\ket{+} + \ket{-})
\]
\[
\rho_{\psi_1} = \ket{\psi_1} \bra{\psi_1} \qquad \text{stato puro}
\]
\[
A = \begin{pmatrix}
0 & 1 \\
1 & 0
\end{pmatrix} \quad \Rightarrow \quad A \ket{\psi_1} = \ket{\psi_1} \quad \Rightarrow \quad \braket{A}_{\psi_1} = \braket{\psi_1|A|\psi_1}
\]
Se invece ho lo stato misto, in cui
\[
p_{\ket{+}} = \frac{1}{2} \qquad p_{\ket{-}} = \frac{1}{2}
\]
\[
\rho_m = \frac{1}{2} \ket{+} \bra{+} + \frac{1}{2} \ket{-} \bra{-}
\]
\[
\braket{A}_m = \Tr(A\rho_m) = \frac{1}{2} \braket{+|A|+} + \frac{1}{2} \braket{-|A|-} = 0
\]

In particolare, per gli stati puri ho
\[
\rho^2 = \ket{\psi} \braket{\psi|\psi} \bra{\psi} = \ket{\psi} \bra{\psi} = \rho
\]
ma per lo stato misto, nella base degli autostati
\begin{align*}
\rho^2 &= \left( \sum_i p_i \ket{\psi_i} \bra{\psi_i} \right) \left( \sum_j p_j \ket{\psi_j} \bra{\psi_j} \right) = \\
&= \sum_i \left( \sum_j p_i p_j \ket{\psi_i} \braket{\psi_i|\psi_j} \bra{\psi_j} \right) = \\
&= \sum_i p_i^2 \ket{\psi_i} \bra{\psi_i} \ne \rho
\end{align*}
Inoltre, $\Tr(\rho) = \sum_i p_i = 1$, ma
\[
\Tr \rho = \Tr \rho^2 = \begin{cases}
1 \qquad \qquad  \qquad\text{stati puri} \\
\sum_i p_i^2 < 1 \qquad \text{stati misti}
\end{cases}
\]
allora per uno stato puro $p_{i_0} = 1$ e $p_i = 0$ per ogni $i \ne i_0$ e poich\'e vale $\det \rho = \prod_i p_i$,  nel caso bipartito ho
\[
\det \rho \begin{cases}
= 0 \qquad \text{puro} \\
\ne 0 \qquad \text{misto}
\end{cases}
\]

\subsubsection{Esempio: ancora matrice densit\`a}
\[
\ket{\psi} = N \left( \left( 1 + \sqrt{2} \right) i \ket{+} + \ket{-} \right)
\]
\[
\rho_\psi = \ket{\psi} \bra{\psi} = |N|^2 \begin{pmatrix}
\left( 1 + \sqrt{2} \right)^2 & \left( 1 + \sqrt{2} \right) i \\
-\left( 1 + \sqrt{2} \right) i & 1
\end{pmatrix}
\]
\[
\Tr(\rho_\psi) = |N|^2 \left[ \left( 1 + \sqrt{2} \right) + 1 \right] = 1 \quad \Rightarrow \quad |N|^2 = \frac{1}{2(2 + \sqrt{2})}
\]
\[
\det \rho_\psi = 0
\]
\[
B = \begin{pmatrix}
1 & 0 \\
0 & -1
\end{pmatrix} \qquad B \ket{\psi} = N \left[ \left( 1 + \sqrt{2} \right) i \ket{+} - \ket{-} \right]
\]
\[
C = \begin{pmatrix}
0 & i \\
-i & 0
\end{pmatrix} = \sigma_2
\]
\[
\rho_\psi = \frac{1}{2} \left( \mathbbm{1} + \sum_i p_i \sigma_i \right) \qquad p_1 = p_3 = \frac{1}{\sqrt{2}}
\]
\[
\mathrm{misura} \Rightarrow \begin{cases}
\ket{+}\bra{+} = \begin{pmatrix} 1 & 0 \\ 0 & 0 \end{pmatrix} \\
\ket{-}\bra{-} = \begin{pmatrix} 0 & 0 \\ 0 & 1 \end{pmatrix}
\end{cases}.
\]
% \subsection{Caratterizzare gli stati attraverso la matrice densit\`a}

Per una matrice $2 \times 2$ hermitiana vale
$$
M = a \mathbbm{1} + \vec{b} \cdot \vec{\sigma} = \begin{pmatrix}
a + b_3 & b_1 - ib_2 \\
b_1 + ib_2 & a - b_3
\end{pmatrix}
$$
e poich\'e una matrice densit\`a ha $\Tr M = 1 = 2a$
$$
\rho = \frac{1}{2} (\mathbbm{1} + \vec{p} \cdot \vec{\sigma})
$$
con $\sigma_i$ Matrici di Pauli
$$
\sigma_1 = \begin{pmatrix}
0 & 1 \\
1 & 0
\end{pmatrix} \qquad
\sigma_2 = \begin{pmatrix}
0 & -i \\
i & 0
\end{pmatrix} \qquad
\sigma_3 = \begin{pmatrix}
1 & 0 \\
0 & -1
\end{pmatrix}
$$
dove $\vec{p}$ \`e tale che
$$
\det \rho = \frac{1}{4} (1 - p_3^2 + p_1^2 -p_2^2) = \frac{1}{4} (1 - ||\vec{p}||^2) = 0 \; \Rightarrow \;  ||\vec{p}|| = 1
$$

Una propriet\`a interessante della matrici di Pauli \`e
$$
\Tr(\sigma_i \sigma_j) = 2 \delta_{ij}
$$
quindi
$$
\Tr(\rho \sigma_i) = \Tr \left( \frac{1}{2} \mathbbm{1} \right) + \frac{1}{2} \Tr \left[ (\vec{p} \cdot \vec{\sigma}) \sigma_i \right] \quad \Rightarrow \quad \braket{\sigma_i}_\psi = \Tr(\rho_\psi \sigma_i)
$$

% \subsection{Clonazione degli stati}

Ho due stati $\ket{\psi}$ e $\ket{\varphi}$ e produco $\ket{\psi} \otimes \ket{x}$. Esiste un operatore unitario $U$ tale che
\[
U \ket{\psi} \otimes \ket{x} = \ket{\psi} \otimes \ket{\psi} \text{?}
\]

Se esistesse, allora
\[
\ket{\alpha} = U \ket{\psi} \otimes \ket{x}
\]
\[
\ket{\beta} = U \ket{\varphi} \otimes \ket{x}
\]
\[
\braket{\beta|\alpha} = \braket{\varphi|\psi} \braket{x|x} = \braket{\varphi|\psi} \braket{\varphi|\psi} \; \Rightarrow \; \braket{\varphi|\psi} = \braket{\varphi|\psi}^2 \; \Rightarrow \; \braket{\varphi|\psi} = \begin{cases}
0 \\
1
\end{cases}
\]
cio\`e o $\ket{\varphi} \equiv \ket{\psi}$ o $ \ket{\varphi} \perp \ket{\psi}$.
% \subsection{Operatore posizione}
Il problema è decomporre \(\ket{\psi}\) su una base che ne dia la posizione. Allora introduco
\[
\hat{x} \ket{x} = x \ket{x}
\]
che in generale dà
\[
\braket{x|\psi} = \psi(x) \qquad \text{funzione d'onda}
\]
dove \(|\psi(x)|^2 = |\braket{x|\psi}|^2\). Dato che \(x\) è un parametro reale continuo, vorrei scrivere
\[
\ket{\psi} = \int \psi(x) \ket{x} \, dx
\]
come estensione del caso discreto. Così \(|\psi(x)|^2\) è una distribuzione di probabilità \(\longrightarrow\) una densità di probabilità del tipo
\[
P_\Delta \left( [x, x + \Delta x] \right) = \int_x^{x + \Delta x} |\psi(x)|^2 \, dx
\]
con la condizione che
\[
\int_A^B |\psi(x)|^2 \, dx = 1
\]
che esteso a tutto lo spazio diventa
\[
\int_{-\infty}^{+\infty} |\psi(x)|^2 \, dx = 1
\]
I problemi sorgono nell'esprimere la condizione di ortonormalità, perché
\[
\begin{cases}
\ket{\psi} = \int_{-\infty}^{+\infty} \ket{x} \braket{x|\psi} \, dx \\
\braket{x'|\psi} = \psi(x')
\end{cases} \quad \Rightarrow \quad \psi(x') = \int_{-\infty}^{+\infty} \braket{x'|x} \braket{x|\psi} \, dx \quad \Rightarrow
\]
\[
\Rightarrow \quad \psi(x') = \int_{-\infty}^{+\infty} \braket{x'|x} \psi(x) \, dx
\]
Cioè ho introdotto un oggetto tale che
\[
\braket{x'|x} = \delta(x',x)
\]
che ha la proprietà per definizione
\[
\int_{-\infty}^{+\infty} \psi(x) \delta(x, x') \, dx = \psi(x')
\]
Da questa discende:
\begin{enumerate}
\item invarianza per traslazione
\[
\delta(x', x) = \delta(x' - x)
\]
Poiché \(\delta(x' - x) = \delta(x' + a, x + a)\)
\[
\int_{-\infty}^{+\infty} \delta(x' + a, x + a) \psi(x) \, dx = [x'' = x + a] =
\]
\[
= \int_{-\infty}^{+\infty} \delta(x' + a, x'') \psi(x'' - a) \, dx'' = \psi(x' + a - a) = \psi(x')
\]
\item simmetria
\[
\delta(x' - x)  = \delta(x - x')
\]
Infatti
\[
\braket{\psi|x'} = \int_{-\infty}^{+\infty} \braket{\psi|x}\braket{x|x'} \, dx = \int_{-\infty}^{+\infty} \psi^*(x) \delta(x,x') \, dx = \psi^*(x')
\]
\item \(f(x) \delta(x) = f(0) \delta(x)\)
\begin{align*}
\int_{-\infty}^{+\infty} f(x) \delta(x) \psi(x) \, dx &= f(0) \psi(0) = \\
&= f(0) \int_{-\infty}^{+\infty} \psi(x) \delta(x) \, dx = \\
&= \int_{-\infty}^{+\infty} f(0) \psi(x) \delta(x) \, dx
\end{align*}
Poiché \(\psi\) è generica, \(f(0) \delta(x) = f(x) \delta(x)\).
\item \(\delta(ax) = \frac{1}{a} \delta(x)\)
\[
\int_{-\infty}^{+\infty} \delta(ax) \psi(x) \, dx = \left[ ax = t \; \Rightarrow \; dx = \frac{dt}{a} \right] = 
\]
\[
= \frac{1}{a} \int_{-\infty}^{+\infty} \delta(t) \psi \left( \frac{t}{a} \right) \, dt = \frac{1}{a} \psi(0) = \frac{1}{a} \int_{-\infty}^{+\infty} \delta(x) \psi(x) \, dx 
\]
\item normalizzazione
\[
\int_{-\infty}^{+\infty} \delta(x) \, dx = 1
\]
\[
\int_{-\infty}^{+\infty} \delta(x) \cdot 1 \, dx = 1
\]
\end{enumerate}
\textbf{Relazione di ortonormalizzazione}
\[
\ket{\psi} = \int_{-\infty}^{+\infty} \ket{x}\braket{x|\psi} \, dx
\]
Nel caso continuo \(\braket{x|x'} = \delta(x - x') \longrightarrow\) ortogonalizzazione impropria, infatti
\[
\begin{cases}
\braket{x|x} = \delta(0) \\
\int_{-\infty}^{+\infty} |\psi(x)|^2 \, dx = 1
\end{cases}
\]
\textbf{Risoluzione dell'identità}
\[
\int_{-\infty}^{+\infty} \ket{x} \bra{x} \, dx = \mathbbm{1}
\]
il cui elemento
\[
\Braket{x'|\int_{-\infty}^{+\infty} \Ket{x} \Bra{x} dx\;|x''} = \int_{-\infty}^{+\infty} \delta(x' - x) \delta(x - x'') \, dx = \delta(x' - x'') = \braket{x'|x''}
\]
\textbf{Prodotto scalare fra stati}
\[
\braket{\varphi|\psi} = \int_{+\infty}^{-\infty} \braket{\varphi|x}\braket{x|\psi} \, dx = \int_{-\infty}^{+\infty} \varphi(x) \psi(x) \, dx
\]
\[
\braket{\psi|\psi} = \int_{-\infty}^{+\infty} |\psi(x)|^2 \, dx
\]
\textbf{Operatori}
\[
A(\hat{x}) = \sum_{i = 0}^\infty a_i \, (\hat{x})^i
\]
\[
\braket{x'|A(\hat{x})|x''} = \sum_i a_i \, (\hat{x})^i \braket{x'|x''} = \delta(x' - x'') \sum_i a_i \, (\hat{x})^i = A(x'') \delta(x' - x'')
\]
Ad es.
\[
\braket{y|\cos \hat{x}|x} = \cos z \; \delta(y - x) = \cos y \; \delta(y - x)
\]
\begin{align*}
\braket{\psi|\cos \hat{x}| \psi} &= \int_{-\infty}^{+\infty} \int_{-\infty}^{+\infty} \braket{\psi|y} \braket{y|\cos \hat{x}|z} \braket{z|\psi} \, dz \, dy = \\
&= \int_{-\infty}^{+\infty} \int_{-\infty}^{+\infty} \psi^*(y) \, \cos z \, \delta(y - z) \, \psi(z) \, dy \, dz = \\
&= \int_{-\infty}^{+\infty} |\psi(z)|^2 \, \cos z \, dz
\end{align*}
\textbf{Riepilogo}
\begin{enumerate}
\item \(\braket{q'|\hat{q}|q} = q \braket{q'|q} = q \delta(q' - q)\)
\item \(\braket{q'|f(\hat{q})|q} = f(q) \braket{q|q'} = f(q) \, \delta(q - q')\)
\end{enumerate}

\subsubsection{Esempio: ulteriori proprietà della \(\delta\)}
\[
\int_{-\infty}^{+\infty} \delta(f(x)) \psi(x) \, dx
\]
\begin{align*}
\int_{-\infty}^{+\infty} \delta(f(x)) \psi(x) \, dx &= \left[ x^\prime = f(x), dx^\prime = f^\prime(x) \, dx, x = f^{-1}(x) \right] = \\
&= \int_{-\infty}^{+\infty} \frac{\delta^\prime(x)}{f^\prime(x)} \psi(f^{-1}(x)) \, dx^\prime = \\
\int_{-\infty}^{+\infty} \delta(f(x)) \psi(x) \, dx &= \sum_i \int_{a_i}^{b_i} \delta(f(x)) \psi(x) \, dx \\ 
\int_{-\infty}^{+\infty} \delta(f(x)) \psi(x) \, dx &= \frac{\psi(x_i)}{f^\prime(x)|_{x : f(x) = 0}} \psi(f^{-1})
\end{align*}

\subsubsection{Esempio: operatore parità}
\( \braket{x|\mathcal{P}|\psi} = \psi(-x) \) con \( \braket{x|\psi} = \psi(x) \Rightarrow \bra{x} \mathcal{P} = \bra{-x} \).
\begin{align*}
(\braket{\psi|\mathcal{P}|x})^* &= \left( \int_{-\infty}^{+\infty} \braket{\psi|x'} \braket{x'|\mathcal{P}|x} \, dx' \right)^* = \\
&= \left( \int_{-\infty}^{+\infty} \psi^*(x') \delta(-x' - x) \, dx' \right)^* = \\
&= \int_{-\infty}^{+\infty} \psi(x') \delta(-x' - x) \, dx' = \\
&= \psi(-x)
\end{align*}
quindi \( \mathcal{P} \) è hermitiano \( \Rightarrow \mathcal{P} \ket{x} = \ket{-x} \Rightarrow \mathcal{P}^2 \ket{x} = \ket{x} \Rightarrow \mathcal{P}^2 = \mathbbm{1} \Rightarrow \mathcal{P} \mathcal{P} = \mathcal{P}^\dag \mathcal{P} = \mathbbm{1} \Rightarrow \mathcal{P} \) è unitario. \( \mathcal{P}^2 = \mathbbm{1} \Rightarrow \) autovalori \( \pm 1 \).


\bigskip
\noindent
Le autofunzioni sono
\[
\braket{x|\mathcal{P}|\varphi_{-}} = \begin{cases}
- \varphi_{-}(x) \\
\varphi_{-}(-x)
\end{cases} \quad \Rightarrow \quad \varphi_{-} \text{ funzione dispari}
\]
\[
\braket{x|\mathcal{P}|\varphi_{+}} = \begin{cases}
\varphi_{+}(x) \\
\varphi_{+}(-x)
\end{cases} \quad \Rightarrow \quad \varphi_{+} \text{ funzione pari}
\]
\[
f(x) = \frac{f(x) + f(-x)}{2} + \frac{f(x) - f(-x)}{2}.
\]
% \subsection{Quantizzazione di un sistema}
Il trucco \`e definire le grandezze attraverso le leggi di conservazione.

Dal formalismo lagrangiano, suppongo che il sistema ha lagrangiana, allora
$$
\delta \mathcal{L} = 0 \qquad \Rightarrow \qquad Q = \frac{\partial \mathcal{L}}{\partial \dot{q}} \, \delta q \; \text{(carica di Noether).}
$$

Infatti l'invarianza della lagrangiana \`e
\begin{align*}
0 = \delta \mathcal{L} &= \frac{\partial \mathcal{L}}{\partial q} \delta q + \frac{\partial \mathcal{L}}{\partial \dot{q}} \delta \dot{q} = \\
&= \frac{d}{dt} \frac{\partial \mathcal{L}}{\partial \dot{q}} \delta q + \frac{\partial \mathcal{L}}{\partial \dot{q}} \delta \dot{q} = \\
&= \frac{d}{dt} \frac{\partial \mathcal{L}}{\partial \dot{q}} \delta q + \frac{\partial \mathcal{L}}{\partial \dot{q}} \frac{d}{dt} \delta q = \\
&= \frac{d}{dt} \left( \frac{\partial \mathcal{L}}{\partial \dot{q}} \delta q \right)
\end{align*}
quindi $ \frac{\partial \mathcal{L}}{\partial \dot{q}} \delta q $ \`e costante. Questo \`e il \textbf{teorema di Noether}.

Per le traslazioni, dal punto di vista quantistico ho un operatore unitario che agisce sui vettori di base
$$
T \ket{q} = \ket{q - \delta} \qquad (T^\dag T = \mathbbm{1})
$$
producendo una nuova base.

Per lo stato $ \ket{\psi} $
\begin{align*}
T \ket{\psi} &= \int_{-\infty}^{+\infty} T \ket{q} \braket{q|\psi} \, dq = \int_{-\infty}^{+\infty} \ket{q - \delta} \braket{q|\psi} \, dq = \\
&= \int_{-\infty}^{+\infty} \ket{q - \delta} \psi(q) \, dq = \int_{-\infty}^{+\infty} \ket{q^\prime} \psi(q^\prime + \delta) \, dq^\prime
\end{align*}
Allora
$$
\bra{q} T \ket{\psi} = \int_{-\infty}^{+\infty} \braket{q|q^\prime} \psi(q^\prime + \delta) \, dq^\prime = \int_{-\infty}^{+\infty} \delta(q^\prime - q) \, \psi(q^\prime + \delta) \, dq^\prime = \psi(q + \delta)
$$
che sviluppata \`e
\begin{align*}
\psi(q + \delta) &= \sum_{k = 0}^\infty \delta^k \frac{1}{k!} \psi^{(k)}(q) = \\
&= \left[ \sum_{k = 0}^\infty \frac{\delta^k}{k!} \left( \frac{d}{dq} \right)^k \right] \psi(q) = \\
&= e^{\delta \frac{d}{dq}} \psi(q)
\end{align*}
Inoltre so che $T_\delta = e^{i \hat{K} \delta}$, dove $\hat{K}^\dag = \hat{K}$, quindi
$$
\psi(q + \delta) = e^{i \delta \left( -i \frac{d}{dq} \right)} \psi(q).
$$

Per trasformazioni infinitesime $ T_\varepsilon = e^{i \varepsilon \hat{K}} = 1 + i \varepsilon \hat{K} + \sigma(\varepsilon^2) $, dove $ \hat{K} $ \`e il \textbf{generatore delle traslazioni}. Allora
$$
\bra{q} 1 + i \delta \hat{K} \ket{\psi} =  \bra{q} e^{i \delta \hat{K}} \ket{\psi} = e^{i \delta \left( -i \frac{d}{dq} \right)} \psi(q) = \psi(q) + i \delta \left( -i \psi^\prime(q) \right)
$$
$$
\boxed{
\braket{q|\hat{K}|\psi} = -i \frac{d}{dq} \psi(q).
}
$$

Ma cosa si conserva di un sistema quando c'\`e invarianza per traslazione?

Supponendo che esista $ S $ tale che $ S^\dag \, S = \mathbbm{1} $ e $ S(t, t_0) |\ket{\psi} = \ket{\psi(t)} $ e quindi $ \braket{\psi(t_0)|\psi(t_0)} = \braket{\psi(t)|\psi(t)} $; avendo invarianza per traslazione per due stati $ \ket{\varphi} $ e $ \ket{\psi} $, segue
$$
\braket{\varphi(t)|\psi(t)} = \braket{\varphi(t)|S(t, t_0)|\psi(t_0)} = \braket{\varphi(t)|T^\dag S(t, t_0) T|\psi(t_0)}
$$
che nel caso infinitesimo \`e
$$
\braket{\varphi(t)|\left( \mathbbm{1} - i \varepsilon \hat{K} \right) \, S(t, t_0) \, \left( \mathbbm{1} + i \varepsilon \hat{K} \right) | \psi(t_0)} =
$$
$$
= \braket{\varphi(t)|S(t, t_0) - i \varepsilon \hat{K} S(t, t_0) + i \varepsilon \hat{K} S(t, t_0)|\psi(t_0)} =
$$
$$
= \braket{\varphi(t)|S(t, t_0)|\psi(t_0)} - i \varepsilon \braket{\varphi(t)|\left[ \hat{K}, S(t, t_0) \right]|\psi(t_0)}
$$
Quindi deve essere $ \left[ \hat{K}, S(t, t_0) \right] = 0 $.

Gli autostati di  $ \hat{K} $ ci sono e $ \hat{K} \ket{k} = k \ket{k} $, quindi
\begin{align*}
0 &= \braket{\varphi(t)|\left[ \hat{K}, S(t, t_0) \right]|\psi(t_0)} = \int_{-\infty}^{+\infty} \braket{\varphi|k} \braket{k|[,]|k^\prime} \braket{k^\prime|\psi} dk \, dk^\prime \\
&= \int_{-\infty}^{+\infty} \braket{\varphi|k} \left( k - k^\prime \right) \braket{k|S(t, t_0)|k^\prime} \braket{k^\prime|\psi} dk \, dk^\prime
\end{align*}
e questo per ogni $\ket{\psi}, \ket{\phi}$, quindi
$$
\text{se } k \neq k^\prime, \text{allora } \braket{k|S(t, t_0)|k^\prime} = 0
$$
cio\`e l'autovalore di $ \hat{K} $ si conserva.

Definisco l'\textbf{operatore impulso} o \textbf{momento}
$$
\hat{p} = \hbar \hat{K}.
$$

\subsubsection{Esempio: commutatore posizione-momento}
$$
\hat{q} \xrightarrow{\text{evoluzione}} T^{-1} \hat{q} \, T = \hat{q} - \delta.
$$
\begin{align*}
T^{-1} \hat{q} \, T \ket{q} &= T^{-1} \hat{q} \ket{q - \delta} = T^{-1} \left( q - \delta \right) \ket{q - \delta} = \\
&= \left( q - \delta \right) T^{-1} \ket{q - \delta} = \left( q - \delta \right) \ket{q}.
\end{align*}
Per trasfromazioni infinitesimali $ T_\varepsilon = \mathbbm{1} + i \varepsilon \hat{K} = \mathbbm{1} + i \varepsilon \frac{1}{\hbar} \hat{p} $, quindi
$$
\left( \mathbbm{1} - i \varepsilon \frac{1}{\hbar} \hat{p} \right) \, \hat{q} \, \left( \mathbbm{1} + i \varepsilon \frac{1}{\hbar} \hat{p} \right) = \hat{q} - \varepsilon
$$
$$
\hat{q} - i \varepsilon \frac{1}{\hbar} \left( \hat{p} \hat{q} - \hat{q} \hat{p} \right) = \hat{q} - \varepsilon
$$
quindi  $ \left[ \hat{p}, \hat{q} \right] = -i \hbar $.

\subsubsection{Esempio: ancora sull'operatore traslazione}
$$
T = \begin{cases}
\mathbbm{1} + i \frac{\varepsilon}{\hbar} \hat{p} = T_\varepsilon \text{ infinitesimo} \\
e^{i \frac{\varepsilon}{\hbar} \hat{p}} = T_\varepsilon
\end{cases}
$$
\begin{align*}
\braket{k'|T_\varepsilon|k} &= \braket{k'|\mathbbm{1} + i \frac{\varepsilon}{\hbar} \hat{p}|k} = \delta(k' - k) (1 + i \varepsilon k), \\
\braket{k'|T_\varepsilon|k} &= e^{i \varepsilon k} \delta(k' - k), \\
\braket{k'|\hat{q}|k} &= - i \frac{d}{dk} \delta(k - k') = i \frac{d}{dk'} \delta(k' - k).
\end{align*}
% \subsection{Elementi di matrice}
\[
\braket{q^\prime|\hat{p}|q} = \text{?}, \qquad \braket{q^\prime|\hat{q}|q} = \text{?}
\]
Poinch\'e \( \braket{q^\prime|q} = \delta(q - q^\prime) = \delta(q^\prime - q) \), esprimo
\[
\ket{q} \equiv \ket{\psi_q} \quad \Rightarrow \quad \braket{q^\prime|\psi_q} = \delta(q - q^\prime) = \psi_q(q^\prime)
\]
da cui
\[
\braket{q^\prime|\hat{p}|q} = -i \hbar \frac{d}{dq^\prime} \delta(q - q^\prime) = i \hbar \frac{d}{dq} \delta(q - q^\prime)
\]
\[
\braket{q^\prime|\hat{q}|q} = q \, \delta(q - q^\prime) = q^\prime \, \delta(q - q^\prime).
\]

Verifico \( \hat{p} \) e \( \hat{q} \) hermitiani.
\begin{align*}
\braket{q^\prime|\hat{p}|q} &= q \, \delta(q - q^\prime) = q^\prime \, \delta(q^\prime - q) = q^\prime \delta(q^\prime - q) = \\
&= \braket{q|\hat{q}|q^\prime} = \left( \braket{q^\prime|\hat{q}|q} \right)^*
\end{align*}
\begin{align*}
\braket{q^\prime|\hat{p}|q} &= -i \hbar \frac{d}{dq^\prime} \delta(q - q^\prime) = i \hbar \frac{d}{dq} \delta(q - q^\prime) = \\
&= i \hbar \frac{d}{dq} \delta(q^\prime - q) = \left( -i \hbar \frac{d}{dq} \delta(q^\prime - q) \right)^* = \\
&= \left( \braket{q|\hat{p}|q^\prime} \right)^*.
\end{align*}

Allora per \( \ket{\psi} \) e \( \ket{\varphi} \) ho \( \braket{q|\psi} = \psi(q) \) e \( \braket{q|\varphi} = \varphi(q) \).
\[
\braket{\varphi|\hat{p}|\psi} = \int_{-\infty}^{+\infty} \braket{\varphi|q} \braket{q|\hat{p}|\psi} \, dq = \int_{-\infty}^{+\infty} \varphi^*(q) (-i \hbar) \frac{d}{dq} \psi(q) \, dq,
\]
ma \( \braket{\varphi|\hat{p}|\psi} = (\braket{\psi|\hat{p}|\varphi})^* \), cio\`e
\[
\left( \braket{\psi|\hat{p}|\varphi} \right)^* = \left( \int_{-\infty}^{+\infty} \psi^*(q) (-i \hbar) \frac{d}{dq} \varphi(q) \right)^* dq = \int_{-\infty}^{+\infty} \psi(q) i \hbar \frac{d}{dq} \varphi^*(q) \, dq.
\]
Quindi
\[
\int_{-\infty}^{+\infty} \varphi^*(q) (-i \hbar) \frac{d}{dq} \psi(q) \, dq = \int_{-\infty}^{+\infty} \psi(q) i \hbar \frac{d}{dq} \varphi^*(q) \, dq,
\]
che integrando per parti d\`a
\[
i \hbar \int_{-\infty}^{+\infty} \varphi^*(q) \, \psi(q) \, dq = 0
\]
allora devo imporre che \( f \rightarrow 0 \) e \( g \rightarrow 0 \), cio\`e in particolare devo
\[
\int_{-\infty}^{+\infty} \left| \psi(q) \right|^2 \, dq < + \infty.
\]
% \subsection{Autofunzioni degli operatori}
So che \( \hat{q} \ket{q} = q \ket{q} \), dunque
\[
\braket{q^\prime|q} = \psi_q(q^\prime) = \delta(q - q^\prime) \; \rightarrow \; \text{autofunzione della posizione.}
\]

Per l'impulso voglio che esista \( \ket{k} \) tale che \( \hat{p} \ket{k} = \hbar k \ket{k} \). Scelgo la base dei \( \ket{q} \) e faccio
\[
\braket{q|\hat{p}|k} = \hbar k \braket{q|k}, \quad \text{con } \braket{q|k} = \psi_k(q) \quad \Rightarrow
\]
\[
\Rightarrow \quad -i \hbar \frac{d}{dq} \psi_k(q) = \hbar k \psi_k(q) \quad \Rightarrow \quad \psi_k(q) = N_k e^{ikq}.
\]
Poiché
\[
\int_{-\infty}^{+\infty} \left| \psi_k(q) \right|^2 \, dq = |N_k|^2 \int_{-\infty}^{+\infty} 1 \, dq,
\]
non posso normalizzare in senso proprio. Coerentemente con le autofunzioni della posizione, ho \( \braket{k|k^\prime} = \delta(k - k^\prime) \), cioè
\begin{align*}
\braket{k|k^\prime} &= \int_{-\infty}^{+\infty} \braket{k^\prime|q} \braket{q|k} \, dq = |N_k|^2 \int_{-\infty}^{+\infty} e^{-i k^\prime q} e^{ikq} \, dq = \\
&= |N_k|^2 \int_{-\infty}^{+\infty} e^{i (k - k^\prime) q} \, dq = \delta(k^\prime - k).
\end{align*}
Allora
\begin{align*}
1 &= \int_{-\infty}^{+\infty} \braket{k|k'} \, dk = \int_{-\infty}^{+\infty} \delta(k - k') \, dk = \\
&= \int_{\mathbb{R}^2} |N_k|^2 e^{i (k - k^\prime) q} \, dq \, dk = \\
&= \lim_{\lambda \rightarrow +\infty} \int_{-\infty}^{+\infty} |N_k|^2 \left[ \frac{e^{i (k - k^\prime) q}}{i (k - k^\prime)} \right]_{-\lambda}^{+\lambda} \, dk = \\
&= \lim_{\lambda \rightarrow +\infty} \int_{-\infty}^{+\infty} |N_k|^2 \frac{e^{i (k - k^\prime) \lambda} - e^{-i (k - k^\prime) \lambda}}{i (k - k^\prime)} \, dk = \\
&= 2 \pi |N_k|^2 \qquad \Rightarrow \qquad |N_k|^2 = \frac{1}{2\pi}
\end{align*}
\[
\psi_k(q) = \frac{1}{\sqrt{2\pi}} e^{ikq}
\]
e ridefinisco
\[
\delta(x) = \frac{1}{2\pi} \int_{-\infty}^{+\infty} e^{ipx} \, dp.
\]
Dalla definizione di aggiunto,
\[
\braket{q|k} = \frac{1}{\sqrt{2\pi}} e^{ikq}, \qquad \braket{k|q} = \frac{1}{\sqrt{2\pi}} e^{-ikq}
\]
quindi analogo ad \( U^\dag U = \mathbbm{1} \), cioè matrici di cambio di base
\[
\braket{q|\psi} = \psi(q), \qquad \braket{k|\psi} = \tilde{\psi}(k)
\]
legate da
\begin{align*}
\braket{k|\psi} &= \tilde{\psi}(k) = \int_{-\infty}^{+\infty} \braket{k|q} \braket{q|\psi} \, dq = \\
&= \int_{-\infty}^{+\infty} \frac{e^{-ikq}}{\sqrt{2\pi}} \psi(q) \, dq \\
\braket{q|\psi} &= \int_{-\infty}^{+\infty} \braket{q|k} \braket{k|\psi} \, dk = \int_{-\infty}^{+\infty} \frac{e^{ikq}}{\sqrt{2\pi}} \tilde{\psi}(k) \, dk.
\end{align*}
N.B.: in generale, \( \tilde{\psi} \neq \psi \).

Poiché
\[
\psi(q) = \braket{q|\psi} = \int_{-\infty}^{+\infty} \braket{q|k} \braket{k|\psi} \, dk = \frac{1}{\sqrt{2\pi}} \int_{-\infty}^{+\infty} e^{-ikq} \psi(k) \, dk
\]
\[
\psi(k) = \braket{k|\psi} = \int_{-\infty}^{+\infty} \braket{k|q} \braket{q|\psi} \, dq = \frac{1}{\sqrt{2\pi}} \int_{-\infty}^{+\infty} e^{ikq} \psi(q) \, dq
\]
ottengo che \( \ket{q} \) e \( \ket{k} \) sono \textbf{linearmente dipendenti}.

Gli elementi di matrice
\[
\braket{k|\hat{p}|\psi} = \hbar k \psi(k).
\]
N.B.: \( \hat{p}\ket{p} = p \ket{p} \Leftrightarrow \braket{p|\hat{p}|\psi}  = p \psi(p) \) con \( \psi(p) = \braket{p|\psi} \), ma allora
\[
\braket{q|p} = \frac{1}{\sqrt{2\pi}} e^{i \frac{pq}{\hbar}}.
\]

\begin{align*}
\braket{k|\hat{q}|\psi} &= \int_{-\infty}^{+\infty} \braket{k|\hat{q}|q} \braket{q|\psi} \, dq = \\
&= \int_{-\infty}^{+\infty} q \braket{k|q} \psi(q) \, dq = \\
&= \int_{-\infty}^{+\infty} q \frac{e^{-ikq}}{\sqrt{2\pi}} \psi(q) \, dq = \frac{1}{\sqrt{2\pi}} \int_{-\infty}^{+\infty} i \frac{d}{dk} e^{-ikq} \psi(q) \, dq = \\
&= i \frac{1}{\sqrt{2\pi}} \frac{d}{dk} \underbrace{\int_{-\infty}^{+\infty} e^{-ikq} \psi(q) \, dq}_{\psi(k)} = i \frac{1}{\sqrt{2\pi}}\frac{d}{dk} \psi(k).
\end{align*}
% \subsection{Legge di evoluzione temporale}
Uso il tempo per parametrizzare (non relativistica), quindi gli stati sono \( \ket{\psi(t)} \). Allora
\[
\braket{q|\psi(t_1)} = \psi(q; t_1) = e^{(t_1 - t_0) \frac{d}{dt}} \psi(q; t)|_{t = t_0}
\]
e vorrei descrivere questa evoluzione temporale con un operatore tale che
\[
\ket{\psi(t_1)} = e^{(t_1 - t_0) \frac{d}{dt}} \ket{\psi(t)}|_{t = t_0} = S(t_1, t_0) \ket{\psi(t)}|_{t = t_0}
\]
dove mi aspetto che \( S \equiv S(\hat{q}, \hat{p}; t_1, t_0) \).

\bigskip
\noindent
Inoltre devo pretendere \( S^\dag S = \mathbbm{1} \) perché
\[
\braket{\psi(t_1)|\psi(t_1)} = \braket{\psi(t_0)|S^\dag S|\psi(t_0)} = \braket{\psi(t_0)|\psi(t_0)}.
\]
Allora è comodo passare all'esponenziale
\[
S(t_0, t_1) = e^{i O_H(t_1, t_0)}
\]
e fermarmi allo sviluppo al primo ordine \( (t_1 - t_0) = t \)
\[
S(t_0, t_1) = \mathbbm{1} + i \varepsilon \underbrace{\left. \frac{d}{dt_1} O_H(t_1, t_0)\right|_{t_1 = t_0}}_{\hat{\mathcal{H}}(t_0)}
\]
e in generale
\[
S(t + \varepsilon, t) = \mathbbm{1} + i \varepsilon \hat{\mathcal{H}}(\hat{p}, \hat{q}; t).
\]
% \subsection{Invarianza per traslazioni temporali di $\mathcal{H}$}
So che $ S(t_1, t_0) = S(t_1 + \delta, t_0 + \delta) $, nel caso infinitesimo
\[
\mathbbm{1} + i \varepsilon \hat{\mathcal{H}}(t_0) = \mathbbm{1} + i \varepsilon \hat{\mathcal{H}}(t_0 + \delta) \quad \Rightarrow \quad \frac{d}{dt} \hat{\mathcal{H}}(t) = 0
\]
cio\`e $ \hat{\mathcal{H}} \equiv \hat{\mathcal{H}}(\hat{p}, \hat{q}) $ (N.B.: \`e un ``se e solo se'').

Se $ \hat{\mathcal{H}} $ non dipende dal tempo, allora $ S $ e $ \hat{\mathcal{H}} $ commutano. La condizione di partenza equivale a $ S(t_1 + \delta, t_0 + \delta) = S(t_1, t_0) $ con $ \delta $ infinitesimo questa volta,
\begin{align*}
S(t_1 + \delta, t_0 + \delta) &= S(t_1 + \delta, t_1) S(t_1, t_0) S(t_0, t_0 + \delta) = \\
&= S(t_1 + \delta, t_1) S(t_1, t_0) S^{-1}(t_0 + \delta, t_0) = \\
&=  S(t_1 + \delta, t_1) S(t_1, t_0) S^\dag(t_0 + \delta, t_0) = \\
&= \left( \mathbbm{1} + i \delta \hat{\mathcal{H}} \right) S(t_1, t_0) \left( \mathbbm{1} - i \delta \hat{\mathcal{H}} \right) = \\
&= S(t_1, t_0) + i \delta \left[ \hat{\mathcal{H}}, S(t_1, t_0) \right] + \mathcal{O}(\delta) \quad \Rightarrow
\end{align*}
\[
\Rightarrow \quad \left[ \hat{\mathcal{H}}, S(t_1, t_0) \right] = 0.
\]
Inoltre gli autovalori di $ \hat{\mathcal{H}} $ si conservano per traslazione temporale
\[
\hat{\mathcal{H}} \ket{E} = E \ket{E}.
\]
\begin{align*}
0 &= \braket{\varphi|[\hat{\mathcal{H}}, S(t, t_0)]|\psi} = \braket{\varphi|(\hat{\mathcal{H}} S(t, t_0) - S(t, t_0) \hat{\mathcal{H}})|\psi} = \\
&= \int_{-\infty}^{+\infty} \int_{-\infty}^{+\infty} \braket{\varphi|E} \braket{E|(\hat{\mathcal{H}} S - S \hat{\mathcal{H}})|E^\prime} \braket{E^\prime|\psi} \, dE \, dE^\prime = \\
&= \int_{-\infty}^{+\infty} \int_{-\infty}^{+\infty} \braket{\varphi|E} \left( E \braket{E|S|E^\prime} - E^\prime \braket{E|S|E^\prime}\right) \braket{E^\prime|\psi} \, dE \, dE^\prime  = \\
&= \int_{-\infty}^{+\infty} \int_{-\infty}^{+\infty} (E - E^\prime) \braket{\varphi|E} \braket{E|S|E^\prime} \braket{E^\prime|\psi} \, dE \, dE^\prime
\end{align*}
cio\`e $ \braket{E|S|E^\prime} = 0 $ se $ E \neq E^\prime $.

Digressione sulla lagrangiana.
\begin{align*}
\frac{d \mathcal{L}}{dt} \left( \underline{q}(t), \dot{\underline{q}}(t) \right) &= \sum_k \frac{\partial \mathcal{L}}{\partial q_k} \dot{q}_k + \sum_k \frac{\partial \mathcal{L}}{\partial \dot{q}_k} \ddot{q}_k = \left[ \frac{d}{dt} \frac{\partial \mathcal{L}}{\partial \dot{q}_k} = \frac{\partial \mathcal{L}}{\partial q_k} \right] = \\
&= \sum_k \dot{p}_k \dot{q}_k + \sum_k p_k \ddot{q}_k = \frac{d}{dt} \sum_k p_k \dot{q}_k
\end{align*}
\[
\frac{d}{dt} \mathcal{H} \left( \underline{q}(t), \underline{p}(t) \right) = \frac{d}{dt} \left( \sum_k p_k \dot{q}_k - \mathcal{L} \right) = 0.
\]
Fine digressione sulla lagrangiana.

Voglio che $\hat{\mathcal{H}} = -c \mathcal{H}$. Dimensionalmente, poich\'e
\[
S(t, t_0) \ket{\psi(t_0)} = \ket{\psi(t)} \; \xrightarrow[\text{infinitesim.}] \; \ket{\psi(t)} = e^{(t - t_0) \frac{d}{dt}} \ket{\psi(t)} |_{t_0} =
\]
\[
= \left( \mathbbm{1} + (t - t_0) \frac{d}{dt} + \mathcal{O}(t - t_0) \right) \ket{\psi(t_0)} \quad \Rightarrow
\]
\[
\Rightarrow \quad \hat{\mathcal{H}}(t) \ket{\psi(t)} = \frac{d}{dt} \ket{\psi(t)}
\]
la costante \`e un'azione, quindi $ c \equiv \frac{1}{\hbar} $. Riscrivo
\[
\hat{\mathcal{H}} = - \frac{1}{\hbar} \mathcal{H} \left( \hat{p}, \hat{q} \right) \quad \Rightarrow \quad \boxed{i \hbar \frac{d}{dt} \ket{\psi(t)} = \mathcal{H}( t; \hat{p}, \hat{q}) \ket{\psi(t)}} \; \text{eq. di Schr\"odinger.}
\]
Per $ \mathcal{H} = T + V = \frac{\hat{p}^2}{2m} + V(\hat{q}) $ per stati $ \ket{q} $
\[
\braket{q|i \hbar \frac{d}{dt}|\psi(t)} = \braket{q|\frac{\hat{p}^2}{2m} + V(\hat{q})|\psi(t)} \quad \Rightarrow
\]
\[
\left[
\begin{aligned}
\braket{q|\hat{p}^2|\psi} &= \int_{-\infty}^{+\infty} \braket{q|\hat{p}|k} \braket{k|\hat{p}|\psi} \, dk = \\
&= \int_{-\infty}^{+\infty} (\hbar k)^2 \braket{q|k} \braket{k|\psi} \, dk = \\
&= \int_{-\infty}^{+\infty} \frac{(\hbar k)^2}{\sqrt{2\pi}} e^{iqk} \psi(k) \, dk \\
\braket{q|\hat{q}|\psi} &= q \psi(q) \\
\braket{q|F(\hat{q})|\psi} &= F(q) \, \psi(q)
\end{aligned}
\right]
\]
\[
\Rightarrow \quad i \hbar \frac{d \psi}{dt}(q; t) = - \frac{\hbar^2}{2m} \frac{d^2 \psi}{dq^2}(q;t) + V(q) \psi(q;t).
\]

Per l'evoluzione temporale, $ \ket{\psi(t)} = S(t, t_0) \ket{\psi(t_0)} $, quindi
\[
i \hbar \frac{d}{dt} S(t, t_0) \ket{\psi(t_0)} = \mathcal{H}(t) S(t, t_0) \ket{\psi(t_0)}
\]
e, a seconda dei casi, ho diverse discussioni.

Caso 1.
\[
\frac{d\mathcal{H}}{dt} = 0 \qquad \text{(stato legato statico)}
\]
\[
i \hbar \frac{d}{dt} S(t, t_0) = \mathcal{H} S(t, t_0)
\]
\[
\boxed{S(t, t_0) = \exp \left\{ \frac{t - t_0}{i \hbar}\mathcal{H} \right\} }
\]

Caso 2.
\[
\frac{d\mathcal{H}}{dt} \neq 0 \text{ e } [\mathcal{H}(t), \mathcal{H}(t^\prime)] = 0
\]
\[
i \hbar \frac{d}{dt} S(t, t_0) = \mathcal{H}(t) S(t, t_0) \quad \Rightarrow \quad \boxed{S(t, t_0) = \exp{\left\{ \frac{1}{i \hbar} \int_{t_0}^{t} \mathcal{H}(t^\prime) \, dt^\prime \right\}}}
\]

Caso 3.
\[
\frac{d\mathcal{H}}{dt} \neq 0 \text{ e } [\mathcal{H}(t), \mathcal{H}(t^\prime)] \neq 0
\]
\begin{figure}
\centering
\begin{tikzpicture}
  \draw[thick] (0,5) -- (0,0) -- (5,0);
  \draw[dashed] (0,0) -- (5,5);
  \draw (0,3) -- (3,3) -- (3,0);

  \node (t0) at (-.2,-.2) {$t_0$};
  \node (t1) at (3,-.3) {$t'$};
  \node (t2) at (-.3,3) {$t''$};
\end{tikzpicture}
\caption{Logica di funzionamento della funzione $T[O_1(t_1)O_2(t_1)]$.}
\label{schema_T}
\end{figure}
Definisco
\[
T \left[ O_1(t_1) O_2(t_1) \right] \equiv \left\{
\begin{aligned} 
    & O_1(t_1) O_2(t_1) & t_1 > t \\ 
    & O_2(t_1) O_1(t_1) & t_1 < t
\end{aligned} \right. .
\]
Sviluppando al secondo ordine, facendo uso dello schema in Figura \ref{schema_T},
\begin{align*}
S(t, t_0) &= T \exp{\left\{ \frac{1}{i \hbar} \int_{t_0}^t \mathcal{H}(t^\prime) \, dt^\prime \right\}} = \\
&= \mathbbm{1} + \frac{1}{i \hbar} \int_{t_0}^t \mathcal{H}(t^\prime) \, dt^\prime + \frac{1}{2} \frac{1}{(i \hbar)^2} \int_{t_0}^t \int_{t_0}^t T \left( \mathcal{H}(t^\prime) \mathcal{H}(t'') \right) \, dt' \, dt'' = \\
&= \mathbbm{1} + \frac{1}{i \hbar} \int_{t_0}^t \mathcal{H}(t^\prime) \, dt^\prime + \frac{1}{2} \frac{1}{(i \hbar)^2} \left( \int_{t_0}^t \int_{t_0}^{t'} \mathcal{H}(t') \mathcal{H}(t'') \, dt' \, dt'' + \right. \\
& \left. \quad + \int_{t_0}^t \int_{t'}^t \mathcal{H}(t'') \mathcal{H}(t') \, dt'' \, dt' \right) = \\
&= \mathbbm{1} + \frac{1}{i \hbar} \int_{t_0}^t \mathcal{H}(t^\prime) \, dt^\prime + \frac{1}{(i \hbar)^2} \int_{t_0}^t \int_{t_0}^{t'} \mathcal{H}(t') \mathcal{H}(t'') \, dt' \, dt''.
\end{align*}
Quindi per l'$n$-esimo termine dello sviluppo in serie di Taylor dell'esponenziale applico lo stesso metodo, cio\`e
\[
T \left( \int_{t_0}^t \mathcal{H}(t_1) \, dt_1 \right) \left( \int_{t_0}^{t_1} \mathcal{H}(t_2) \, dt_2 \right) \cdots \left( \int_{t_0}^{t_{n - 1}} \mathcal{H}(t_n) \, dt_n \right) =
\]
\[
= n! \int_{t_0}^t \int_{t_0}^{t_1} \int_{t_0}^{t_2} \cdots \int_{t_0}^{t_{n - 1}} \mathcal{H}(t_1) \mathcal{H}(t_2) \cdots \mathcal{H}(t_n) \, dt_1 \, dt_2 \cdots dt_n
\]
e gli $ n! $ si semplificano. Adesso controllo se $ S(t, t_0) $ \`e soluzione dell'equazione di Schr\"odinger.
\begin{align*}
i \hbar \frac{d}{dt} S(t, t_0) &= \mathcal{H}(t) + \frac{1}{i \hbar} \mathcal{H}(t) \int_{t_0}^t \mathcal{H}(t_2) \, dt_2 + \cdots + \\
&\quad + \frac{1}{(i \hbar)^{n + 1}} \mathcal{H}(t) \int_{t_0}^t \int_{t_0}^{t_2} \cdots \int_{t_0}^{t_{n - 1}} \mathcal{H}(t_2) \cdots \mathcal{H}(t_n) \, dt_2 \cdots dt_n = \\
&= \mathcal{H}(t) S(t, t_0)
\end{align*}

Nel caso pi\`u semplice in cui $ \frac{d}{dt} \mathcal{H} = 0 $ so che
\[
S(t, t_0) = \exp\left\{ \frac{1}{i \hbar} (t - t_0) \mathcal{H} \right\}
\]
allora l'operatore di evoluzione temporale lo posso calcolare solo quando $ \mathcal{H} \ket{n} = E_n \ket{n} $ (diagonalizzabilit\`a dell'operatore), da cui poi $ \braket{n|\psi} = c_n $ e $ p_n = |c_n|^2$.

Se ho un potenziale per cui so calcolare gli stati di energia, posso calcolare gli elementi di matrice, ottenendo
\[
S(t, t_0) = \sum_{n,m} \ket{m} \braket{m|S(t, t_0)|n} \bra{n} =
\]
\[
\left[ \exp{\left\{ \frac{1}{i \hbar} (t - t_0)  \mathcal{H} \right\} } \ket{n} = \exp{\left\{ \frac{1}{i \hbar} (t - t_0) E_n \right\}} \ket{n} \right]
\]
\[
= \sum_{n,m} \ket{m} \braket{m|\exp{\left\{ \frac{1}{i \hbar} (t - t_0) E_n \right\}}|n} \bra{n} =
\]
\[
= \sum_{n,m} \exp{\left\{ \frac{1}{i \hbar} (t - t_0) E_n \right\}} \ket{m} \braket{m|n} \bra{n} = 
\]
\[
= \sum_n \exp{\left\{ \frac{1}{i \hbar} (t - t_0) E_n \right\}} \ket{n} \bra{n}
\]
operatore diagonale nella base degli $ \ket{n} $.

Quindi per $ \ket{\psi(t_0)} = \ket{k} $ al tempo $ t $
\[
\ket{\psi(t)} = S(t, t_0) \ket{k} = \sum_n \exp{\left\{ \frac{1}{i \hbar} (t - t_0) E_n \right\}} \ket{n} \braket{n|k} = e^{\frac{1}{i \hbar}(t-t_0)E_k}\ket{k}
\]
cio\`e $ \ket{\psi(t_0)} $ e $ \ket{\psi(t)} $ \textbf{differiscono per una fase}. Inoltre
\[
\ket{\psi(t_0)} = \sum_k \ket{k} \overbrace{\braket{k|\psi(t_0)}}^{c_k(t_0)} = \sum_k c_k(t_0) \ket{k} \quad \Rightarrow
\]
\[
\Rightarrow \quad \ket{\psi(t)} = \sum_n \exp{\left\{ \frac{1}{i \hbar} (t - t_0) E_n \right\}} \ket{n} \overbrace{\braket{n|\psi(t)}}^{c_n(t)} \; \Rightarrow
\]
\[
\Rightarrow \; c_n(t) = c_n(t_0) \exp{\left\{ \frac{1}{i \hbar} (t - t_0) E_n \right\}}
\]
cio\`e la probabilit\`a si conserva perch\'e 
\[
p_n(t) = |c_n(t)|^2 = |c_n(t_0)|^2 = p_n(t_0).
\]

Per lo stato $ \ket{\psi} $ il valor medio di un operatore sulla base degli $ \ket{n} $ autostati dell'energia 
\[
\braket{n|A|n} = \braket{n(t_0)|e^{- \frac{1}{i \hbar} (t_n - t_0) E_n} A e^{\frac{1}{i \hbar} (t_n - t_0) E_n}|n(t_0)} = \braket{n(t_0)|A|n(t_0)}.
\]
Gli stati $ \ket{n} $ sono detti \textbf{stati stazionari}.

\subsubsection{Esempio: hamiltoniana e matrici di Pauli}
$$ \mathcal{H} = B \sigma_z, \qquad \sigma_z = \begin{pmatrix}
1 & 0 \\
0 & -1
\end{pmatrix}.
\[
\]
S = \exp \left\{ t \frac{B \sigma_z}{i\hbar} \right\}
\[
\]
\left| \braket{\psi_1|\psi_2} \right|^2 = \left| \braket{\psi_1|\psi_2(t)} \right|^2 = \left| \braket{\psi_1|S|\psi_2(0)} \right|^2
\[
\]
\ket{\psi(t)} = \exp \left\{ \frac{Bt\sigma_z}{i\hbar} \right\} \frac{\ket{+} + \ket{-}}{\sqrt{2}} = \frac{1}{\sqrt{2}} \left( e^{\frac{Bt}{i\hbar}} \ket{+} + e^{-\frac{Bt}{i\hbar}} \ket{-} \right)
\[
\]
\bra{\psi_2} = \frac{1}{\sqrt{2}} \left( \bra{+} + \bra{-} \right)
\[
\]
\left| \braket{\psi_2|\psi_1} \right|^2 = \frac{1}{4} \left| e^{\frac{Bt}{i\hbar}} + e^{-\frac{Bt}{i\hbar}} \right|^2 = \cos^2 \frac{Bt}{\hbar}
\[

\subsubsection{Esempio: ancora su hamiltoniana e matrici di Pauli}
\]
\mathcal{H} = t B \sigma_z.
\[
\]
S = \exp \left\{ \frac{1}{i\hbar} \int_{0}^t \mathcal{H}(t') \, dt' \right\} = \exp \left\{ \frac{t^2 B \sigma_z}{2i\hbar} \right\}.
\[
\]
\ket{\psi(t)} = \exp \left\{ \frac{Bt^2\sigma_z}{i\hbar} \right\} \frac{\ket{+} + \ket{-}}{\sqrt{2}} = \frac{1}{\sqrt{2}} \left( e^{\frac{Bt^2}{i\hbar}} \ket{+} + e^{-\frac{Bt^2}{i\hbar}} \ket{-} \right)
\[
\]
\left| \braket{\psi_2|\psi_1} \right|^2 = \frac{1}{4} \left| e^{\frac{Bt^2}{i\hbar}} + e^{-\frac{Bt^2}{i\hbar}} \right|^2 = \cos^2 \frac{Bt^2}{\hbar}
$$
% \subsection{Rappresentazione di Heisenberg}
\[
A^H(t) = S^{-1}(t, t_0) A^S(t_0) S(t, t_0), \quad A^S(t_0) \text{ costante nel tempo.}
\]
Per \( \ket{\psi} \) ho \( \ket{\psi(t)} = S(t, t_0) \ket{\psi(t_0)} \) e ho elemento di matrice
\begin{align*}
\braket{\psi_H(t_0)|A^H(t)|\psi_H(t_0)} &= \braket{\psi_S(t_0)|S^{-1}(t, t_0) A^S(t_0) S(t, t_0)|\psi_S(t_0)} = \\
&= \braket{\psi_S(t)|A^S(t_0)|\psi_S(t)}
\end{align*}
Quindi passo da \( S \) ad \( H \) e nel primo l'operatore \`e costante nel tempo ed evolvono gli stati, nel secondo \`e il contrario.

Allora posso prendere gli \( \ket{n(t)} \) (fissi nel tempo) e uso
\[
A^H(t) \ket{n(t)} = \lambda_n \ket{n(t)}
\]
\[
S^{-1}(t, t_0) A^S(t_0) S(t, t_0) \ket{n(t)} = \lambda_n \ket{n(t)}
\]
\[
A^S(t_0) \underbrace{S(t, t_0) \ket{n(t)}}_{\ket{n(t_0)}} = \lambda_n \underbrace{S(t, t_0) \ket{n(t)}}_{\ket{n(t_0)}}
\]
\[
\ket{n_H(t)} = S^{-1}(t, t_0) \ket{n_S(t_0)}
\]
Quindi una generica misura
\[
\braket{n_S(t_0)|\psi_S(t)} = \braket{n_S(t_0)|S(t, t_0)|\psi_S(t_0)} = \braket{n_H(t)|\psi_H(t)}.
\]

Per studiare l'evoluzione temporale di \( A^H(t) \) faccio
\begin{align*}
\frac{d}{dt} A^H(t) &= \frac{d}{dt} \left( S^{-1}(t, t_0) A^S(t_0) S(t, t_0) \right) = \\
&= \frac{\partial S^{-1}}{\partial t} A^S(t_0) S + S^{-1} A^S \frac{\partial S}{\partial t} + S^{-1} \frac{\partial A^S}{\partial t} S,
\end{align*}
ma so che
\[
i \hbar \frac{\partial S}{\partial t} = \mathcal{H} S \; \Rightarrow \; -i \hbar \frac{\partial S^\dag}{\partial t} = S^\dag \mathcal{H}^\dag \; \Rightarrow \; -i \hbar \frac{\partial S^{-1}}{\partial t} = S^{-1} \mathcal{H}
\]
quindi 
\[
\frac{d}{dt} A^H(t) = -\frac{1}{i \hbar} S^{-1} \mathcal{H} A^S S + S^{-1} A^S \frac{\mathcal{H} S}{i \hbar} + \frac{\partial}{\partial t} A^H(t).
\]
Poich\'e per definizione \( S(t, t_0) A^H(t) = A^S(t_0) S(t, t_0) \), ho
\begin{align*}
\frac{d}{dt} A^H(t) &= -\frac{1}{i \hbar} S^{-1} \mathcal{H} S A^H + \frac{1}{i \hbar} A^H S^{-1} \mathcal{H} S + \frac{\partial}{\partial t} A^H(t) = \\
&= \frac{1}{i \hbar} \left[ A^H, \underbrace{S^{-1} \mathcal{H} S}_{\mathcal{H}^H} \right] + \frac{\partial}{\partial t} A^H(t)
\end{align*}
\begin{align*}
\frac{d}{dt} A_{ij}(t) &= \frac{d}{dt} \braket{\psi_i^S(t)|A^S(t_0)|\psi_j^S(t)} = \\
&= \frac{d}{dt} \braket{\psi_i^H(t_0)|A^H(t)|\psi_j^H(t_0)} = \\
&= \frac{1}{i \hbar} \braket{\psi_i^H(t_0)|\left\{ \left[ A^H(t), \mathcal{H}^H \right] + \frac{\partial}{\partial t} A^H(t) \right\}|\psi_j^H(t_0)}
\end{align*}
per\`o \( A^H = S^{-1} A^S S \) e \( A^H \mathcal{H}^H = S A^S S S^{-1} \mathcal{H} S = S A^S \mathcal{H} S \)
\[
\frac{d}{dt} A_{ij}(t) = \braket{\psi_i^S(t)|\frac{1}{i \hbar} \left[ A^S(t_0), \mathcal{H}^S \right]|\psi_j^S(t)} + \braket{\psi_i^S(t)|\frac{\partial A^S}{\partial t}|\psi_j^S(t)}.
\]

Esempio: \( A \) e \( B \) costanti del moto \( \Rightarrow [A, B] \) costante del moto. Infatti
\[
[A, \mathcal{H}] = [B, \mathcal{H}] = 0 \quad \Rightarrow \quad \left[ [A, B], \hat{\mathcal{H}} \right] = 0.
\]
% \subsection{Applicazione a $\hat{p}$ e $\hat{q}$}
Non avendo dipendenza dal tempo,
\[
\frac{d \hat{q}^H}{dt} = \frac{1}{i \hbar} \left[ q^H, \mathcal{H}^H \right], \quad \frac{d \hat{p}^H}{dt} = \frac{1}{i \hbar} \left[ p^H, \mathcal{H}^H \right]
\]
e specifico $ \mathcal{H} = \frac{\hat{p}^2}{2m} + V(\hat{q}) $. Poi so che $ [A, BC] = B[A, C] + [A, B] C $ e inoltre
\begin{align*}
\left[ \hat{q}, f(\hat{p}) \right] &= \left[ \hat{q}, \sum_k a_k \hat{p}^k \right] = \sum_k f_k \left[ \hat{q}, \hat{p}^k \right] = \\
&= \sum_k f_k \left( \hat{p} \left[ \hat{q} \hat{p}^{k - 1} \right] + \left[ \hat{q}, \hat{p} \right] \hat{p}^{k - 1} \right) = \\
&= i \hbar \sum_k f_k k \hat{p}^{k - 1} = i \hbar \frac{\partial f}{\partial \hat{p}} (\hat{p}) \\
\left[ \hat{p}, g(\hat{q}) \right] &= - i \hbar \frac{\partial g}{\partial \hat{q}}(\hat{q})
\end{align*}
quindi
\[
\frac{d \hat{q}^H}{dt} = \frac{\hat{p}}{m}, \qquad \frac{d \hat{p}^H}{dt} = -\frac{\partial V}{\partial q}.
\]

Nel caso in cui $ V(\hat{p}, \hat{q}) = \hat{p} f(\hat{q}) + f(\hat{q})\hat{p} $
\begin{align*}
\frac{d\hat{q}^H}{dt} &= \frac{1}{i \hbar} \left[ \hat{q}, \frac{\hat{p}^2}{2m} + \hat{p} f(\hat{q}) + f(\hat{q})\hat{p} \right] = \\
&= \frac{\hat{p}}{m} + \frac{1}{i \hbar} \left[ \hat{q}, \hat{p} f(\hat{q}) + f(\hat{q})\hat{p} \right] = \\
&= \frac{\hat{p}}{m} + \frac{1}{i \hbar} \left( \; \left[ \hat{q} f(\hat{q}) \right] + \text{MANCANTE} \right)
\end{align*}

In particolare, se $ A = A(\hat{p}, \hat{q}) $,
\[
\frac{d}{dt} \hat{A}_H(\hat{p}, \hat{q}) = \frac{1}{i \hbar} \left[ \hat{A}, \mathcal{H} \right] + \frac{\partial \hat{A}}{\partial t}
\]
molto simile all'approccio classico dove
\[
\frac{d}{dt} A(p, q) = \left\{ A, \mathcal{H} \right\} + \frac{\partial A}{\partial t}.
\]

Anche in Meccanica Quantistica ho il \textbf{commutatore fondamentale}
\[
\left[ \hat{p}, \hat{q} \right] = -i \hbar.
\]
Questo si ripercuote sui valori medi
\[
\frac{d}{dt} \braket{\hat{p}} = \braket{-\frac{\partial \mathcal{H}}{\partial \hat{q}}}, \qquad \frac{d}{dt} \braket{\hat{q}} = \braket{\frac{\partial \mathcal{H}}{\partial \hat{p}}}
\]
che \`e vero sia per $ H $ che per $ S $.

Esempio: $ \hat{\mathcal{H}} = \frac{\hat{p}^2}{2m} + a\hat{x} $
\[
\frac{d}{dt} \hat{x} = \frac{1}{i\hbar} [\hat{x}, \mathcal{H}] = \frac{\hat{p}}{m}, \quad \frac{d}{dt} \hat{p} = \frac{1}{i\hbar} [\hat{p}, \mathcal{H}] = -a
\]
\[
\frac{d}{dt} \hat{x}^2 = \frac{1}{i\hbar} [\hat{x}^2, \mathcal{H}] = \frac{1}{i\hbar} \left( \hat{x} [\hat{x}, \mathcal{H}] + [\hat{x}, \mathcal{H}] \hat{x} \right) = \frac{1}{m} (\hat{x}\hat{p} + \hat{p}\hat{x})
\]
\[
\frac{d}{dt} \hat{p}^2 = \frac{1}{i\hbar} \left( \hat{p}[\hat{p}, \mathcal{H}] + [\hat{p}, \mathcal{H}] \mathcal{H} \right) = -2a\hat{p}
\]

Esempio: $\mathcal{H} = \frac{\hat{p}^2}{2m} + a\hat{x}$. $ \Delta^2 \hat{x}(t) $ collegandolo all'indetrminazione.

Metodo 1:
\[
\frac{d}{dt} \Delta^2 \hat{x}(t) = \frac{1}{i\hbar} [\Delta^2 \hat{x}, \mathcal{H}] \text{ ecc.}
\]

Metodo 2: $ \Delta^2 \hat{x}(t) = \braket{\hat{x}^2} - (\braket{x})^2 $. So che la soluzione dell'equazione del moto 
\[
\hat{x}(t) = \underbrace{\hat{x}(t_0)}_{\hat{x}^S} - \frac{1}{2}\frac{a}{m}t^2 + \frac{\overbrace{\hat{p}_0(t)}^{\hat{p}^S}}{m}t
\]
\[
[\hat{p}^S, \hat{x}^S] = -i\hbar
\]
allora esplicitamente
\[
\braket{\left( \hat{x}^S +\frac{\hat{p}^S}{m}t - \frac{1}{2} \frac{a}{m} t^2 \right) (\dots)} - \left( \braket{(\dots)} \right)^2.
\]
% \subsection{Particella libera}
\[
\braket{q|\psi} = \psi(q) \quad \text{e} \quad \mathcal{H} = \frac{\hat{p}^2}{2m}.
\]
\[
\hat{p}\ket{k} = \hbar k \ket{k} \quad \Rightarrow \quad \mathcal{H} \ket{k} = \frac{\hbar^2 k^2}{2m} \ket{k}.
\]
Adesso passo agli autostati $ \ket{q} $ con $ \braket{q|\mathcal{H}|k} = E_k \braket{q|k} $, dove
\begin{align*}
\braket{q|k} &= \frac{1}{\sqrt{2\pi}} e^{ikq} \\
\braket{k|k'} &= \delta(k - k') \\
E_k &= \frac{\hbar^2 k^2}{2m}
\end{align*}
In rappresentazione di Schr\"odinger
\begin{align*}
\braket{q|k; t} &= \braket{q|e^{\frac{1}{i \hbar}t \mathcal{H}}|k; t_0 = 0} = \\
&= e^{\frac{1}{i \hbar}t E_k} \braket{q|k; t_0 = 0} = \\
&= \frac{1}{\sqrt{2\pi}} \exp \left\{ i \left( kq - \frac{E_k t}{\hbar} \right) \right\}, \qquad \omega_k = \frac{\hbar k^2}{2m} = \frac{E_k}{\hbar}.
\end{align*}
In rappresentazione di Heisenberg
\[
\braket{q|k} = \frac{1}{\sqrt{2\pi}} e^{ikq} \; \text{indipendente da } t,
\]
\[
\hat{p}_H(t) = \hat{p}_S,
\]
\[
\hat{q}_H(t) = \hat{q}_H(t_0) + (t - t_0) \frac{\hat{p}_H}{m} = \hat{q}_S(t_0) + (t - t_0) \frac{\hat{p}_S}{m}.
\]

In concreto, se risolvo alla Schr\"odinger l'elemento di matrice nella BON degli autostati di $ \mathcal{H} $
\begin{align*}
\braket{k'; t|\hat{p}|k; t} &= \hbar k \braket{k'; t|k; t} = \\
&= \hbar k \int_{-\infty}^{+\infty} \braket{k'; t|q} \braket{q|k; t} \, dq = \\
&= \frac{\hbar k}{2\pi} \int_{-\infty}^{+\infty} \exp \left\{ i \left[ (k - k')q - (\omega_k - \omega_{k'})t \right] \right\} \, dq = \\
&= \frac{\hbar k}{2\pi} e^{-i(\omega_k - \omega_{k'})t} \int_{-\infty}^{+\infty} e^{i(k - k')q} \, dq = \\
&= \hbar k \delta(k - k') e^{-i(\omega_k - \omega_{k'})t} = \\
&= \hbar k \delta(k - k') = \\
&= \hbar k' \delta(k - k')
\end{align*}
che era ovvio nella rappresentazione di Heisenberg.
\begin{align*}
\braket{k'; t|\hat{q}|k; t} &= \int_{-\infty}^{+\infty} \braket{k'; t|\hat{q}|q} \braket{q|k; t} \, dq = \\
&= \int_{-\infty}^{+\infty} q \braket{k'; t|q} \braket{q|k; t} \, dq = \\
&= \frac{1}{2\pi} \int_{-\infty}^{+\infty} q \exp \left\{ i \left[ (k - k')q - (\omega_k - \omega_{k'})t \right] \right\} \, dq = \\
&= \frac{1}{2\pi} \int_{-\infty}^{+\infty} \frac{1}{i} \frac{d}{dk} \left( \exp \left\{ i(k - k')q \right\} \right) \exp \left\{ -i(\omega_k - \omega_{k'})t \right\} \, dq = \\
&=  \frac{1}{2\pi i} \int_{-\infty}^{+\infty} \left[ \frac{d}{dk} \left( \exp \left\{ i(k - k')q \right\} \exp \left\{ -i(\omega_k - \omega_{k'})t \right\} \right)  + \right. \\
&\quad + \left. \exp \left\{ i(k - k')q \right\} \frac{i\hbar k t}{m} \exp \left\{ -i(\omega_k - \omega_{k'})t \right\} \right] \, dq = \\
&= \int_{-\infty}^{+\infty} \frac{1}{2\pi} \left( -i \frac{d}{dk} + \frac{\hbar k t}{m} \right) \exp \left\{ i \left[ (k - k')q - (\omega_k - \omega_{k'})t \right] \right\} \, dq = \\
&= -i \frac{d}{dk} \delta(k - k') + \frac{\hbar k t}{m} \delta(k - k') = \\
&= \braket{k|\hat{q}_S|k'} + \frac{t}{m} \braket{k|\hat{p}_S|k'}
\end{align*}

La cosa interessante \`e che
\[
\left[ \hat{q}_H(t), \hat{q}_H(t') \right] = \frac{t - t'}{m} \left[ \hat{q}_S, \hat{p}_S \right] = i \frac{\hbar(t - t')}{m}
\]
e per il principio di indeterminazione
\[
\Delta^2 q(t) \, \Delta^2 q(t') \geq \frac{\hbar^2 (t - t')^2}{4m^2}.
\]
% \subsection{Particella libera (parte 2)}
\[
\braket{x|\psi(t)} = \psi(x, t), \quad \hat{\mathcal{H}} = \frac{\hat{p}}{2m} \rightarrow \begin{cases} 
\hat{p} \ket{k} = \hbar k \ket{k} \\
\hat{\mathcal{H}} \ket{k} = \hbar \omega_k \ket{k}
\end{cases}, \quad \hbar \omega_k = \frac{\hbar^2 k^2}{2m}
\]
\[
S = e^{(t - t_0)\frac{\mathcal{H}}{i \hbar}} \quad \braket{x|k} = \frac{1}{\sqrt{2\pi}} e^{ikx} \quad \braket{x|k; t} = \frac{1}{\sqrt{2\pi}} e^{i(kx - \omega_k t)}
\]
\begin{align*}
\psi(x, t) &= \int_{-\infty}^{+\infty} \braket{x|k} \braket{k|\psi(t)} \, dk \\
&= \int_{-\infty}^{+\infty} \braket{x|k} \braket{k|S(t, 0)|\psi(0)} \, dk \\
&= \int_{-\infty}^{+\infty} \frac{1}{\sqrt{2\pi}} e^{i(kx - \omega_k t)} \psi(k; 0) \, dk
\end{align*}
\( \psi(x, t) \) pacchetto d'onda. Pacchetto d'onda ad indeterminazione minima \( \Delta^2 \hat{x} \Delta^2 \hat{p} = \frac{\hbar^2}{4} \), in \( t = 0 \).
\[
\psi(x, 0) = \frac{1}{\sqrt{2\pi}} \int_{-\infty}^{+\infty} e^{ikx} \psi(k; 0) \, dk
\]
Allora anche la disuguaglianza di Schwarz deve diventare un'uguaglianza, cioè
\[
\Delta \hat{x} = \hat{x} - \braket{\hat{x}}, \qquad \Delta \hat{p} = \hat{p} - \braket{\hat{p}}
\]
cioè per \( \ket{\psi} \)
\[
\left( \hat{p} - \braket{\hat{p}} \right) \ket{\psi} = i \lambda \left( \hat{x} - \braket{\hat{x}} \right) \ket{\psi}.
\]
Anche l'altra deve essere uguale, quindi \( \lambda \in \mathbb{R} \).
\[
\braket{x|(\hat{p} - \underbrace{\braket{\hat{p}}}_{p_0})|\psi} = i \lambda \braket{x|(\hat{x} - \underbrace{\braket{\hat{x}}}_{x_0})| \psi}
\]
\[
-i \hbar \frac{d}{dx} \psi(x) - p_0 \psi(x) = i \lambda \left( x\psi(x) - x_0 \psi(x) \right)
\]
\[
\frac{d}{dx} \psi(x) = i \frac{p_0}{\hbar} \psi(x) - \frac{\lambda}{\hbar} (x - x_0) \psi(x)
\]
\[
\psi(x) = N \exp \left\{ i\frac{p_0}{\hbar} x - \frac{\lambda}{2\hbar} (x - x_0)^2 \right\}
\]
che è una gaussiana modulata (Figura \ref{gaussiana_modulata}).
\begin{figure}
\centering
\begin{tikzpicture}
  \draw[thick] plot[smooth] file {./grafici/wave_packet.txt};
  \draw[->] (-2.5,0) -- (6.5,0);
  \draw[->] (0,-4) -- (0,4.5);
\end{tikzpicture}
\caption{Andamento qualitativo della funzione d'onda \(\psi\) per la particella libera. è una gaussiana modulata, pertanto la distribuzione è quella tipica del pacchetto d'onda.}
\label{gaussiana_modulata}
\end{figure}
\begin{align*}
\braket{\hat{x}} &= \braket{\psi|\hat{x}|\psi} = \int_{-\infty}^{+\infty} x \braket{\psi|x} \braket{x|\psi} \, dx = \\
&= \int_{-\infty}^{+\infty} x \left| \psi(x) \right|^2 \, dx = \\
&= |N|^2 \int_{-\infty}^{+\infty} e^{-\frac{\lambda}{\hbar}(x - x_0)^2} (x - x_0 + x_0) \, dx \\
&= |N|^2 \int_{-\infty}^{+\infty} e^{-\frac{\lambda}{\hbar} {x'}^2} (x' + x_0) \, dx' = \\
&= x_0 \braket{\psi|\psi} = x_0
\end{align*}
\begin{align*}
\braket{\hat{p}} &= \int_{-\infty}^{+\infty} \psi^*(x) (-i\hbar) \frac{d}{dx} \psi(x) \, dx = \\
&= \int_{-\infty}^{+\infty} \left[ \psi^*(x) \left( -i\hbar \frac{i p_0}{\hbar} \right) \psi(x) + \psi^*(x) (-i\hbar) \left(\frac{-\lambda}{2\hbar}\right)(x - x_0) 2 \psi(x) \right] \, dx = \\
&= p_0 \braket{\psi|\psi} = p_0
\end{align*}
\[
\braket{\psi|\left( \hat{p} - p_0 \right)^2|\psi} = \int_{-\infty}^{+\infty} \psi^*(x) \left( -i\hbar \frac{d}{dx} - p_0 \right) \left( -i\hbar \frac{d}{dx} - p_0 \right) \psi(x) \, dx =
\]
MANCANTE

\bigskip
\noindent
Per un pacchetto gaussiano calcolo la dipendenza di indeterminazione dal tempo. In rappresentazione di Schr\"odinger MANCANTE.

\bigskip
\noindent
In rappresentazione di Heisenberg in \( t = 0 \) (\( \psi_H(t_0) = \psi_S(0) \))
\[
\braket{x|\psi} = \psi(x) = N \exp \left\{ i k_0 x \right\} \exp \left\{ -\frac{\lambda}{2} \frac{(x - x_0)^2}{\hbar} \right\}
\]
\[
\Delta^2 x = \Delta_0^2 x = \frac{\hbar}{2\lambda}, \; \Delta^2 p = \frac{\hbar^2}{4 \Delta_0^2 x} = \Delta_0^2 p \; \implies \; \braket{x} = x_0, \; \braket{p} = \hbar k_0
\]
Al \( t \) generico
\[
\braket{p} = \braket{\psi|p_H|\psi} = \braket{\psi|p_S|\psi} = \hbar k_0
\]
\[
\braket{x} = \braket{\psi|x_S|\psi} + \frac{t}{m} \braket{\psi|p_S|\psi} = x_0 + v_G t, \quad v_G = \frac{\hbar k_0}{m} \text{ velocità di gruppo}
\]
\[
\frac{d}{dt} p_H^2 = \frac{1}{i\hbar} \left[ p_H^2, \mathcal{H} \right] = 0
\]
\[
\frac{d}{dt} \Delta^2 p_H = 0 \quad \implies \quad \Delta_0^2 p = \Delta^2 p = \frac{\hbar^2}{4 \Delta_0^2 x}
\]
\[
\braket{\Delta^2 x} = \braket{x^2} - (\braket{x})^2 = \braket{\left( x_H(t) \right)^2} - \left( \braket{x}(t) \right)^2 = \braket{\left( x_H(t) - \braket{x}(t) \right)^2}
\]
\[
\frac{d}{dt} \braket{\Delta^2 x} = \braket{\frac{d}{dt} x_H^2(t)} - 2 \braket{x} \frac{d}{dt} \braket{x}(t)
\]
\begin{align*}
\braket{\frac{d}{dt} x_H^2(t)} &= \frac{1}{\hbar} [x_H^2, \mathcal{H}] = \\
&= \frac{1}{i\hbar} \left( [x_H, \mathcal{H}] x_H + x_H [x_H, \mathcal{H}] \right) = \\
&= \frac{1}{m} (p_H x_H + x_H p_H) = \\
&= \frac{1}{m} \left[ p_S \left( x_S + \frac{p_S}{m} t \right) + \left( x_S + \frac{p_S}{m} t \right) p_S \right] = \\
&= \frac{1}{m} \left( x_S p_S + p_S x_S \right) + \frac{2}{m} p_S^2 t
\end{align*}
\[
\braket{x} \frac{d}{dt} \braket{x}(t) = \left( \braket{x_S} + \frac{t}{m} \braket{p_S} \right) \frac{1}{m} \braket{p_S}
\]
\[
\frac{d}{dt} \braket{\Delta^2 x} = \frac{1}{m} \braket{x_S p_S + p_S x_S} - \frac{2}{m} \braket{x_S} \braket{p_S} + \frac{2t}{m^2} \underbrace{\left( \braket{p_S^2 - (\braket{p_S})^2} \right)}_{\Delta_0^2 p}
\]
Nel caso gaussiano
\[
\frac{d}{dt} \braket{\psi|\Delta^2 x|\psi} = \frac{1}{m} \braket{\psi|(x_S p_S + p_S x_S)|\psi} - \text{ecc.}
\]
\begin{align*}
\braket{x_S p_S + p_S x_S} &= \underbrace{2 \text{Re} \braket{x_S p_S}}_{2 x_0 p_0} \\
(p_S - p_0) \ket{\psi} &= i \lambda (x - x_0) \ket{\psi} \\
p_S \ket{\psi} &= p_0 \ket{\psi} + \lambda (x - x_0) \ket{\psi} 
\end{align*}
\[
\frac{d}{dt} \braket{\psi|\Delta^2 x|\psi} = \frac{2t}{m} \Delta_0^2 p 
\]
Impongo la condizione iniziale \( \Delta_0^2 x \)
\[
\Delta^2 x = \Delta_0^2 x + \frac{t^2}{m^2} \Delta_0^2 p
\]
Poiché
\[
\begin{cases}
\Delta^2 x(t) \Delta^2 x(0) \geq \frac{\hbar^2 t^2}{4m^2} \\
\Delta_0^2 x = \frac{\hbar^2}{4 \Delta_0^2 p} \text{ (minima indeterminazione)}
\end{cases} \implies \quad \Delta^2 x(t) \geq \frac{t^2}{m^2} \Delta_0^2 p
\]
Esempio: particella libera con 
\[
\psi(x) = N e^{-\frac{x^2}{\beta - i\gamma}}, \qquad \omega = \frac{\hbar k^2}{2m}.
\]
\begin{align*}
\psi(x, t) &= N' \int_{-\infty}^{+\infty} e^{-\frac{k^2}{4}(\beta - i\gamma)} e^{ikx} e^{-i\frac{\hbar k^2}{2m}t} \, dk = \\
&= \text{ MANCANTE} \\
\tilde{\psi}(k, t) &= \text{ MANCANTE}.
\end{align*}
Esempio:
\[
\psi(x, t) = \int_{-\infty}^{+\infty} e^{i(kx - \omega(k) t)} \tilde{\psi}(k) \, dk
\]
velocità di gruppo
\[
v_G = \frac{d}{dt} \braket{\hat{x}}(t)
\]
\[
\braket{\hat{x}} = \int_{-\infty}^{+\infty} \psi^*(k) i \frac{d}{dk} \psi(k) \, dk.
\]
% \subsection{Operatore parit\`a}
Si indica con $ \hat{\mathcal{P}} $ e per ogni stato nella base delle posizioni, \`e tale che
\[
\braket{x|\hat{\mathcal{P}}|\psi} = \braket{-x|\psi}.
\]
Questo operatore \`e di grande importanza nella Meccanica Quantistica perch\'e \`e allo stesso tempo autoaggiunto e unitario e agisce in modo particolare sugli operatori posizione, impulso e hamiltoniano.

Vediamo la sua autoaggiuntezza:
\[
\braket{x|\hat{\mathcal{P}}^\dag|\psi} = \int_{-\infty}^{+\infty} \braket{x|-x'} \psi(x') \, dx' = \braket{-x|\psi}
\]
quindi $ \hat{\mathcal{P}} = \hat{\mathcal{P}}^\dag $. L'unitariet\`a discende da 
\[
\braket{x|\hat{\mathcal{P}}\hat{\mathcal{P}}^\dag|\psi} = \int_{-\infty}^{+\infty} \braket{-x|-x'} \psi(x') \, dx' = \braket{x|\psi}
\]
da cui $ \hat{\mathcal{P}}\hat{\mathcal{P}}^\dag = \hat{\mathbb{I}} $. Un'altra propriet\`a interessante \`e l'idempotenza:
\[
\braket{x|\hat{\mathcal{P}}^2|\psi} = \int_{-\infty}^{+\infty} \braket{-x|x'} \psi(-x') \, dx' = \braket{x|\psi}
\]
quindi anche $ \hat{\mathcal{P}}^2 = \hat{\mathbb{I}} $, cos\`{i} si ha
\[
\hat{\mathcal{P}} = \hat{\mathcal{P}}^{-1} = \hat{\mathcal{P}}^\dag.
\]
Dall'idempotenza discende immediatamente che gli autovalori di $ \hat{\mathcal{P}} $ sono $ \pm 1 $ e per trovare le autofunzioni
\[
\braket{x|\hat{\mathcal{P}}|\psi} = \braket{x|\psi} \quad \Rightarrow \quad \psi \text{ funzione pari},
\]
\[
\braket{x|\hat{\mathcal{P}}|\psi} = -\braket{x|\psi} \quad \Rightarrow \quad \psi \text{ funzione dispari}.
\]

Per quanto riguarda gli operatori posizione e impulso, studiamo l'azione degli operatori composti (vogliamo arrivare a dire qualcosa sui commutatori) $ \hat{\mathcal{P}} \hat{x} \hat{\mathcal{P}}^{-1} $ e $ \hat{\mathcal{P}} \hat{p} \hat{\mathcal{P}}^{-1} $ attraverso la decomposizione dello stato sulla base delle coordinate.
\[
\braket{x|\hat{\mathcal{P}} \hat{x} \hat{\mathcal{P}}^{-1}|\psi} = -x \int_{-\infty}^{+\infty} \braket{-x|x'} \psi(x') = -x \braket{x|\psi} = - \braket{x|\hat{x}|\psi}
\]
\begin{align*}
\braket{x|\hat{\mathcal{P}} \hat{p} \hat{\mathcal{P}}^{-1}|\psi} &= \int_{-\infty}^{+\infty} \braket{-x|\hat{p}|x'} \braket{x'|\psi} \, dx' = \\
&= i \hbar \int_{-\infty}^{+\infty} \frac{d}{d(-x')} \delta(x - x') \psi(x') \, dx' = \\
&= \int_{-\infty}^{+\infty} \frac{d}{dx'} \delta(x + x') \psi(-x') \, dx' = \\
&= i \hbar \int_{-\infty}^{+\infty} \delta(x + x') \frac{d}{dx'} \psi(-x') \, dx' = \\
&= -i \hbar \frac{d}{dx} \psi(x) = -\braket{x|\hat{p}|\psi}
\end{align*}
da cui risulta $ \hat{\mathcal{P}} \hat{x} \hat{\mathcal{P}}^{-1} = -\hat{x} $ e $ \hat{\mathcal{P}} \hat{p} \hat{\mathcal{P}}^{-1} = -\hat{p} $. Da ci\`o segue che l'operatore parit\`a non commuta con gli operatori posizione e impulso. Esplicitamente
\[
[\hat{x}, \hat{\mathcal{P}}] = \hat{x} \hat{\mathcal{P}} - \hat{\mathcal{P}} \hat{x} = 2 \hat{x} \hat{\mathcal{P}}, \quad [\hat{p}, \hat{\mathcal{P}}] = \hat{p} \hat{\mathcal{P}} - \hat{\mathcal{P}} \hat{p} = 2 \hat{p} \hat{\mathcal{P}}.
\]

Scrivendo la formula completa dell'operatore hamiltoniano come 
\[
\mathcal{H} = \frac{\hat{p}^2}{2m} + V(\hat{x}) = \frac{\hat{p}^2}{2m} + \sum_j V_j \hat{x}^j
\]
e studiando l'operatore composto $ \hat{\mathcal{P}} \mathcal{H} \hat{\mathcal{P}}^{-1} $, si ha
\[
\hat{\mathcal{P}} \hat{\mathcal{H}} \hat{\mathcal{P}}^{-1} = \frac{1}{2m} \hat{\mathcal{P}} \hat{p} \hat{\mathcal{P}}^{-1} \hat{\mathcal{P}} \hat{p} \hat{\mathcal{P}}^{-1} + \sum_j \hat{\mathcal{P}} \hat{x}^j \hat{\mathcal{P}}^{-1} = \frac{\hat{p}^2}{2m} + \sum_j (-1)^j V_j \hat{x}^j
\]
che \`e uguale ad $ \mathcal{H} $ se e solo se il potenziale \`e una funzione pari (cfr. particella libera, oscillatore armonico, ecc.).

Nel caso degli operatori creazione e distruzione definiti da 
\[
\hat{a} = \sqrt{\frac{m\omega}{2\hbar}} \left( \hat{x} + i \frac{\hat{p}}{m\omega} \right), \qquad \hat{a}^\dag = \sqrt{\frac{m\omega}{2\hbar}} \left( \hat{x} - i \frac{\hat{p}}{m\omega} \right)
\]
si vede subito che 
\[
\hat{\mathcal{P}} \hat{a} \hat{\mathcal{P}}^{-1} = -\hat{a}, \qquad \hat{\mathcal{P}} \hat{a}^\dag \hat{\mathcal{P}}^{-1} = -\hat{a}^\dag
\]
quindi presi gli autostati $ \ket{n} $ dell'energia nell'oscillatore armonico, sapendo che 
\[
\hat{\mathcal{P}} \ket{0} = -\hat{a} \hat{\mathcal{P}} \ket{1} = \hat{a} \hat{a}^\dag \hat{\mathcal{P}} \ket{0} = \left( \frac{\hat{\mathcal{H}}}{\hbar \omega} + \frac{1}{2} \hat{\mathbb{I}} \right) \hat{\mathcal{P}} \ket{0} = \left( \frac{1}{2} \hat{\mathbb{I}} \frac{1}{2} \hat{\mathcal{P}} \right) \ket{0},
\]
cio\`e $ \hat{\mathcal{P}} \ket{0} = \ket{0} $, si ha
\[
\hat{\mathcal{P}} = \frac{1}{\sqrt{n!}} \hat{\mathcal{P}} \left( \hat{a}^\dag \right)^n \ket{0} = \frac{(-1)^n}{\sqrt{n!}} \left( \hat{a}^\dag \right)^n \ket{0} = (-1)^n \ket{n}.
\]

Esempio: ancora sul commutatore $ [\hat{p}, \mathcal{P}] $.
\[
\braket{q|[\hat{p}, \mathcal{P}]|\psi} = \braket{q|\hat{p} \mathcal{P}|\psi} - \braket{q|\mathcal{P} \hat{p}|\psi}
\]
\[
[\hat{p}, \mathcal{P}] = \mathcal{P} (\mathcal{P}\hat{p}\mathcal{P} - \hat{p}) = -2 \mathcal{P}\hat{p}
\]
\[
\braket{q|[\hat{p}, \mathcal{P}]|\psi} = -2 \braket{q|\mathcal{P}\hat{p}|\psi} = -2\braket{-q|\hat{p}|\psi} = 2i \hbar \frac{d}{dq} \psi(-q)
\]
% \subsection{Buca di potenziale a pareti infinite}
\[
\mathcal{H} = \frac{\hat{p}^2}{2m} + V(x), \quad V(x) = \left\{ \begin{matrix*}[l]
0 & |x| < a \\
V_0 \to +\infty & |x| > a
\end{matrix*} \right.
\]

La particella è intrappolata nella buca indipendentemente dall'energia che può avere e non supera la barriera (questa è la teoria classica).

Voglio sapere lo spettro di autostati di \( \mathcal{H} \) nella base delle coordinate. Allora:
\[
\mathcal{H} \ket{\psi_E} = E \ket{\psi_E}
\]
dove so \( \braket{x|\psi_E} = \psi_E(x) \).

Per \( |x| < a \)
\[
\braket{x|\mathcal{H}|\psi_E} = - \frac{\hbar^2}{2m} \frac{d^2}{dx^2} \psi_E(x) = E \psi_E(x).
\]

Per \( |x| > a \)
\[
\braket{x|\mathcal{H}|\psi_E} = - \frac{\hbar^2}{2m} \frac{d^2}{dx^2} \psi_E(x) = (E - V_0) \psi_E(x).
\]
\[
V_0 \to +\infty \Rightarrow \frac{d^2}{dx^2} \psi_E(x) = -\frac{2m}{\hbar^2} (E - V_0) \psi_E(x) \to +\infty, \forall E \Rightarrow \psi_E(x) = 0.
\]
Infatti fuori dalla buca una soluzione dell'equazione differenziale (con normalizzazione) è
\[
\psi_E(x) = N_E \exp \left\{ \pm \sqrt{\frac{2m}{\hbar^2}(V_0 - E)} x \right\}, \quad |x| > a.
\]
Deve restare valida la condizione
\[
\int_{-\infty}^{+\infty} |\psi(x)|^2 \, dx < +\infty
\]
quindi
\[
\psi_E(x) = N_E \exp \left\{ -\sqrt{\frac{2m}{\hbar^2}(V_0 - E)} x \right\} \quad (x > a),
\]
\[
\psi_E(x) = N_E \exp \left\{ \sqrt{\frac{2m}{\hbar^2}(V_0 - E)} x \right\} \quad (x < -a)
\]
Allora in generale ponendo \( k = \sqrt{\frac{2m}{\hbar^2}(V_0 - E)} \),
\[
\psi_E(x) = N_E e^{-k |x|} \quad (x > a)
\]
che per \( V_0 \to +\infty \) dà \( \psi_E \to 0 \). Questo impone come condizioni al contorno \( \psi_E(\pm a) = 0 \).

Per \( |x| < a \) invece
\[
\psi_E(x) = A_E' e^{ikx} + B_E' e^{-ikx}, \qquad E = \frac{\hbar^2 k^2}{2m}
\]
ovvero una combinazione lineare di seno e coseno
\[
\psi_E(x) = A_E \sin(kx) + B_E \cos(kx)
\]
e si presentano due casi.
\begin{enumerate}
\item \( B_E = 0 \Rightarrow \psi_E^I(x) = A_E \sin(kx) \) quindi per \( x = \pm a \) vale
\begin{align*}
\sin(ka) = -\sin(-ka) = 0 &\Rightarrow ka = \pi n \Rightarrow k_n = \frac{\pi}{a} n \Rightarrow \\
&\Rightarrow E_n = \frac{\hbar^2 \pi^2}{2m} \frac{n^2}{a^2}
\end{align*}
\item \( A_E = 0 \Rightarrow \psi_E^{II}(x) = B_E \cos(kx) \) e per \( x = \pm a \) ho
\begin{align*}
\cos(ka) = 0 \Rightarrow ka = (2n + 1) \frac{\pi}{2} &\Rightarrow k_n = \frac{\pi}{2a}(2n + 1) \Rightarrow \\ 
&\Rightarrow E_n = \frac{\hbar^2 \pi^2}{2m} \frac{(2n + 1)^2}{4a^2}
\end{align*}
\end{enumerate}

Quindi lo \textbf{spettro degli autovalori} è
\[
E_n = \frac{\hbar^2 k_n^2}{2m}, \qquad k_n = \frac{n\pi}{2a}
\]
e nel caso in cui \( n \) è pari ho la soluzione \( \psi_E^I \), mentre quando \( n \) è dispari ho la \( \psi_E^{II} \). Poiché infine \( L = 2a \) è la lunghezza della buca
\[
E_n = \frac{\hbar^2 n^2 \pi^2}{2 m L^2}.
\]
è anche possibile definire una \textbf{lunghezza d'onda} per le \( \psi_E \):
\[
\lambda_n = \frac{2\pi}{k_n} = \frac{2L}{n} \quad \Rightarrow \quad \frac{\lambda_n}{2} = \frac{L}{n}
\]
cioè la buca contiene \( n \) semilunghezze d'onda.

Definiamo allora lo \textbf{stato fondamentale} \( \psi^{II} \) come
\[
n = 1, \quad k = \frac{\pi}{2a}, \quad E = \frac{\hbar^2 \pi^2}{2 m L^2}
\]
e lo \textbf{stato eccitato} \( \psi^I \) per
\[
n = 2, \quad k = \frac{\pi}{a}, \quad E = \frac{2 \hbar^2 \pi^2}{m L^2}.
\]
Gli andamenti sono riportati in Figura \ref{buca_infinita}.
\begin{figure}
\centering
\begin{tikzpicture}
  \draw[color=red, domain=-pi/2:pi/2, samples=200, smooth] plot(\x,{cos(\x r)}) node[right] {};
  \node (A) at (-1.2,1) {\(n = 1\)};

  \draw[color=blue, domain=-pi/2:pi/2, samples=200, smooth] plot(\x,{sin(2*\x r)}) node[right] {};
  \node (B) at (-1,-1.2) {\(n = 2\)};

  \draw[thick] (-3,0) -- (3,0);
  \draw[thick] (0,-1.5) -- (0,1.5);

  \node (C) at (pi/2+.2,-.2) {\(a\)};
  \node (D) at (-pi/2-.3,-.2) {\(-a\)};
\end{tikzpicture}
\caption{Andamento dello stato fondamentale (rosso) e del primo stato eccitato (blu) per la buca di potenziale a pareti infinite.}
\label{buca_infinita}
\end{figure}

Pensando alle soluzioni come azione dell'operatore parità sulla funzione d'onda si ha
\[
\braket{x|\mathcal{P}|\psi} = \braket{-x|\psi} \; \Rightarrow \; \mathcal{P} \ket{\psi_E} = \begin{cases}
+ \ket{\psi_E} \text{ se ho } \psi^{II} \\
- \ket{\psi_E} \text{ se ho } \psi^I 
\end{cases}
\]
Inoltre \([\mathcal{H}, \mathcal{P}] = 0\), ed essendo \( \mathcal{P}^{-1} \mathcal{H} \mathcal{P} = \mathcal{H} \), le autofunzioni \( \psi_E \) sono gli autostati di \( \mathcal{H} \) e di \( \mathcal{P} \).

Era possibile dedurre lo stato ad energia minima dal principio di indeterminazione perché 
\[
\begin{cases}
\Delta x \leq L \Rightarrow \Delta^2 x \leq L^2 \\
\Delta^2 \hat{p} \Delta^2 \hat{x} \geq \frac{\hbar^2}{4}
\end{cases} \; \Rightarrow \; \Delta^2 \hat{p} \geq \frac{\hbar^2}{4L^2}
\]
avendo che l'impulso ha un valore minimo. Dunque non posso che avere un'energia minima.
% \subsection{Potenziale a gradino}
\begin{figure}
\centering
\begin{tikzpicture}
  \draw (-5,0) -- (5,0);
  \draw (0,-1) -- (0,3);

  \draw[thick] (-4, 0) -- (0,0);
  \draw[thick] (0,2) -- (4,2);

  \node (1) at (-2.5,2.5) {1};
  \node (2) at (2.5,2.5) {2};
  \node (V0) at (-.3,2) {$V_0$};
\end{tikzpicture}
\caption{Schema del potenziale a gradino.}
\label{gradino}
\end{figure}
Lo schema \`e riportato in Figura \ref{gradino}.
$$
V(x) = \left\{ \begin{matrix}
V_0 & x > 0 \\
0 & x < 0
\end{matrix} \right..
$$
Supponendo che
$$
\delta(x) = \frac{1}{V_0} \frac{d}{dx} V(x),
$$
perch\'e
\begin{align*}
\int_{-\infty}^{+\infty} V(x) f'(x) \, dx &= \left[ V(x) f(x) \right]_{-\infty}^{+\infty} - \int_{-\infty}^{+\infty} f(x) \frac{d}{dx}V(x) \, dx = \\
&= f(+\infty) - \int_{-\infty}^{+\infty} f(x) \frac{d}{dx}V(x) \, dx \\
\int_{-\infty}^{+\infty} V(x) f'(x) \, dx &= \int_0^{+\infty} f'(x) \, dx = f(+\infty) - f(0).
\end{align*}

Scrivo l'equazione di Schr\"odinger
$$
-\frac{\hbar^2}{2m} \psi''(x) + V(x) \psi(x) = E \psi(x)
$$
$$
\psi''(x) = -\frac{2m}{\hbar^2} (E - V(x)) \psi(x).
$$
Allora la $ \psi'' $ \`e discontinua perch\'e lo \`e $ V $, ma quindi $ \psi $ e $ \psi' $ sono continue, cos\`{i} posso dividere il problema in due parti e poi imporre delle condizioni di raccordo.

Per $ x < 0 $ ho $ \psi_1 $ tale che 
$$
\psi_1''(x) = -\frac{2m}{\hbar^2} E \psi_1(x) \; \Rightarrow \; \psi_1(x) = A e^{ikx} + B e^{-ikx}, \; k = \frac{\sqrt{2mE}}{\hbar}.
$$

Per $ x > 0 $
\begin{align*}
\psi_2''(x) &= -\frac{2m}{\hbar^2} (E - V_0) \psi_2(x) \; \Rightarrow \\
&\; \Rightarrow \; \psi_2(x) = C e^{ik'x} + D e^{-ik'x}, \; k' = \sqrt{\frac{2m}{\hbar^2}(E - V_0)}
\end{align*}
e per ora impongo $ E > V_0 $.

Poich\'e $ \psi_1(0) = \psi_2(0) $ e $ \psi_1'(0) = \psi_2'(0) $, allora mi restano due condizioni da imporre. Sfrutto la normalizzazione per la terza equazione e poi impongo $ D = 0 $. Allora 
$$
\begin{cases}
A + B = C \\
ik(A - B) = ik'C 
\end{cases} \quad \Rightarrow \quad \begin{cases}
B = \frac{k - k'}{k + k'} A \\
C = \frac{2k}{k + k'} A
\end{cases}.
$$

Introduco la \textbf{corrente di probabilit\`a}
$$
J(x, t) = -\frac{i\hbar}{2m} \left( \psi^*(x, t) \frac{\partial}{\partial x} \psi(x, t) - \left( \frac{\partial}{\partial x} \psi^*(x, t) \right) \psi(x, t) \right)
$$
tale che
$$
\frac{\partial}{\partial x} J(x, t) + \frac{\partial}{\partial t} \rho(x, t) = 0, \quad \text{con } \rho(x, t) = |\psi(x, t)|^2 = |\braket{x|\psi(t)}|^2.
$$
Infatti
$$
\frac{\partial J}{\partial x} = -\frac{i\hbar}{2m} \left( \cancel{\frac{\partial \psi^*}{\partial x} \frac{\partial \psi}{\partial x}} + \psi^* \frac{\partial^2 \psi}{\partial x^2} - \frac{\partial^2 \psi^*}{\partial x^2} \psi - \cancel{\frac{\partial \psi^*}{\partial x} \frac{\partial \psi}{\partial x}} \right)
$$
poi per l'equazione di Schr\"odinger
$$
\frac{\partial^2 \psi}{\partial x^2} = \left( i\hbar\frac{\partial \psi}{\partial t} -V(x)\psi \right) \left( -\frac{2m}{\hbar^2} \right).
$$

Essendo $ \mathcal{H} $ hermitiano, $ \mathcal{H} = \mathcal{H}^\dag $, quindi 
$$
\frac{\partial^2 \psi^*}{\partial x^2} = \left( -i\hbar\frac{\partial \psi^*}{\partial t} - V(x) \psi^* \right) \left( -\frac{2m}{\hbar^2} \right)
$$
allora 
\begin{align*}
\frac{\partial}{\partial x} J(x, t) &= -\frac{i\hbar}{2m} \left( \psi^* \left( i\hbar\frac{\partial \psi}{\partial t} - V \psi \right) - \left( -i\hbar\frac{\partial \psi^*}{\partial t} - V \psi^* \right) \psi \right) \left( -\frac{2m}{\hbar^2} \right) = \\
&= - \left( \psi^* \frac{\partial \psi}{\partial t} + \frac{\partial \psi^*}{\partial t} \psi \right) = - \frac{\partial}{\partial t} \left( \psi^* \psi \right) = \\
&= -\frac{\partial}{\partial t} \rho(x, t).
\end{align*}

Integrando
$$
\int_a^b \frac{\partial \rho}{\partial t} \, dx = -\int_a^b \frac{\partial J}{ \partial x} \, dx = - \left( J(b) - J(a) \right)
$$
$$
\frac{\partial \rho}{\partial t} > 0 \qquad \Rightarrow \qquad J(a) > J(b).
$$

Il valore medio di $ J $
\begin{align*}
\int_{-\infty}^{+\infty} J(x, t) \, dx &= \frac{1}{m} \int_{-\infty}^{+\infty} \braket{\psi|x} \braket{x|\hat{p}|\psi} \, dx = \\
&= \frac{1}{m} \braket{\psi|\hat{p}|\psi} = \\
&= \frac{\braket{\hat{p}}_\psi}{m}
\end{align*}
corrispondente al valore medio di $ \hat{p} $ diviso per $ m $
$$
\braket{J} = \frac{1}{m} \braket{k|\hat{p}|k} = \frac{k}{m} \braket{k|k}.
$$
Quindi data
$$
\psi(x, t) = A \exp \left\{ \frac{i}{\hbar} (kx - \omega t) \right\} + B \exp \left\{ -\frac{i}{\hbar} (kx - \omega t) \right\}
$$
la $ J $ associata \`e 
\begin{align*}
J(x, t) &= -\frac{i\hbar}{2m} \left( \left( A^* \exp \left\{ -\frac{i}{\hbar} (kx - \omega t) \right\} + B^* \exp \left\{ \frac{i}{\hbar} (kx - \omega t) \right\} \right) \right. \cdot \\
&\cdot \frac{ik}{\hbar} \left( A \exp \left\{ \frac{i}{\hbar} (\dots) \right\} + B \exp \left\{ -\frac{i}{\hbar} (\dots) \right\} \right) + \\
&- \frac{ik}{\hbar} \left( -A^* \exp \left\{ -\frac{i}{\hbar} (\dots) \right\} + B^* \exp \left\{ -\frac{i}{\hbar} (\dots) \right\}  \right) \cdot \\
&\cdot \left( A \exp \left\{ \frac{i}{\hbar} (\dots) \right\} + B \exp \left\{ -\frac{i}{\hbar} (\dots) \right\} \right) = \\
&= \frac{k}{2m} \left( |A|^2 - |B|^2 \right) 2 = \frac{k}{m} \left( |A|^2 - |B|^2 \right) = \\
&= J_A(x, t) + J_B(x, t).
\end{align*}

Poich\'e $ \ket{\psi} = \ket{k} \braket{k|\psi} + \ket{-k} \braket{-k|\psi} $, allora 
$$
\braket{J} = \frac{k}{m} \left( |\braket{k|\psi}|^2 - |\braket{-k|\psi}|^2 \right).
$$
Inoltre, essendo 
$$
\rho = |\psi|^2 \; \Rightarrow \; \frac{\partial \rho}{\partial t} = 0 \; \Rightarrow \; \frac{\partial J}{\partial x} = 0,
$$
quindi $ J $ \`e invariante per spostamento, quindi 
$$
J_A + J_B = J_C \text{ costante}
$$
e questa condizione deve valere nella prima e nella seconda regione su cui poi impongo le equazioni di continuit\`a. Allora 
$$
J_{\text{inc.}} + J_{\text{rifl.}} = J_{\text{trasm.}}
$$
e i coefficienti
$$
T = \left| \frac{C}{A} \right|^2, \qquad R = \left| \frac{B}{A} \right|^2.
$$

Nel caso in cui $ E < V_0 $
$$
k' = \sqrt{\frac{2m}{\hbar}(E - V_0)} \quad \Rightarrow \quad k' = i \beta, \, \beta \in \mathbb{R}
$$
quindi $ \psi_2(x) = C e^{-\beta x} $, cio\`e l'andamento \`e esponenziale e la corrente $ J_C(x, t) = 0 $, quindi $ J_A = - J_B $ e 
$$
R = \left| \frac{B}{A} \right|^2 = \left| \frac{k - i\beta}{k + i\beta} \right|^2 = 1.
$$
Se $ V_0 \to +\infty $, allora $ \beta \to +\infty $, allora $ \psi_2(x) = 0 $, quindi essendo $ \psi_1(0) = \psi_2(0) $, $ A + B = 0 $, $ \psi_1(x) = A' (e^{ikx} - e^{-ikx}) $.
% \subsection{Barriera di potenziale}
\begin{figure}
\centering
\begin{tikzpicture}
  \draw (-5,0) -- (5,0);
  \draw (0,-1) -- (0,3);

  \draw[thick] (-4,0) -- (-2,0);
  \draw[thick] (2,0) -- (4,0);
  \draw[thick] (-2, 2) -- (2,2);

  \draw[dotted] (-2,-1) -- (-2,3);
  \draw[dotted] (2,-1) -- (2,3);

  \node (-a) at (-2.3,-.3) {$-a$};
  \node (a) at (2.2,-.3) {$a$};

  \node (I) at (-3,2.5) {I};
  \node (II) at (.5,2.5) {II};
  \node (II) at (3,2.5) {III};
\end{tikzpicture}
\caption{Schema della barriera di potenziale.}
\label{barriera}
\end{figure}
Nella regione I e III ho la particella libera (Figura \ref{barriera}), quindi
$$
\frac{\hat{p}^2}{2m}\ket{\psi} = E \ket{\psi} \quad \Rightarrow \quad \psi''(x) = -\frac{2m}{\hbar^2} E \psi(x),
$$
mentre nella regione II 
$$
\left( \frac{\hat{p}^2}{2m} + V_0 \right) \ket{\psi} = E \ket{\psi} \quad \Rightarrow \quad \psi''(x) = \frac{2m}{\hbar^2}(V_0 - E) \psi(x).
$$
e come nel caso del gradino, ho varie situazioni a seconda che la particella abbia energia maggiore o minore di $ V_0 $.

Si tratta di raccordare le soluzioni delle tre regioni, ma solo quando $ E > V_0 $? In particolare, c'\`e una probabilit\`a che per $ E < V_0 $ la particella passi al di l\`a di una barriera?

La forma della $ \psi $ \`e data da 
$$
\psi_I = A e^{-ikx} + B e^{ikx}, \quad \psi_{III} = F e^{-ikx} + G e^{ikx} \quad (\text{sempre})
$$
$$
\psi_{II} = C e^{-ik'x} + D e^{ik'x}, \qquad \text{se } E > V_0
$$
$$
\psi_{II} = C e^{-\beta x} + D e^{\beta x}, \qquad \text{se } E < V_0
$$
con
$$
k = \sqrt{\frac{2mE}{\hbar^2}}, \quad k' = \sqrt{\frac{2m(E - V_0)}{\hbar^2}}, \quad \beta = \sqrt{\frac{2m(V_0 - E)}{\hbar^2}}
$$
con le condizioni ai bordi
$$
\psi_I(-a) = \psi_{II}(-a), \; \psi_I'(-a) = \psi_{II}'(-a), \; \psi_{II}(a) = \psi_{III}(a), \; \psi_{II}'(a) = \psi_{III}'(a).
$$

Anche in questo caso ho sei costanti arbitrarie, le quattro condizioni pi\`u quella di normalizzazione, quindi resta un parametro: allora ho due soluzioni indipendenti che si combinano per dare la soluzione generica. Mettiamo a 0 una delle due costatnti di $ \psi_{III} $ a seconda che l'onda provenga da destra o da sinistra. Poniamo che sia $ F = 0 $.

Vogliamo gli indici di riflessione e trasmissione, quindi 
$$
\begin{cases}
B \text{ incidente} \\
A \text{ riflessa} \\
G \text{ trasmessa}
\end{cases}, \quad T = \frac{J_T}{J_I}, \quad R = \frac{J_R}{J_I}.
$$
Poich\'e le $ J $ e le $ \rho $ sono termini del tipo $ \psi^* \psi $, allora 
$$
T = \left| \frac{G}{B} \right|^2, \qquad R = \left| \frac{A}{B} \right|^2
$$
e poi $ J_I = J_R + J_T $ e costruisco $ R $, ecc. Operando le sostituzioni
$$
\xi = \frac{\beta}{k} - \frac{k}{\beta}, \qquad \eta = \frac{\beta}{k} + \frac{k}{\beta}
$$
ho i coefficienti
$$
T = \left[ 1 + \left( 1 + \frac{\xi^2}{4} \right) + \sinh(2\beta a) \right]^{-1}
$$
Se $ \beta a \gg 1 $, allora
$$
T \simeq \alpha e^{-2\beta a}
$$
cio\`e quando la barriera \`e molto grande rispetto a $ E $ la probabilit\`a che la particella passi dall'altra parte della barriera diminuisce (effetto tunnel).

Il caso di $ E > V_0 $ \`e quello della trasmissione della buca attraverso due mezzi diversi.
% \subsection{Trattazione qualitativa di un particolare potenziale}
\begin{figure}
\centering
\begin{tikzpicture}
  \draw[->] (-5,0) -- (5,0);
  \draw[->] (0,-3) -- (0,1);

  \draw[thick,smooth] plot file {./grafici/potenziale_particolare.dat};
\end{tikzpicture}
\caption{Andamento di un potenziale simmetrico, nullo all'infinito e con un punto di equilibrio stabile.}
\label{particolare}
\end{figure}
Il potenziale (Figura \ref{particolare}) è attrattivo e tale che
\[
V(x) = V(-x), \quad \lim_{x \to \pm \infty} V(x) = 0, \quad V(x) < 0.
\]
Quindi il solito \( \hat{\mathcal{H}} \ket{\psi} = E \ket{\psi} \), cioè 
\[
\psi''(x) = \frac{2m}{\hbar^2} \left( V(x) - E \right) \psi(x).
\]
Da quello che già conosco mi aspetto uno spettro discreto non degenere, cioè degli stati legati (cioè normalizzabili), per quello che ottengo dalla buca infinita di potenziale. Inoltre ci sarà anche una parte continua dello spettro (cioè stati non normalizzabili).

\bigskip
\noindent
\( \boxed{V(x) > E} \) se \( \psi > 0 \), allora \( \psi'' > 0 \); se \( \psi < 0 \), allora \( \psi'' < 0 \) (Figura \ref{potenziale_generico_V_gt_E}).
\begin{figure}
\centering
\begin{tikzpicture}
  \draw (-2,3) -- (-2,-3);
  \node[draw] at (1,3) {\( \psi > 0 \)};
  \draw (-1.8,0.5) -- (1.8,0.5);
  \draw (0,0.2) -- (0,3);
  \draw[thick] (.5,2) .. controls (.75,1.5) and (1,1) .. (1.7,.59);
  \draw[thick] (-1.7,.59) .. controls (-1,1) and (-.75,1.5) .. (-.5,2);

  \draw (-1.8,-2.7) -- (1.8,-2.7);
  \draw (0,-0.2) -- (0,-3);
  \draw[thick] (.5,-2.52) .. controls (.75,-2.5) and (1.25,-2.3) .. (1.7,-1.64);
  \draw[thick] (-.5,-2.52) .. controls (-.75,-2.5) and (-1.25,-2.3) .. (-1.7,-1.64);

  \draw (2,3) -- (2,-3);
  \node[draw] at (5,3) {\( \psi < 0 \)};
  \draw (2.2,2.5) -- (5.8,2.5);
  \draw (4,0.2) -- (4,3);
  \begin{scope}[shift={(4,2.5)}]
    \draw[thick] (.5,-2) .. controls (.75,-1.5) and (1,-1) .. (1.7,-.59);
    \draw[thick] (-.5,-2) .. controls (-.75,-1.5) and (-1,-1) .. (-1.7,-.59);
  \end{scope}

  \draw (2.2,-0.7) -- (5.8,-0.7);
  \draw (4,-0.2) -- (4,-3);
  \begin{scope}[shift={(4,0)}]
    \draw[thick] (.5,-1.08) .. controls (.75,-1.2) and (1.25,-1.5) .. (1.7,-1.96);
    \draw[thick] (-.5,-1.08) .. controls (-.75,-1.2) and (-1.25,-1.5) .. (-1.7,-1.96);
  \end{scope}

  \node[draw] at (-3,3) {\( V > E \)};
  \draw (-5.8,0) -- (-2.2,0);
  \draw (-5.8,-.5) -- (-2.2,-.5);
  \node (E) at (-2.5,-.8) {\(E\)};
  \draw (-4,-1) -- (-4,.5);
  \begin{scope}[scale=0.4, shift={(-10,0)}]
    \draw[thick,smooth] plot file {./grafici/potenziale_particolare.dat};
  \end{scope}
\end{tikzpicture}
\caption{Trattazione qualitativa delle funzioni d'onda per potenziali \( V(x) > E \).}
\label{potenziale_generico_V_gt_E}
\end{figure}
\( \boxed{V(x) < E} \) se \( \psi > 0 \), allora \( \psi'' < 0 \); se \( \psi < 0 \), allora \( \psi'' > 0 \) (Figura \ref{potenziale_generico_V_lt_E}).
\begin{figure}
\centering
\begin{tikzpicture}
  \draw (-2,3) -- (-2,-3);
  \node[draw] at (1,3) {\( \psi > 0 \)};
  \draw (-1.8,0.5) -- (1.8,0.5);
  \draw (0,0.2) -- (0,3);
  \draw[thick] (.5,1.92) .. controls (.75,1.85) and (1.25,1.6) .. (1.7,1.04);
  \draw[thick] (-.5,1.92) .. controls (-.75,1.85) and (-1.25,1.6) .. (-1.7,1.04);

  \draw (-1.8,-2.7) -- (1.8,-2.7);
  \draw (0,-0.2) -- (0,-3);
  \begin{scope}[shift={(4,1)}]
    \draw[thick] (.5,.08) .. controls (.75,.1) and (1.5,.75) .. (1.7,.96);
    \draw[thick] (-.5,.08) .. controls (-.75,.1) and (-1.5,.75) .. (-1.7,.96);
  \end{scope}
  \draw (2,3) -- (2,-3);
  \node[draw] at (5,3) {\( \psi < 0 \)};
  \draw (2.2,2.5) -- (5.8,2.5);
  \draw (4,0.2) -- (4,3);
  \begin{scope}[shift={(0,-0.5)}]
    \draw[thick,smooth] (.5,-2) .. controls (.75,-1.25) and (1.25,-.75) .. (1.7,-.59);
    \draw[thick,smooth] (-.5,-2) .. controls (-.75,-1.25) and (-1.25,-.75) .. (-1.7,-.59);
  \end{scope}

  \draw (2.2,-0.7) -- (5.8,-0.7);
  \draw (4,-0.2) -- (4,-3);
  \begin{scope}[shift={(4,-3)}]
    \draw[thick,smooth] (.5,2) .. controls (.75,1.25) and (1.25,.75) .. (1.7,.59);
    \draw[thick,smooth] (-.5,2) .. controls (-.75,1.25) and (-1.25,.75) .. (-1.7,.59);
  \end{scope}

  \node[draw] at (-3,3) {\( V < E \)};
  \draw (-5.8,0) -- (-2.2,0);
  \draw (-5.8,-.5) -- (-2.2,-.5);
  \node (E) at (-2.5,-.8) {\(E\)};
  \draw (-4,-1) -- (-4,.5);
  \begin{scope}[scale=0.4, shift={(-10,0)}]
    \draw[thick,smooth] plot file {./grafici/potenziale_particolare.dat};
  \end{scope}
\end{tikzpicture}
\caption{Trattazione qualitativa delle funzioni d'onda per potenziali \( V(x) < E \).}
\label{potenziale_generico_V_lt_E}
\end{figure}

\bigskip
\noindent
Se \( V(x) - E > 0 \) sempre, allora ho Figura \ref{V_gt_E_sempre} quindi non c'è alcuna soluzione.
\begin{figure}
\centering
\begin{tikzpicture}
  \draw (-5,0) -- (-1,0);
  \draw (-3,3) -- (-3,-.5);
  \begin{scope}[shift={(-3,1)}]
    \draw[thick] (0,1) .. controls (.5,.5) and (1.5,.25) .. (2,0.13);
    \draw[thick] (0,1) .. controls (-.5,.5) and (-1.5,.25) .. (-2,0.13);
  \end{scope}
  \node (A) at (-3,-1) {non derivabile};

  \draw (1,0) -- (5,0);
  \draw (3,3) -- (3,-.5);
  \begin{scope}[shift={(3,1)}]
    \draw[thick] (-2,1) .. controls (-1,-.5) and (1,-.5) .. (2,1);
  \end{scope}
  \node (B) at (3,-1) {non normalizzabile};
\end{tikzpicture}
\caption{Soluzioni nel caso \(V(x) > E\) sempre. Nessuna di esse è accettabile.}
\label{V_gt_E_sempre}
\end{figure}

\bigskip
\noindent
Se \( V(x) - E > 0 \) o \( V(x) - E < 0 \) a seconda delle regioni, ho Figura \ref{potenziale_variabile}.
\begin{figure}
\centering
\begin{tikzpicture}
  \draw (-5,0) -- (5,0);
  \draw (0,3) -- (0,-3);

  \draw[thick,smooth] plot file {./grafici/potenziale_particolare_1.dat};
  \draw[thick,dashed,smooth] plot file {./grafici/potenziale_particolare_2.dat};
  \draw[thick,dotted,smooth] plot file {./grafici/potenziale_particolare_3.dat};

  \draw[dashed] (-4,-3) -- (-4,3);
  \node (mx0) at (-4.5,-.5) {\(-x_0\)};

  \draw[dashed] (4,-3) -- (4,3);
  \node (x0) at (4.5,-.5) {\(x_0\)};

  \draw[thick, dotted] (.5,-2) -- (2,-2);
  \node (disc) at (1.5,-2.5) {discontinua};
\end{tikzpicture}
\caption{\( V(x) - E > 0 \) o \( V(x) - E < 0 \) a seconda delle regioni.}
\label{potenziale_variabile}
\end{figure}

\bigskip
\noindent
Poiché se \( E \) cresce, cresce anche \( x_0 \), lo stato a minima energia è quello che ha meno oscillazione.

\bigskip
\noindent
Nel caso \( V(x) > E \) sempre, ho soluzioni oscillanti \( \Rightarrow \) spettro continuo, cioè c'è almeno uno stato legato.

\bigskip
\noindent
Quindi ho:
\begin{enumerate}
\item almeno uno stato legato;
\item \( E_0 > V_0 \), quindi \(\Delta x\) è massimo solo localmente e \(\Delta p\) è minimo nello stato fondamentale;
\item spettro non degenere, quindi ogni operatore che commuta con \( \hat{\mathcal{H}} \) è diagonale nella base degli \( \ket{E} \), cioè 
\[
[A, \hat{\mathcal{H}}] = 0 \quad \Rightarrow \quad A \ket{E} = a_E \ket{E}
\]
\end{enumerate}
Infatti se ho \( \psi_1 \) e \( \psi_2 \) tali che
\[
\psi_1'' = \frac{2m}{\hbar^2} \left( V(x) - E \right) \psi_1, \quad \psi_2'' = \frac{2m}{\hbar^2} \left( V(x) - E \right) \psi_2
\]
e sono normalizzabili, ho che
\[
\psi_2 \psi_1'' - \psi_1 \psi_2'' = 0 = \frac{d}{dx} (\psi_2 \psi_1' - \psi_1 \psi_2') \; \Rightarrow \; \psi_2 \psi_1' - \psi_1 \psi_2' = \text{ costante}
\]
e poiché \( \psi_1 \) e \( \psi_2 \) sono normalizzate,
\[
\lim_{x \to \pm \infty} \psi_1(x) = \lim_{x \to \pm \infty} \psi_2(x) = 0
\]
Quindi la costante è 0. Allora 
\[
\frac{\psi_1'}{\psi_1} = \frac{\psi_2'}{\psi_2} \; \Rightarrow \; \frac{d}{dx} \ln \psi_1(x) = \frac{d}{dx} \ln \psi_2(x) \; \Rightarrow \; \ln \psi_1(x) = \ln \psi_2(x) + \text{ cost.} \; \Rightarrow
\]
\[
\; \Rightarrow \psi_1(x) = e^{\text{cost.}} \psi_2(x)
\]
quindi \( \psi_1 \) e \( \psi_2 \) sono la stessa autofunzione e ho spettro non degenere.
% \subsection{Oscillatore armonico}
Se prendo un potenziale continuo con un minimo locale stabile, allora
\[
V(x) = V(x_0) + \frac{1}{2}(x - x_0)^2 V''(x_0) + \mathcal{O}((x - x_0)^2),
\]
\[
\mathcal{H} = \frac{\hat{p}^2}{2m} + \frac{1}{2} m \omega^2 \hat{x}^2,
\]
pongo \( x_0 = 0 \) e \( V(x_0) = 0 \), allora determino gli autostati dell'energia
\[
\mathcal{H} \ket{n} = E_n \ket{n}
\]
e anche \( \braket{x|n} = \psi_n(x) \).

\textbf{Metodo di Heisenberg}: avrò alcune proprietà sugli \( E_n \).
\begin{enumerate}
\item \( \left( \frac{\hat{p}^2}{2m} + \frac{1}{2} m \omega^2 \hat{x}^2 \right) \ket{n} = E_n \ket{n}  \Rightarrow \braket{n|\left( \frac{\hat{p}^2}{2m} + \frac{1}{2} m \omega^2 \hat{x}^2 \right)|n} = E_n \braket{n|n} \). 
\[
E_n = \sum_k \frac{1}{2m} \braket{n|\hat{p}|k} \braket{k|\hat{p}|n} + \sum_k \frac{1}{2} m \omega^2 \braket{n|\hat{x}|k} \braket{k|\hat{x}|n} =
\]
\[
= \frac{1}{2m} \sum_k \left| \braket{n|\hat{p}|k} \right|^2 + \frac{1}{2} m \omega^2 \sum_k \left| \braket{n|\hat{x}|k} \right|^2 \geq 0.
\]
Ma se per assurdo \( E_n = 0 \Rightarrow \braket{n|\hat{p}|k} = \braket{n|\hat{x}|k} = 0, \forall \ket{k} \), quindi
\[
\sum_k \left( \braket{n|\hat{x}|k} \braket{k|\hat{p}|n} - \braket{n|\hat{p}|k} \braket{k|\hat{x}|n} \right) = 0,
\]
quindi \( [\hat{p}, \hat{x}] = 0 \), che è falso. Quindi \( \boxed{E_n > 0} \).
\item Sulla base degli \( \ket{n} \), voglio costruire \( \braket{n|\hat{x}|k} \) e \( \braket{n|\hat{p}|k} \). Costruisco i commutatori
\begin{align*}
[\hat{x}, \mathcal{H}] &= \frac{1}{2m} [\hat{x}, \hat{p}^2] = i\hbar \frac{\hat{p}}{m} \\
[\hat{p}, \mathcal{H}] &= \frac{1}{2} m \omega^2 [\hat{p}, \hat{x}^2] = -i\hbar m \omega^2 \hat{x}
\end{align*}
e poi so che
\begin{align*}
\braket{n|[\hat{x}, \mathcal{H}]|k} &=  \braket{n|(\hat{x}\mathcal{H} - \mathcal{H}\hat{x})|k} = \\
&= \braket{n|\hat{x}|k}(E_k - E_n) \\
&= i\hbar \braket{n|\frac{\hat{p}}{m}|k}
\end{align*}
\begin{align*}
\braket{n|[\hat{p}, \mathcal{H}]|k} &= \braket{n|\hat{p}|k} (E_k - E_n) \\
&= -i\hbar m \omega^2 \braket{n|\hat{x}|k}.
\end{align*}
Allora ho due equazioni, da cui 
\[
\braket{n|\hat{x}|k} = \frac{i\hbar}{m(E_k - E_n)} \frac{(-i\hbar m \omega^2)}{E_k - E_n} \braket{n|\hat{x}|k} \; \Rightarrow
\]
$$ \Rightarrow \; (E_k - E_n)^2 \braket{n|\hat{x}|k} = \hbar^2\omega^2 \braket{n|\hat{x}|k}
\[
cioè l'elemento
\]
\braket{n|\hat{x}|k} \begin{cases}
= 0 \text{ sulla diagonale e anche dove la differenza è da } (\hbar \omega) \\
\neq 0 \text{ dove } E_k - E_n = \pm \hbar \omega
\end{cases}
\[
Introduco l'operatore
\]
\hat{a} = \sqrt{\frac{m\omega}{2\hbar}} \left( \hat{x} + i\frac{\hat{p}}{m\omega} \right), \quad \hat{a}^\dag = \sqrt{\frac{m\omega}{2\hbar}} \left( \hat{x} - i\frac{\hat{p}}{m\omega} \right)
$$
\begin{align*}
[\hat{a}^\dag, \hat{a}] &= \frac{m\omega}{2\hbar} \left[ \hat{x} - i\frac{\hat{p}}{m\omega}, \hat{x} + i\frac{\hat{p}}{m\omega} \right] = \frac{m\omega}{\hbar} i \left[ \hat{x}, \frac{\hat{p}}{m\omega} \right] = -1 \\
\hat{a}^\dag\hat{a} &= \frac{m\omega}{2\hbar} \left( \hat{x} - i\frac{\hat{p}}{m\omega} \right) \left( \hat{x} + i\frac{\hat{p}}{m\omega} \right) = \\
&= \frac{m\omega}{2\hbar} \left( \hat{x}^2 + \frac{\hat{p}^2}{m^2\omega^2} - \frac{\hbar}{m\omega} \right) = \\
&= \frac{1}{\hbar\omega} \mathcal{H} - \frac{1}{2} \equiv N.
\end{align*}
\end{enumerate}
% \subsection{Spettro dell'oscillatore armonico}
Le relazioni che ho ottenuto sono 
\[
[\hat{a}^\dag, \hat{a}] = -1, \qquad \mathcal{H} = \hbar \omega \left( N + \frac{1}{2} \right).
\]
Voglio determinare lo spettro di $ \mathcal{H} $ e passo attraverso $ [N, \hat{a}] $ e $ [N, \hat{a}^\dag] $.
\[
[N, \hat{a}] = [\hat{a}^\dag \hat{a}, \hat{a}] = [\hat{a}^\dag, \hat{a}]\hat{a} = -\hat{a},
\]
\[
[N, \hat{a}^\dag] = [\hat{a}^\dag \hat{a}, \hat{a}^\dag] = \hat{a}^\dag[\hat{a}, \hat{a}^\dag] = \hat{a}^\dag.
\]

Se conosco un autostato di $ N $, sia $ \ket{z} $, allora
\[
N \ket{z} = \lambda_z \ket{z} \; \Rightarrow \; \begin{cases}
N \hat{a}^\dag \ket{z} = (\lambda_z + 1) \hat{a}^\dag \ket{z} \\
N \hat{a} \ket{z} = (\lambda_z - 1) \hat{a} \ket{z}
\end{cases}.
\]
Dimostrazione:
\begin{align*}
N \hat{a}^\dag \ket{z} &= [N \hat{a}^\dag = \hat{a}^\dag N + \hat{a}^\dag] = (\hat{a}^\dag N + \hat{a}^\dag) \ket{z} = \\
&= \hat{a}^\dag \lambda_z \ket{z} + \hat{a}^\dag \ket{z} = (\lambda_z + 1) \hat{a}^\dag \ket{z} \\
N \hat{a} \ket{z} &= (\hat{a}N - \hat{a}) \ket{z} = (\lambda_z - 1) \hat{a} \ket{z}. \qed
\end{align*}

Inoltre, $ N \ket{z} = \lambda_z \ket{z} \Rightarrow \lambda_z \geq 0 $.

Dimostrazione: $ \braket{z|N|z} = \lambda_z $ quando $ \braket{z|z} = 1 $, quindi
\begin{align*}
\lambda_z &= \braket{z|\hat{a}^\dag \hat{a}|z} = \\
&= \sum_n \braket{z|\hat{a}^\dag|n} \braket{n|\hat{a}|z} = \\
&= \sum_n \left| \braket{n|\hat{a}|z} \right|^2 \geq 0. \qed
\end{align*}

Deve esistere $ \ket{0} $ tale che $ \hat{a} \ket{0} = 0 $, perch\'e se non esistesse, cosiderando
\[
N \hat{a}^P \ket{z} = (\lambda_z - P) \hat{a}^P \ket{z},
\]
per $ P $ abbastanza grande avrei autovalori negativi. Inoltre, se $ \hat{a} \ket{0} = 0 $, allora $ N \ket{0} = \hat{a}^\dag \hat{a} \ket{0} = \lambda_0 \ket{0} $, cio\`e $ \lambda_0 = 0 $.

Allora posso definire
\[
\ket{n} = K_n (\hat{a}^\dag)^n \ket{0}
\]
$ K_n $ costante di normalizzazione tale che 
\[
N \ket{n} = n \ket{n},
\]
\[
\hat{a}^\dag \ket{n} = \frac{K_n}{K_{n + 1}} \ket{n + 1}, \qquad \hat{a} \ket{n} = n \frac{K_n}{K_{n - 1}}\ket{n - 1}.
\]
Dico che $ \ket{0} $ \`e il \textbf{vuoto}, $ N $ \`e l'\textbf{operatore numero}, $ \hat{a} $ e $ \hat{a}^\dag $ sono gli \textbf{operatori di ditruzione e creazione}.

Per calcolare $ K_n $ normalizzo $ \braket{0|0} = 1 $. Allora
\[
\begin{cases}
\ket{n} = K_n (\hat{a}^\dag)^n \ket{0} \\
\braket{n|n} = 1
\end{cases} \quad \Rightarrow
\]
\begin{align*}
\Rightarrow \quad 1 &= |K_n|^2 \braket{0|(\hat{a})^n(\hat{a}^\dag)^n|0} = \\
&= |K_n|^2 \braket{0|(\hat{a})^{n - 1} \hat{a} \hat{a}^\dag (\hat{a}^\dag)^{n - 1}|0} = \\
&= |K_n|^2 \braket{0|(\hat{a})^{n - 1} (\hat{a}^\dag \hat{a} + 1) (\hat{a}^\dag)^{n - 1}|0} = \\
&= |K_n|^2 \braket{0|(\hat{a})^{n - 1} (N + 1) (\hat{a}^\dag)^{n - 1}|0} = \\
&= |K_n|^2 n \braket{0|(\hat{a})^{n - 1} (\hat{a}^\dag)^{n - 1}|0} = \\
&= |K_n|^2 n (n - 1) \braket{0|(\hat{a})^{n - 2} (\hat{a}^\dag)^{n - 2}|0} = \\
&= |K_n|^2 n! \braket{0|0}
\end{align*}
\[
\boxed{K_n = \frac{1}{\sqrt{n!}}} 
\]
da cui segue
\[
\ket{n} = \frac{1}{\sqrt{n!}} (a^\dag)^n \ket{0}
\]
\[
\hat{a}^\dag \ket{n} = \sqrt{n + 1} \ket{n + 1}
\]
\[
\hat{a} \ket{n} = \sqrt{n} \ket{n - 1}
\]
\[
E_n = \hbar \omega \left( n + \frac{1}{2} \right) \text{ autovalori di } \mathcal{H}.
\]

N.B.: il formalismo di Heisenberg ha permesso di riscrivere tutto senza dover scegliere una base.

Gli elementi di matrice li calcolo attraverso $ \hat{a} \pm \hat{a}^\dag $
\[
\hat{x} = \sqrt{\frac{\hbar}{2m\omega}}(\hat{a} + \hat{a}^\dag), \qquad \hat{p} = -i \sqrt{\frac{m\omega\hbar}{2}}(\hat{a} - \hat{a}^\dag)
\]
e usando $ \braket{n|m} = \delta_{nm} $, allora
\[
\begin{cases}
\braket{m|\hat{a}^\dag|n} = \delta_{m, n + 1}\sqrt{n + 1} \\
\braket{m|\hat{a}|n} = \delta_{m, n - 1}\sqrt{n}
\end{cases} \quad \Rightarrow 
\]
\[
\Rightarrow \quad \begin{cases}
\braket{m|\hat{x}|n} = \sqrt{\frac{\hbar}{2m\omega}} (\delta_{m, n - 1}\sqrt{n} + \delta_{m, n + 1}\sqrt{n + 1}) \\
\braket{m|\hat{p}|n} = -i \sqrt{\frac{m\hbar\omega}{2}} (\delta_{m, n - 1}\sqrt{n} - \delta_{m, n + 1}\sqrt{n + 1})
\end{cases}
\]
Allora i valori medi
\[
\braket{n|\hat{x}|n} = \braket{n|\hat{p}|n} = 0.
\]
Per gli operatori al quadrato
\begin{align*}
\braket{n|\hat{x}^2|n} &= \frac{\hbar}{2m\omega} \braket{n|(\hat{a} + \hat{a}^\dag)(\hat{a} + \hat{a}^\dag)|n} = \\
&= \frac{\hbar}{2m\omega} \braket{n|(\hat{a}\hat{a}^\dag + \hat{a}^\dag\hat{a})|n} = \\
&= \frac{\hbar}{2m\omega} \braket{n|(2N + 1)|n} = \\
&= \frac{\hbar}{m\omega} \left( n + \frac{1}{2} \right). \\
\braket{n|\hat{p}^2|n} &= -\frac{m\omega\hbar}{2} \braket{n|(\hat{a}\hat{a}^\dag + \hat{a}^\dag\hat{a})|n} = \\
&= m\omega\hbar \left( n + \frac{1}{2} \right).
\end{align*}

Passo all'energia cinetica e al potenziale (valori medi)
\[
\braket{T} = \frac{\omega\hbar}{2}\left( n + \frac{1}{2} \right),
\]
\[
\braket{V} = \frac{1}{2} m\omega^2 \frac{\hbar}{m\omega} \left( n + \frac{1}{2} \right) = \braket{T},
\]
cio\`e un caso di equipartizione dell'energia.

Le indeterminazioni sono 
\[
\begin{cases}
\braket{n|\Delta^2 \hat{x}|n} = \braket{n|\hat{x}^2|n} \\
\braket{n|\Delta^2 \hat{p}|n} = \braket{n|\hat{p}^2|n}
\end{cases} \quad \Rightarrow
\]
\[
\Rightarrow \quad \Delta^2 \hat{x} \Delta^2 \hat{p} = \hbar^2 \left( n + \frac{1}{2} \right)^2 \geq \frac{\hbar^2}{4}
\]
e $ n = 0 $ d\`a lo stato a minima indeterminazione.
% \subsection{Autofunzioni dell'oscillatore armonico}
\begin{align*}
a \ket{0} &= 0 \text{ è l'equazione che devo risolvere.} \\
\braket{x|n} &= \psi_n(x)
\end{align*}
\begin{align*}
\braket{x|a|0} = 0 &= \sqrt{\frac{m\omega}{2\hbar}} \braket{x|\left( \hat{x} + i\frac{\hat{p}}{m\omega} \right)|0} = \\
&= \sqrt{\frac{m\omega}{2\hbar}} \left( x \psi_0(x) + \frac{i}{m\omega} (-i) \hbar \frac{d}{dx}\psi_0(x) \right),
\end{align*}
\[
\frac{d}{dx} \psi_0(x) = -\frac{m\omega}{\hbar} x \psi_0(x),
\]
\[
\psi_0(x) = N_0 \exp \left\{ -\frac{m\omega}{\hbar} \frac{x^2}{2} \right\}.
\]
cioè è una gaussiana e rappresenta \textbf{lo stato fondamentale dell'atomo di idrogeno}.

Per gli stati superiori
\begin{align*}
\braket{x|1} = \psi_1(x) &= \braket{x|\sqrt{\frac{m\omega}{2\hbar}} \left( \hat{x} - i \frac{\hat{p}}{m\omega} \right)|0} = \\
&= N_1 \left( \exp \left\{ -\frac{m\omega}{\hbar} \frac{x^2}{2} \right\} - \frac{\hbar}{m\omega} \frac{d}{dx} \exp \left\{ -\frac{m\omega}{\hbar} \frac{x^2}{2} \right\} \right) = \\
&= N_1 \left( -\frac{\hbar}{m\omega} \right) \left( -\frac{m\omega}{\hbar} x \right) 2 \exp \left\{ \frac{m\omega}{\hbar} \frac{x^2}{2} \right\} = \\
&= \frac{N_1}{N_0} 2 x \psi_0(x).
\end{align*}
e poi iterando
\[
\psi_n(x) = \braket{x|n} = N_n \psi_0(x) H_n(x)
\]
dove \( H_n \) è un potenziale di grado \( n \). Sono i polinomi di Hermite.
% \subsection{Evoluzione temporale}
\[
\ket{\psi} = \sum_n \ket{n}\braket{n|\psi}
\]
\[
\ket{n, t} = \exp \left\{ \frac{\mathcal{H}t}{i\hbar} \right\} \ket{n} =  e^{-i\omega\left( n + \frac{1}{2} \right) t} \ket{n}
\]
\[
\ket{\psi(t)} = \sum_n c_n^\psi e^{-i\omega \left( n + \frac{1}{2} \right) t} \ket{n}
\]
e questa \`e un'evoluzione temporale alla Schr\"odinger.

Per la dipendenza temporale alla Heisenberg, ho le equazioni differenziali
\[
\frac{d}{dt} \hat{p} = -m\omega^2\hat{x}, \quad \frac{d}{dt}\hat{x} = \frac{\hat{p}}{m}
\]
e devo risolverle.

\textbf{Metodo 1}:
\begin{align*}
\hat{x}_H(t) &= S^\dag(t) \hat{x}_S S(t) = e^{\frac{-\mathcal{H}t}{i\hbar}} \hat{x}_S e^{\frac{\mathcal{H}t}{i\hbar}}, \\
\hat{p}_H(t) &= S^\dag(t) \hat{p}_S S(t) = \hat{p}_S,
\end{align*}
ma in generale non siamo in grado di fare i calcoli. In questo caso c'\`e la formula di Baker-Campbell-Hausdorff. Se \( [A, B] \neq 0 \),
\[
e^{-A} B e^A = B + [-A, B] + \frac{1}{2!} \left[ -A, [-A, B] \right] + \frac{1}{3!} \left[ -A, \left[ -A, [-A, B] \right] \right] + \cdots
\]
\[
e^A e^B = e^{A + B} e^{\frac{1}{2}[A, B]}
\]

Dimostrazione:
\begin{align*}
e^{-\lambda A} B e^{\lambda A} |_{\lambda = 1} &= \sum_{n = 0}^\infty \frac{1}{n!} \left. \frac{d^n}{d\lambda^n} \left( e^{-\lambda A} B e^{\lambda A} \right) \right|_{\lambda = 0} = \\
&= \left. e^{-\lambda A} \left( B + [-A, B]  + \frac{1}{2!} [-A, [-A, B]] + \cdots \right) e^{\lambda A} \right|_{\lambda = 0}. \qed
\end{align*}
\begin{align*}
\hat{x}_H(t) &= \hat{x}_S + \left[ -\frac{1}{i\hbar} \mathcal{H} t, \hat{x}_S \right] + \frac{1}{2} \left(\frac{t}{i\hbar}\right)^2 \left[ \mathcal{H}, \left[ \mathcal{H}, \hat{x}_S \right] \right] + \cdots = \\
&= \hat{x}_S + t \left[ -\frac{\mathcal{H}}{i\hbar}, \hat{x}_S \right] + \frac{t^2}{2} \left[ \frac{\mathcal{H}}{i\hbar}, \left[ \frac{\mathcal{H}}{i\hbar}, \hat{x}_S \right] \right] + \cdots = \\
& \left[ -\frac{\mathcal{H}}{i\hbar}, \hat{x}_S \right] = \frac{\hat{p}}{m} \\
& \left[ -\frac{\mathcal{H}}{i\hbar}, \frac{\hat{p}_S}{m} \right] = -\omega^2\hat{x}_S \\
& \left[ -\frac{\mathcal{H}}{i\hbar}, -\omega^2\hat{x}_S \right] = -\omega^2 \frac{\hat{p}_S}{m} \\
&= \frac{\hat{p}_S}{m} \left( t + \frac{t^3}{3!} (-\omega^2) + \frac{t^5}{5!} (-\omega^2)^2 + \cdots \right) + \\
&+\hat{x}_S \left( 1 + \frac{t^2}{2!} (-\omega^2) + \frac{t^4}{4!} (-\omega^2)^2 + \cdots \right) = \\
&=\hat{x}_S \cos (\omega t) + \frac{\hat{p}_S}{m\omega} \sin(\omega t).
\end{align*}

\textbf{Metodo 2}: uso gli operatori di creazione e distruzione.
\[
\hat{a} = \sqrt{\frac{m\omega}{2\hbar}} \left( \hat{x} + i \frac{\hat{p}}{m\omega} \right), \; [N, \hat{a}] = -\hat{a}, [\mathcal{H}, \hat{a}] = \left[ \hbar \omega \left( N + \frac{1}{2} \right), \hat{a} \right] = -\hbar \omega \hat{a}
\]
\[
\frac{d \hat{a}_H}{dt} = \frac{1}{i\hbar} [\hat{a}, \mathcal{H}] = -i\omega \hat{a} \quad \Rightarrow \quad \hat{a}_H(t) = e^{-i\omega t} \hat{a}_S
\]
\[
\frac{d \hat{a}_H^\dag}{dt} = \frac{1}{i\hbar} [\hat{a}^\dag, \mathcal{H}] = i\omega \hat{a}^\dag \quad \Rightarrow \quad \hat{a}_H^\dag(t) = e^{i\omega t} \hat{a}_S^\dag
\]
cos\`{i} disaccoppio \( \hat{x} \) e \( \hat{p} \) nelle equazioni differenziali ed ho lo stesso risultato.

Gli elementi di matrice
\begin{align*}
\braket{n + 1|\hat{a}_H^\dag(t)|n} &= e^{i\omega t} \braket{n + 1|\hat{a}_S^\dag|n} = e^{i\omega t} \sqrt{n + 1} \\
&= \braket{n + 1; t|\hat{a}_S^\dag|n; t} = \\
&= \braket{n + 1|e^{-\frac{\hbar\omega}{i\hbar}\left( n + 1 + \frac{1}{2} \right)} \hat{a}_S^\dag e^{\frac{\hbar\omega}{i\hbar} \left( n + \frac{1}{2} \right)}|n} = \\
&= e^{i\hbar\omega}\braket{n + 1|\hat{a}_S^\dag|n}.
\end{align*}
% \subsection{Stati coerenti}
Sovrapposizione di stati di oscillatore armonico \( \Rightarrow \) pacchetti d'onda ad indeterminazione minima posizione-impulso che non si allargano perch\'e sono soggetti a potenziale. Definisco questi stati in modo che 
\[
\hat{a} \ket{z} = z \ket{z}, \quad z \in \mathbb{C}, \hat{a} \text{ operatore distruzione},
\]
\[
\mathcal{H} = \frac{\hat{p}^2}{2m} + \frac{1}{2} m\omega^2 \hat{x}^2,
\]
\[
\ket{z} = \sum_n \ket{n}\braket{n|z}.
\]
Ricordo che \( N \ket{n} = n \ket{n}, N = \hat{a}^\dag \hat{a} \), quindi
\begin{align*}
\hat{a} \ket{z} &= \underbrace{\sum_{n = 0}^\infty \hat{a} \ket{n} \braket{n|z}}_{\text{vuoto annichilato}} = \sum_{n = 1}^\infty \sqrt{n} \ket{n - 1} \braket{n|z} \\
&= \sum_{n = 0}^\infty \sqrt{n + 1} \ket{n} \braket{n + 1|z}.
\end{align*}
Allora la condizione di autostato \`e data da 
\[
z \braket{n|z} = \sqrt{n + 1} \braket{n + 1|z}
\]
\[
\braket{n - 1|z} = \sqrt{n} \braket{n|z} \frac{1}{z}
\]
\begin{align*}
\braket{n|z} &= \frac{z}{\sqrt{n}} \braket{n - 1|z} = \\
&= \frac{z^2}{\sqrt{n}\sqrt{n - 1}} \braket{n - 2|z} = \\
&= \frac{z^n}{\sqrt{n!}} \braket{0|z}
\end{align*}
quindi scrivo \( \ket{z} \) come
\begin{align*}
\ket{z} &= \sum_{n = 0}^\infty \ket{n} \frac{z^n}{\sqrt{n!}} \braket{0|z} = \sum_{n = 0}^\infty \frac{(\hat{a}^\dag)^n}{\sqrt{n!}\sqrt{n!}} \ket{0} z^n \braket{0|z} = \\
&= \sum_{n = 0}^\infty \frac{(z\hat{a}^\dag)^n}{n!} \ket{0} \braket{0|z} = e^{z\hat{a}^\dag} \ket{0} \braket{0|z}.
\end{align*}
Impongo che \( \braket{z|z} = 1 \), allora 
\[
1 = |\braket{0|z}|^2 \sum_{n, m} \braket{m|n} \frac{z^n {z^*}^m}{\sqrt{n!m!}} = \sum_{n = 0}^\infty \frac{|z|^{2n}}{n!} |\braket{0|z}|^2 = e^{|z|^2} |\braket{0|z}|^2
\]
\[
\braket{0|z} = e^{-\frac{1}{2}|z|^2}
\]
\[
\boxed{\ket{z} = e^{-\frac{1}{2}|z|^2} e^{z\hat{a}^\dag} \ket{0}},
\]
questa \`e la forma degli \textbf{stati coerenti}.

I valori medi sono 
\[
\braket{z|\hat{x}|z} = \sqrt{\frac{\hbar}{2m\omega}} \braket{z|(\hat{a} + \hat{a}^\dag)|z} = \sqrt{\frac{\hbar}{2m\omega}} (z + z^*) = \sqrt{\frac{2\hbar}{m\omega}} \, \text{Re} \,  z
\]
\[
\braket{z|\hat{p}|z} = -i\sqrt{\frac{m\hbar\omega}{2}}\braket{z|\hat{a} - \hat{a}^\dag|z} = \sqrt{2m\omega\hbar} \, \text{Im} \, z
\]
e quindi sono valori oscillanti. Inoltre
\[
\braket{\hat{x}} = 2 \, \text{Re}z = \Delta \hat{x},
\]
\[
\braket{\hat{p}} = 2 \, \text{Im}z = \Delta\hat{p}.
\]

Gli scarti quadratici
\begin{align*}
\braket{z|\Delta^2 \hat{x}|z} &= \braket{z|(\hat{x} - \braket{\hat{x}})^2|z} = \\
&= \frac{\hbar}{2m\omega} \braket{z|((\hat{a} + \hat{a}^\dag) - (\braket{\hat{a}} + \braket{\hat{a}^\dag}))((\hat{a} + \hat{a}^\dag) - (\braket{\hat{a}} + \braket{\hat{a}^\dag}))|z} = \\
&= \frac{\hbar}{2m\omega} \braket{z|(\hat{a} + z^* - z - z^*)(z + \hat{a}^\dag - z - z^*)|z} = \\
&= \frac{\hbar}{2m\omega} \braket{z|(\hat{a} - z)(\hat{a}^\dag - z^*)|z} = \\
&= 0 + \frac{\hbar}{2m\omega} \braket{z|[\hat{a} - z, \hat{a}^\dag - z^*]|z} = \\
&= \frac{\hbar}{2m\omega}, \\
\braket{z|\Delta^2 \hat{p}|z} &= \frac{\hbar m\omega}{2}
\end{align*}
danno l'indeterminazione minima 
\[
\boxed{\Delta^2 \hat{x} \Delta^2 \hat{p} = \frac{\hbar^2}{4}}.
\]
Ai tempi successivi, poich\'e \( \hat{a}(t) = e^{-i\omega t}\hat{a} \), l'indeterminazione si conserva, mentre i valori medi oscillano.

Per due stati coerenti 
\[
\braket{z_1|z_2} = \sum_{n = 0}^\infty \frac{(z_1^* z_2)^n}{n!} e^{-\frac{1}{2}|z_1|^2} e^{\frac{1}{2}|z_2|^2} = e^{z_1^* z_2} e^{-\frac{|z_1|^2 - |z_2|^2}{2}},
\]
ad esempio, per \( \braket{z|-z} = e^{-2|z|^2} \), cio\`e ho un decadimento esponenziale di \( \braket{z|-z} \), ma poich\'e 
\[
|z|^2 = (\text{Re}z)^2 + (\text{Im}z)^2 = \frac{\braket{\hat{p}}^2}{4(\Delta\hat{p})^2} + \frac{\braket{\hat{x}}^2}{4(\Delta\hat{x})^2},
\]
\[
\braket{z|-z} = e^{2 \left( \frac{1}{4} + \frac{1}{4} \right)} = e.
\]
% \pagebreak

% \section{Secondo modulo}
% \subsection{Sistemi \(n\)-dimensionali}

\[
\hat{x} \ket{x} = x \ket{x}, \quad p_\psi = |\braket{x|\psi}|^2, \quad \braket{x|\psi} = \psi(x).
\]
Introduco vettore di operatori
\[
\hat{\vec{x}} = \begin{pmatrix}
\hat{x}_{(1)} \\ \hat{x}_{(2)} \\ \vdots \\ \hat{x}_{(n)}
\end{pmatrix}
\]
e agisce
\[
\hat{\vec{x}}\ket{\vec{x}} = \vec{x} \ket{\vec{x}}, \quad \text{dove } \ket{\vec{x}} = \ket{x^{(1)}} \otimes \ket{x^{(2)}} \otimes \cdots \otimes \ket{x^{(n)}}.
\]
Ma allora per
\[
\braket{\vec{x}|\psi} = \psi(x^{(1)}, x^{(2)}, \dots, x^{(n)}) = \psi(\vec{x}),
\]
però in generale
\[
\psi(\vec{x}) \neq \psi^{(1)}(x_1) \psi^{(2)}(x_2) \cdots \psi^{(n)}(x_n).
\]
Nel caso di \( n \) paricelle, all'\(i\)-esima particella associo
\[
\ket{e_i^{(1)}} \otimes \cdots \otimes \ket{e_i^{(n)}}.
\]

Ad esempio, per il qubit ho \( \ket{+} \) e \( \ket{-} \), ma quattro possibili stati se ho due qubit:
\[
\ket{+} \otimes \ket{+}, \ket{+} \otimes \ket{-}, \ket{-} \otimes \ket{+}, \ket{-} \otimes \ket{-}
\]
\[
\ket{++}, \quad \ket{+-}, \quad \ket{-+}, \quad \ket{--}
\]
e un operatore \( \hat{A} \) che agisce su questo sistema
\[
\braket{mn|\hat{A}|ij} = A_{mn,ij} \quad \rightarrow \quad A_{ab} \begin{cases}
\ket{ij} = \ket{b} \\
\ket{mn} = \ket{a}
\end{cases}.
\]

Allora \( \mathcal{H} \) va generalizzato
\[
\mathcal{H} = \frac{\left| \hat{\vec{p}} \right|^2}{2m} + V(\hat{\vec{x}}).
\]
Partendo da 1 dimensione
\begin{align*}
\braket{x|\hat{p}|\psi} &= -i\hbar \frac{d}{dx} \psi(x), \quad \braket{\vec{x}|\hat{\vec{p}}|\psi} = -i\hbar \vec{\nabla} \psi(\vec{x}), \\
\braket{x|\frac{\hat{p}^2}{2m}|\psi} &= \left( -i\hbar \frac{d}{dx} \right)^2 \frac{1}{2m} \psi(x) = -\frac{\hbar^2}{2m} \frac{d^2}{dx^2} \psi(x) \\
\braket{\vec{x}|\frac{\hat{\vec{p}}}{2m}|\psi} &= \frac{1}{2m} \braket{\vec{x}|\hat{\vec{p}} \boldsymbol{\cdot} \hat{\vec{p}}|\psi} = \frac{1}{2m} \braket{\vec{x}|({p^{(1)}}^2 + {p^{(2)}}^2 + \cdots {p^{(n)}}^2)|\psi} = \\
&= -\frac{\hbar^2}{2m} \left( \frac{d^2}{d{x^{(1)}}^2} + \cdots + \frac{d^2}{d{x^{(n)}}^2} \right) \psi(\vec{x}) = \\
&= -\frac{\hbar^2}{2m} \vec{\nabla} \boldsymbol{\cdot} \vec{\nabla} \psi(\vec{x}) = \\
&= -\frac{\hbar^2}{2m} \Delta \psi(\vec{x}).
\end{align*}
% \subsection{Potenziali separabili in coordinate cartesiane}
Il problema $ n $-dimensionale si riduce a $ n $ problemi 1-dimensionali.
\[
V(\vec{x}) = \sum_i V_i(x^{(i)})
\]
N.B.: $ \hat{T} $ \`e sempre separabile.
\[
\mathcal{H} = \frac{{\hat{\vec{p}}}^2}{2m} + \sum_i V_i(\hat{x}^{(i)}) = \sum_i \mathcal{H}_i,
\]
\[
\mathcal{H}_i = \frac{{\hat{p}^{(i)}}^2}{2m} + V_i(\hat{x}^{(i)}).
\]

Allora faccio
\begin{align*}
\braket{\vec{x}|\mathcal{H}|\psi} &= \braket{\vec{x}|\mathcal{H}_1|\psi} + \braket{\vec{x}|\mathcal{H}_2|\psi} + \cdots = \\
&= \left( -\frac{\hbar^2}{2m} \frac{d^2}{d{x^{(1)}}^2} + V(\hat{x}_1) \right) \psi(\vec{x}) \\
&+ \left( -\frac{\hbar^2}{2m} \frac{d^2}{d{x^{(2)}}^2} + V(\hat{x}_2) \right) \psi(\vec{x}) + \cdots
\end{align*}

Se inoltre ho $ \mathcal{H} \ket{\psi} = E \ket{\psi} $,
\[
\mathcal{H} \ket{\psi_n} = E_n \ket{\psi_n}, \quad \braket{x^{(i)}|\psi_n} = \psi_n(x^{(i)}).
\]
Se \`e $ d $ il numero di dimensioni, allora 
\begin{align*}
\psi_{n_1, \dots, n_d}(\vec{x}) &= \psi_{n_1}(x^{(1)}) \cdots \psi_{n_d}(x^{(d)}) \\
\braket{\vec{x}|\mathcal{H}|\psi_{n_1, \dots, n_d}}  &= \frac{\hbar^2}{2m} \left( \frac{d^2}{d{x^{(1)}}^2} \psi_{n_1}(x^{(1)}) \right) \psi_{n_2}(x^{(2)}) \cdots \psi_{n_d}(x^{(d)}) + \\
&+ V_1(x^{(1)}) \psi_{n_1}(x^{(1)}) \cdots \psi_{n_d}(x^{(d)}) + \\
&- \frac{\hbar^2}{2m} \psi_{n_1}(x^{(1)}) \left( \frac{d^2}{d{x^{(2)}}^2} \psi_{n_2}(x^{(2)}) \right) \psi_{n_3}(x^{(3)}) \cdots \psi_{n_d}(x^{(d)}) + \cdots \\
&= (E_{n_1} + \cdots + E_{n_d}) \psi_{n_1, n_2, \dots, n_d}(\vec{x}) \quad \Rightarrow
\end{align*}
\[
\Rightarrow \quad -\frac{\hbar^2}{2m} \frac{d^2}{d{x^{(i)}}^2} \psi_{n_i}(x^{(i)}) + V_i(x^{(i)})\psi_{n_i}(x^{(i)}) = E_{n_i} \psi_{n_i}(x^{(i)})
\]
% \subsection{Relazioni di commutazione}
1 dimensione $ \rightarrow [\hat{p}, \hat{x}] = -i\hbar $ \newline
$ d $ dimensioni $ \rightarrow [\hat{p}^{(i)}, \hat{x}^{(j)}] = -i\hbar \delta^{(ij)} $
$$
\left( \hat{p}^{(i)} \rightarrow -i\hbar \frac{\partial}{\partial x^{(i)}} \qquad -i\hbar\left[ \frac{\partial}{\partial x^{(i)}}, \hat{x}^{(j)} \right] = -i\hbar \delta_{ij} \right)
$$
% \subsection{Alcune semplificazioni}
$ \hat{\vec{p}} \to -i\hbar\vec{\nabla} $ vuol dire
$$
\braket{\vec{x}|\hat{\vec{p}}|\psi} = -i\hbar\vec{\nabla}\psi(\vec{x})
$$
nella rappresentiazione delle coordinate scrivo sempre $ \vec{p} = -i\hbar \vec{\nabla} $, quindi
$$
-i\hbar\left[ \frac{\partial}{\partial x^{(i)}}, \hat{x}^{(j)} \right] = [\hat{p}^{(i)}, \hat{x}^{(j)}] = -i\hbar \left( \frac{\partial}{\partial x^{(i)}} x^{(j)} - x^{(j)} \frac{\partial}{\partial x^{(i)}} \right) =
$$
$$
= -i\hbar\left( \delta_{ij} + x^{(j)} \frac{\partial}{\partial x^{(i)}} - x^{(j)} \frac{\partial}{\partial x^{(i)}} \right) = -i\hbar\delta_{ij}.
$$
% \subsection{Hamiltoniane}
$$
[\mathcal{H}_i(\hat{x}^{(i)}, \hat{p}^{(i)}), \mathcal{H}_j(\hat{x}^{(j)}, \hat{p}^{(j)})] = 0, \quad \text{se } i \neq j.
$$
La $ \mathcal{H} $ \`e somma di hamiltoniane che commutano fra loro, quindi gli autostati della $ \mathcal{H} $ sono tutti gli autostati di $ \mathcal{H}_i $. Allora
$$
\mathcal{H} \ket{\psi_E} = E \ket{\psi_E} \quad \Rightarrow \quad E_{k_1, \dots, k_d} = E_{k_1} + \cdots + E_{k_d}.
$$
Inoltre,
$$
\mathcal{H}_i \ket{\psi_k^{(i)}} = E_k \ket{\psi_k^{(i)}}, \qquad \braket{x^{(i)}|\psi_k^{(i)}} = \psi_k^{(i)}(x^{(i)})
$$
e ne segue che 
$$
\psi(\vec{x}) = \psi_{k_1}(x^{(1)}) \cdots \psi_{k_d}(x^{(d)}).
$$
Se il problema \`e separabile, ho questa utile condizione.

Esempio: 
$$
V(\vec{x}) = \sum_{i = 1}^3 V_i(x_i), \qquad V_i(x_i) = \begin{cases}
0 & |x_i| < a_i \\
+\infty & |x_i| > a_i
\end{cases}.
$$
Il problema si riduce a tre problemi unidimensionali.
$$
\left( \frac{{\hat{p}^{(i)}}^2}{2m} + V_i(x^i) \right) \ket{\psi_{n_i}} = E_{n_i} \ket{\psi_{n_i}},
$$
$$
E_{n_i} = \frac{\hbar^2 k_{n_i}^2}{2m}, \qquad k_{n_i} = n_i \frac{\pi}{2a_i},
$$
$$
\psi_{n_i}(x^{(i)}) = \begin{cases}
A_{n_i} \sin(k_{n_i} x_i) \; n_i \text{ pari} \\
B_{n_i} \cos(k_{n_i} x_i) \; n_i \text{ dispari}
\end{cases}.
$$

Nel caso di buca cubica, $ a_1 = a_2 = a_3 = a $ e lo spettro
$$
E_{n_1, n_2, n_3} = \frac{\hbar^2 \pi^2}{8ma^2} (n_1^2 + n_2^2 + n_3^2),
$$
$$
\text{stato fonamentale} \, E_{111} = \frac{3}{8} \frac{\hbar^2 \pi^2}{ma^2},
$$
$$
1^{\text{o}} \text{ stato eccitato (tre volte degenere) } E_{211} = E_{121} = E_{112} = \frac{3}{4} \frac{\hbar^2 \pi^2}{ma^2}.
$$
% \subsection{Oscillatore armonico 3-dimensionale}
\[
V_i(x^{(i)}) = \frac{1}{2} m\omega^2{x^{(i)}}^2,
\]
\[
\mathcal{H} \ket{\psi_{n_1 n_2 n_3}} = E_{n_1 n_2 n_3} \ket{\psi_{n_1 n_2 n_3}}.
\]
Per tre dimensioni
\[
E_{n_i} = \hbar \omega_i \left( n_i + \frac{1}{2} \right),
\]
\[
E_{n_i} = \hbar \left( \omega_1 n_1 + \omega_2 n_2 + \omega_3 n_3 + \frac{1}{2} (\omega_1 + \omega_2 + \omega_3) \right).
\]
Nel caso di $ \omega_1 = \omega_2 = \omega_3 = \omega $, allora
\[
V(\vec{x}) = \frac{1}{2} m\omega^2 \sum_{i = 1}^3 {x^{(i)}}^2 = \frac{1}{2} m\omega^2 \left\Vert \vec{x} \right\Vert^2
\]
e ho potenziale armonico sferico, quindi il caso dell'oscillatore armonico isotropo (sferico).

In questo caso la degenerazione si ha da 
\[
E_{n_1 n_2 n_3} = \hbar \omega \left( n_ 1 + n_2 + n_3 + \frac{3}{2} \right) = \hbar \omega \left( N + \frac{3}{2} \right), \quad N = 0, 1, 2, \dots
\]
ed \`e calcolabile esplicitamente perch\'e per ogni $ N $ fissato scelgo
\[
n_1 = 0, 1, 2, \dots, N \text{ e poi ho su } n_2 \text{ la condizione} \\
\]
\[
n_2 = 0, 1, \dots, N - n_1, \text{ dopodich\'e} \\
\]
\[
n_3 \text{ \`e dato univocamente.}
\]
quindi degenerazione si ha quando
\[
\sum_{n_1 = 0}^N (N - n_1 + 1) = (N + 1) \sum_{n_1 = 0}^N 1 - \sum_{n_1 = 0}^N n_1 = \frac{1}{2} (N + 1)(N + 2).
\]
% \subsection{Problema dei due corpi}
\[
\mathcal{H} = \frac{{\vec{p}_1}^2}{2m_1} + \frac{{\vec{p}_2}^2}{2m_2} + V(\vec{x_1} - \vec{x_2}), \quad \text{corpi 1 e 2.}
\]
Voglio ridurmi al problema ad un solo corpo e poi passo al potenziale centrale, dove uso le coordinate sferiche.

So che
\[
[\hat{p}_a^{(i)}, \hat{x}_b^{(j)}] = -i\hbar \delta^{ij} \delta^{ab},
\]
\[
[\hat{p}_a^{(i)}, \hat{p}_b^{(j)}] = [\hat{x}_a^{(i)}, \hat{x}_b^{(j)}] = 0.
\]
Poi introduco
\[
\begin{dcases}
\vec{r} = \vec{x}_1 - \vec{x}_2 \\
\vec{R} = \frac{m_1 \vec{x}_1 + m_2 \vec{x}_2}{m_1 + m_2}
\end{dcases} \quad \text{come nuove coordinate}
\]
e poi costruisco i momenti coniugati tali che
\[
\begin{cases}
[\hat{p}^{(i)}, \hat{r}^{(j)}] = -i\hbar\delta_{ij} = [P^{(i)}, R^{(j)}] \\
[\hat{p}^{(i)}, R^{(j)}] = [P^{(i)}, r^{(j)}] = 0
\end{cases}.
\]
La scelta che funziona è
\[
\begin{dcases}
\hat{\vec{p}} = \frac{m_2 \vec{p}_1 - m_1 \vec{p}_2}{m_1 + m_2} \\
\vec{P} = \hat{\vec{p}}_1 + \hat{\vec{p}}_2
\end{dcases}
\]
che verificano le regole di commutazione (cambio di variabili canonico). In generale
\[
\begin{cases}
\vec{x_1'} = \vec{r} \\
\vec{x_2'} = \vec{R}
\end{cases}, \quad \begin{cases}
\vec{p_1'} = \vec{p} \\
\vec{p_2'} = \vec{P}
\end{cases}
\]
le scrivo come 
\[
\begin{pmatrix}
\vec{x_1'} \\
\vec{x_2'}
\end{pmatrix} = M \begin{pmatrix}
\vec{x_1} \\
\vec{x_2}
\end{pmatrix}, \quad \begin{pmatrix}
\vec{p_1'} & \vec{p_2'}
\end{pmatrix} = N \begin{pmatrix}
\vec{p_1} & \vec{p_2}
\end{pmatrix}, \quad \text{con } N = M^{-1}
\]
\[
\boxed{N = M^{-1} \Leftrightarrow [p_a^{(i)}, p_b^{(j)}] = -i\hbar\delta_{ij}\delta_{ab}} \quad \text{(voglio verificarlo)}.
\]
\begin{align*}
[{p'}_{a'}^{(i)}, {x'}_{b'}^{(j)}] &= [p_a^{(i)}N_{aa'}, M_{b'b} x_b^{(j)}] = \\
&= [p_a^{(i)}, x_b^{(j)}] M_{b'b} N_{aa'} = \\
&= -i\hbar\delta_{ij}\delta_{ab}N_{aa'}M_{b'b} = \\
&= -i\hbar\delta_{ij}M_{b'a}N_{aa'} = \\
&= i\hbar\delta_{ij} (MN)_{b'a'} \stackrel{\text{per ipotesi}}{=} -i\hbar\delta_{ij}(\mathbbm{1})_{b'a'} \Leftrightarrow M = N^{-1}.
\end{align*}
Quindi, come avevo prima,
\[
{p'}_{a'}^{(i)} = -i\hbar\frac{\partial}{\partial {x'}_{a'}^{(i)}},
\]
infatti, per \( (i) \) fissato,
\[
-i\hbar \frac{\partial}{\partial {x'}_{a'}^{(i)}} = -i\hbar \frac{\partial}{\partial x_a^{(i)}} \frac{\partial x_a^{(i)}}{\partial {x'}_{a'}^{(i)}} \quad \Rightarrow \quad M_{aa'} = M_{aa'}^{-1}
\]
\[
\frac{\partial {x'}_{a'}^{(i)}}{\partial x_a^{(i)}} = M_{aa'} \Rightarrow -i\hbar\frac{\partial}{\partial {x'}_{a'}^{(i)}} = -i\hbar \frac{\partial}{\partial x_a^{(i)}} M_{a'a} = p_a^{(i)} M_{aa'}^{-1}
\]
e poi si dimostra che 
\[
\mathcal{H} = \frac{\vec{P}^2}{2M} + \frac{\vec{p}^2}{2\mu} + V(\vec{r}), \quad M = m_1 + m_2, \quad \mu = \frac{m_1 m_2}{m_1 + m_2}.
\]

D'ora in poi studierò solo sistemi ad un corpo, infatti, nel sistema del centro di massa,
\[
\mathcal{H} = \frac{\vec{p}^2}{2m} + V(r)
\]
e tipicamente userò coordinate sferiche 
\[
\begin{cases}
x = r \sin \theta \cos \varphi \\
y = r \sin \theta \sin \varphi \\
z = r \cos \theta
\end{cases}
\]
però adesso devo sepapare \( \vec{p} = -i\hbar \vec{\nabla} \). Per l'impostazione radiale
\[
\hat{p} = \frac{\vec{x}}{\left\Vert \vec{x} \right\Vert} \mathbf{\cdot} \vec{p}
\]
e vedo se vado da qualche parte. Introduco 
\[
\frac{\vec{x}}{\left\Vert \vec{x} \right\Vert} \mathbf{\cdot} \vec{\nabla} = ?
\]
Prendo la componente \( \partial_i \)
\[
\begin{cases}
\partial_i = \frac{\partial}{\partial r} \partial_i r + \frac{\partial}{\partial \theta} \partial_i \theta + \frac{\partial}{\partial \varphi} \partial_i \varphi \\
\frac{\partial \theta}{\partial x_i} x^i = \frac{\partial \varphi}{\partial x_i} x_i = 0
\end{cases} \quad \Rightarrow
\]
\[
\Rightarrow \quad \frac{x_i}{r} \frac{\partial}{\partial r} \frac{\partial}{\partial x_i} |\vec{x}| = \frac{x_i}{r} \frac{\partial}{\partial r} \frac{x_i}{|\vec{x}|} = \frac{x_i^2}{r^2} \frac{\partial}{\partial r} = \frac{\partial}{\partial r}.
\]

Quindi ottengo che \( [\tilde{p}_r, r] = -i\hbar \). Noto che anche \( [\tilde{p}_r + f(r), r] = -i\hbar \). Il secondo termine è critico, quindi 
\[
T = \frac{{\vec{p}}^2}{2m} = \frac{\Delta}{2m} = \sum_i \frac{\partial_i^2}{2m}.
\]
Uso l'equazione classica
\[
(a \mathbf{\cdot} b) = \left\Vert \vec{a} \right\Vert^2 \left\Vert \vec{b} \right\Vert^2 - \left\Vert \vec{a} \wedge \vec{b} \right\Vert^2
\]
e la dimostro, introducendo
\[
\varepsilon^{ijk} = \begin{cases}
0 \quad \text{se due indici hanno lo stesso valore} \\
1 \quad \varepsilon^{123} \\
\text{antisimmetrico rispetto a scambio di indici}
\end{cases} .
\]
Allora
\[
\left( \vec{a} \wedge \vec{b} \right)_i = \varepsilon_{ijk} a_j b_k
\]
\[
\varepsilon^{ijk} \varepsilon^{iab} = \delta^{ia}\delta^{kb} - \delta^{jb} \delta^{ka}
\]
\[
\left\Vert \vec{a} \wedge \vec{b} \right\Vert^2 = \varepsilon^{imn} \varepsilon^{ijk} a_n b_n a_i b_k
\]

Voglio sapere che cos'è \( \hat{\vec{p}}^2 \), partendo dal caso classico, per arrivare a quello quantistico. Uso l'identità vettoriale, che nel caso classico è 
\[
(\vec{a} \mathbf{\cdot} \vec{b})^2 = |\vec{a}|^2 + |\vec{b}|^2 - |\vec{a} \wedge \vec{b}|^2 \quad (\text{da dimostrare}).
\]
Con ciò visto sopra
\[
\varepsilon^{ijk} \varepsilon^{iab} = \delta^{ja} \delta^{kb} - \delta^{jb} \delta^{ka},
\]
\[
(\vec{a} \mathbf{\cdot} \vec{b})_i = \varepsilon_{ijk} a^j b_k
\]
\begin{align*}
|\vec{a} \wedge \vec{b}|^2 &= \varepsilon^{ijk} \varepsilon^{ilm} a^j b^l a^k b^k = \\
&= (\delta^{jl} \delta^{km} - \delta^{jm} \delta^{kl}) a^j b^k a^l b^m = \\
&= a^j a^j b^k b^k - a^j b^j a^k b^k = \\
&= |\vec{a}|^2 + |\vec{b}|^2 - (\vec{a} \mathbf{\cdot} \vec{b})^2
\end{align*}
e ho ciò che cercavo. Poi sostituisco \( \vec{a} \) e \( \vec{b} \) con \( \vec{x} \) e \( \vec{p} \) e ho la relazione classica 
\[
\vec{p}^2 = \vec{p}_r^2 + \frac{\vec{L}^2}{r^2}.
\]

Passo alla quantistica, con \( \hat{\vec{x}} \) e \( \hat{\vec{p}} \) (che non commutano)
\begin{align*}
\varepsilon^{ijk} \varepsilon^{ilm} \hat{x}^j \hat{p}^m &= (\delta^{jl}\delta^{km} - \delta^{jm}\delta^{kl}) \hat{x}^j \hat{p}^k \hat{x}^l \hat{p}^m = \\
&= \hat{x}^j \hat{p}^k \hat{x}^j \hat{p}^k - \hat{x}^j \hat{p}^k \hat{x}^k \hat{p}^j = \\
&= \hat{x}^j \hat{x}^j \hat{p}^k \hat{p}^k + \hat{x}^j [\hat{p}^k, \hat{x}^j] \hat{p}^k - (\hat{x}^j \hat{p}^k \hat{p^j} \hat{x}^k + \hat{x}^j \hat{p}^k [\hat{x}^k, \hat{p}^j]) = \\
&= \hat{r}^2 \hat{\vec{p}}^2 - i\hbar \hat{x}^j \hat{p}^j - (\hat{x}^j \hat{p}^j \hat{p}^k \hat{x}^k + i\hbar \hat{x}^j \hat{p}^j) = \\
&= \hat{r}^2 \hat{\vec{p}}^2 - i\hbar \hat{x}^j \hat{p}^j - ((\hat{\vec{x}} \mathbf{\cdot} \hat{\vec{p}})(\hat{\vec{x}} \mathbf{\cdot} \hat{\vec{p}}) + \hat{x}^j \hat{p}^j [\hat{p}^k, \hat{x}^k] + i\hbar \hat{x}^j \hat{p}^j) = \\
&= \hat{r}^2 \hat{\vec{p}}^2 -i\hbar\hat{\vec{x}} \mathbf{\cdot} \hat{\vec{p}} - (\hat{\vec{x}} \mathbf{\cdot} \hat{\vec{p}}) (\hat{\vec{x}} \mathbf{\cdot} \hat{\vec{p}}) - (-3i\hbar\hat{\vec{x}} \mathbf{\cdot} \hat{\vec{p}} + i\hbar\hat{\vec{x}} \mathbf{\cdot} \hat{\vec{p}}) = \\
&= \hat{r}^2 \hat{\vec{p}}^2 + i\hbar \hat{\vec{x}} \mathbf{\cdot} \hat{\vec{p}} - (\hat{\vec{x}} \mathbf{\cdot} \hat{\vec{p}})^2 = |\vec{x} \wedge \vec{p}|^2 \equiv \vec{L}^2
\end{align*}
dove \( \vec{L} \) ha la proprietà che
\[
\vec{x} \mathbf{\cdot} \vec{L} = \varepsilon^{ijk} x^i x^j p^k = \frac{1}{2} \varepsilon^{ijk} [x^i, x^j] p^k = 0.
\]
\end{document}